% Suites numériques — Mathématiques Tle C — Cours signé OnBuch+
% Programme officiel MINESEC. Compiler avec : tectonic M04-suites-numeriques.tex
\newcommand{\DOCMATIERE}{Mathématiques}
\newcommand{\DOCNIVEAU}{Tle C}
\newcommand{\DOCMODULE}{Relations et opérations fondamentales dans l'ensemble des nombres réels}
\newcommand{\DOCLECON}{04}
\newcommand{\DOCTITRE}{Suites numériques}
% OnBuch+ — préambule « cours signés » ENRICHI (compilé avec tectonic / XeLaTeX)
% Système visuel « Pop » d'OnBuch+ : fond crème (jamais blanc), bordures encre
% épaisses, ombres « dures » décalées, titres Archivo Black en capitales.
% Version MATHÉMATIQUES : graphes (pgfplots/tikz), boîtes pédagogiques
% (définition, propriété, méthode, attention, le savais-tu, exercice résolu,
% exercice, TP, bilan), page de garde et signature OnBuch+ ancrée en bas.
%
% Macros attendues dans main.tex AVANT \input{preamble} :
%   \DOCMATIERE  \DOCNIVEAU  \DOCMODULE  \DOCLECON (numéro)  \DOCTITRE
\documentclass[11pt]{article}

\usepackage{fontspec}
\usepackage[french]{babel}
\usepackage[a4paper,margin=2cm,top=3.4cm,bottom=2.8cm,headheight=1.4cm,footskip=1.3cm]{geometry}
\usepackage{amsmath,amssymb,amsfonts}
\usepackage{mathtools} % \xmapsto, \coloneqq... souvent utilisés par les modèles
\usepackage{cancel} % simplifications barrées
% \abs{z} (module, valeur absolue) et \norm{u} (norme), utilisés par les modèles
\providecommand{\abs}{}\let\abs\relax\DeclarePairedDelimiter{\abs}{\lvert}{\rvert}
\providecommand{\norm}{}\let\norm\relax\DeclarePairedDelimiter{\norm}{\lVert}{\rVert}
\usepackage[table]{xcolor} % [table] : fournit aussi \rowcolors (alternance de lignes)
\usepackage{pagecolor}
\usepackage{tikz}
\usetikzlibrary{arrows.meta,positioning,calc,shapes.geometric,shapes.misc,shapes.arrows,decorations.pathmorphing,decorations.markings,patterns,shadows,mindmap,trees,fit,backgrounds,matrix,intersections}
\usepackage{pgfplots}
\usepackage{pgf-pie} % diagrammes circulaires \pie (statistiques)
\pgfplotsset{compat=1.18}
\usepgfplotslibrary{fillbetween,groupplots} % aires entre courbes (calcul intégral)
\usepackage{enumitem}
\usepackage{fancyhdr}
% [most] charge toutes les bibliothèques tcolorbox (listings, minted, raster,
% documentation...) et gonfle inutilement la mémoire TeX sur un document long
% (~300 boîtes) : on ne charge que ce qui est réellement utilisé.
\usepackage[skins,breakable]{tcolorbox}
\usepackage{graphicx}
\usepackage{colortbl}
\usepackage{booktabs}
\usepackage{tabularx}
\usepackage{diagbox} % case d'en-tête diagonale des tableaux à double entrée
% Colonnes extensibles centrée / gauche / droite (C, L, R), souvent utilisées par les modèles
\newcolumntype{C}{>{\centering\arraybackslash}X}
\newcolumntype{L}{>{\raggedright\arraybackslash}X}
\newcolumntype{R}{>{\raggedleft\arraybackslash}X}
\usepackage{array}
\usepackage{multicol}
\usepackage{siunitx}
% parse-numbers=false : accepte aussi \SI{2,5 \times 10^{-3}}{...}
\sisetup{locale=FR,output-decimal-marker={,},group-minimum-digits=4,parse-numbers=false}
\DeclareSIUnit\ohm{\ensuremath{\symup{\Omega}}} % U+2126 absent de la police
\usepackage{qrcode}
% \degree : commande fréquemment utilisée par les modèles pour un angle ou une
% température en dehors de \SI{}{} (non fournie par les paquets chargés).
\providecommand{\degree}{^{\circ}}
% \qed (fin de démonstration), utilisé par les modèles sans amsthm
\providecommand{\qed}{\hfill$\square$}
% \barre : double barre des valeurs interdites dans un tableau de variations (tkz-tab)
\providecommand{\barre}{\ensuremath{\|}}
% environnement proof (amsthm non chargé) utilisé par les modèles
\ifdefined\proof\else\newenvironment{proof}[1][Démonstration]{\par\noindent\textit{#1.}\ }{\qed\par}\fi
\usepackage{lastpage}
\usepackage{hyperref}
\hypersetup{hidelinks}

% ============================================================
% Palette « Pop » OnBuch+ — mêmes hex que la classe P (plus_ui.dart)
% ============================================================
\definecolor{poporange}{HTML}{F59321}
\definecolor{popdark}{HTML}{E07A0C}
\definecolor{popcream}{HTML}{FFF6E8}
\definecolor{popink}{HTML}{1C1714}
\definecolor{poppurple}{HTML}{7A5AE0}
\definecolor{popgreen}{HTML}{2E9E6A}
\definecolor{poppink}{HTML}{E0457B}
\definecolor{popblue}{HTML}{2F7DE1}
\definecolor{popgold}{HTML}{C98A12}
\definecolor{popmuted}{HTML}{8A7F76}
\definecolor{popline}{HTML}{EADFD2}
\definecolor{popgray}{HTML}{8A7F76}
\colorlet{popgrey}{popgray}
% Alias tolérants (le modèle nomme parfois les couleurs différemment)
\colorlet{popwhite}{white}
\colorlet{popblack}{popink}
\colorlet{popred}{poppink}
\colorlet{popyellow}{popgold}
% Teintes claires pour fonds de tableaux / zones
\colorlet{popblueL}{popblue!10!white}
\colorlet{poporangeL}{poporange!14!white}
\colorlet{popgreenL}{popgreen!12!white}
\colorlet{poppurpleL}{poppurple!12!white}
\colorlet{poppinkL}{poppink!12!white}
\colorlet{popgoldL}{popgold!14!white}

\pagecolor{popcream}

% ============================================================
% Typographie
% ============================================================
\defaultfontfeatures{Ligatures=TeX}
\setmainfont{PlusJakartaSans-Medium.ttf}[Path=../../../fonts/,BoldFont=PlusJakartaSans-Bold.ttf]
% Maths en sans-serif (Fira Math) avec chiffres et lettres droites de Plus
% Jakarta, pour que les formules se fondent dans le texte.
\usepackage[math-style=ISO,bold-style=ISO]{unicode-math}
\setmathfont{FiraMath-Regular.otf}[Path=../../../fonts/]
\setmathfont{PlusJakartaSans-Medium.ttf}[Path=../../../fonts/,range={up/num,up/{Latin,latin},\mathup}]
\usepackage{icomma} % virgule décimale sans espace en mode math
% Caractères mathématiques Unicode absents des polices de texte (Plus Jakarta,
% Archivo Black) : rendus par leur équivalent mathématique, en texte comme en
% formule — sinon ils s'affichent en carré vide (ex. « Arithmétique dans ℤ »).
\usepackage{newunicodechar}
\newunicodechar{ℝ}{\ensuremath{\mathbb{R}}}
\newunicodechar{ℤ}{\ensuremath{\mathbb{Z}}}
\newunicodechar{ℂ}{\ensuremath{\mathbb{C}}}
\newunicodechar{ℕ}{\ensuremath{\mathbb{N}}}
\newunicodechar{ℚ}{\ensuremath{\mathbb{Q}}}
\newunicodechar{𝔻}{\ensuremath{\mathbb{D}}}
\newunicodechar{α}{\ensuremath{\alpha}}
\newunicodechar{β}{\ensuremath{\beta}}
\newunicodechar{γ}{\ensuremath{\gamma}}
\newunicodechar{δ}{\ensuremath{\delta}}
\newunicodechar{ε}{\ensuremath{\varepsilon}}
\newunicodechar{θ}{\ensuremath{\theta}}
\newunicodechar{λ}{\ensuremath{\lambda}}
\newunicodechar{μ}{\ensuremath{\mu}}
\newunicodechar{σ}{\ensuremath{\sigma}}
\newunicodechar{τ}{\ensuremath{\tau}}
\newunicodechar{φ}{\ensuremath{\varphi}}
\newunicodechar{ψ}{\ensuremath{\psi}}
\newunicodechar{ω}{\ensuremath{\omega}}
\newunicodechar{Σ}{\ensuremath{\Sigma}}
% majuscules produites par \MakeUppercase dans les titres (α-aminés -> Α)
\newunicodechar{Α}{\ensuremath{\alpha}}
\newunicodechar{Β}{\ensuremath{\beta}}
\newunicodechar{Γ}{\ensuremath{\Gamma}}
\newunicodechar{Δ}{\ensuremath{\Delta}}
\newunicodechar{Ε}{\ensuremath{\varepsilon}}
\newunicodechar{Θ}{\ensuremath{\Theta}}
\newunicodechar{Λ}{\ensuremath{\Lambda}}
\newunicodechar{Μ}{\ensuremath{\mu}}
\newunicodechar{Τ}{\ensuremath{\tau}}
\newunicodechar{Φ}{\ensuremath{\Phi}}
\newunicodechar{Ψ}{\ensuremath{\Psi}}
\newunicodechar{Ω}{\ensuremath{\Omega}}
\newunicodechar{↦}{\ensuremath{\mapsto}}
\newunicodechar{∈}{\ensuremath{\in}}
\newunicodechar{∉}{\ensuremath{\notin}}
\newunicodechar{∧}{\ensuremath{\wedge}}
\newunicodechar{⇔}{\ensuremath{\Leftrightarrow}}
\newunicodechar{⟺}{\ensuremath{\Longleftrightarrow}}
\newunicodechar{⇒}{\ensuremath{\Rightarrow}}
\newunicodechar{⟹}{\ensuremath{\Longrightarrow}}
\newunicodechar{∘}{\ensuremath{\circ}}
\newunicodechar{∪}{\ensuremath{\cup}}
\newunicodechar{∩}{\ensuremath{\cap}}
\newunicodechar{≡}{\ensuremath{\equiv}}
\newunicodechar{⊂}{\ensuremath{\subset}}
\newunicodechar{⊥}{\ensuremath{\perp}}
\newunicodechar{∃}{\ensuremath{\exists}}
\newunicodechar{∀}{\ensuremath{\forall}}
\newunicodechar{∛}{\ensuremath{\sqrt[3]{\,}}}
\newunicodechar{∜}{\ensuremath{\sqrt[4]{\,}}}
\newunicodechar{⃗}{\ensuremath{{}^{\to}}}
\newunicodechar{ⁿ}{\ifmmode^{n}\else\textsuperscript{\ensuremath{n}}\fi}
\newunicodechar{ˣ}{\ifmmode^{x}\else\textsuperscript{\ensuremath{x}}\fi}
\newunicodechar{ᵇ}{\ifmmode^{b}\else\textsuperscript{\ensuremath{b}}\fi}
\newunicodechar{ʳ}{\ifmmode^{r}\else\textsuperscript{\ensuremath{r}}\fi}
\newunicodechar{ᵘ}{\ifmmode^{u}\else\textsuperscript{\ensuremath{u}}\fi}
\newunicodechar{ʸ}{\ifmmode^{y}\else\textsuperscript{\ensuremath{y}}\fi}
\newunicodechar{ᵃ}{\ifmmode^{a}\else\textsuperscript{\ensuremath{a}}\fi}
\newunicodechar{ᵉ}{\ifmmode^{e}\else\textsuperscript{\ensuremath{e}}\fi}
\newunicodechar{ᵏ}{\ifmmode^{k}\else\textsuperscript{\ensuremath{k}}\fi}
\newunicodechar{ᵖ}{\ifmmode^{p}\else\textsuperscript{\ensuremath{p}}\fi}
\newunicodechar{ᴺ}{\ifmmode^{N}\else\textsuperscript{\ensuremath{N}}\fi}
\newunicodechar{ᶜ}{\ifmmode^{c}\else\textsuperscript{\ensuremath{c}}\fi}
\newunicodechar{ʰ}{\ifmmode^{h}\else\textsuperscript{\ensuremath{h}}\fi}
\newunicodechar{ᵠ}{\ifmmode^{q}\else\textsuperscript{\ensuremath{q}}\fi}
\newunicodechar{ₙ}{\ifmmode_{n}\else\textsubscript{\ensuremath{n}}\fi}
\newunicodechar{ₐ}{\ifmmode_{a}\else\textsubscript{\ensuremath{a}}\fi}
\newunicodechar{ᵢ}{\ifmmode_{i}\else\textsubscript{\ensuremath{i}}\fi}
\newunicodechar{ᵣ}{\ifmmode_{r}\else\textsubscript{\ensuremath{r}}\fi}
\newunicodechar{ₖ}{\ifmmode_{k}\else\textsubscript{\ensuremath{k}}\fi}
\newunicodechar{ᵦ}{\ifmmode_{\beta}\else\textsubscript{\ensuremath{\beta}}\fi}
\newunicodechar{ₓ}{\ifmmode_{x}\else\textsubscript{\ensuremath{x}}\fi}
\newfontfamily\popdisplay{ArchivoBlack-Regular.ttf}[Path=../../../fonts/]
\newfontfamily\poplogo{SpaceGrotesk-Bold.ttf}[Path=../../../fonts/]
\newfontfamily\popsemi{PlusJakartaSans-SemiBold.ttf}[Path=../../../fonts/]
\newfontfamily\popxbold{PlusJakartaSans-ExtraBold.ttf}[Path=../../../fonts/]
\newfontfamily\popxboldls{PlusJakartaSans-ExtraBold.ttf}[Path=../../../fonts/,LetterSpace=3.0]

\setlist[enumerate]{leftmargin=1.7em,itemsep=5pt,topsep=4pt}
\setlist[itemize]{leftmargin=1.7em,itemsep=4pt,topsep=4pt,label={\color{poporange}\textbullet}}
\setlength{\parindent}{0pt}
\setlength{\parskip}{5pt}
\renewcommand{\arraystretch}{1.25}

% pgfplots : style Pop par défaut
\pgfplotsset{
  every axis/.append style={axis line style={popink,line width=1pt},tick style={popink},
    grid style={popline,line width=0.5pt},label style={font=\small},tick label style={font=\footnotesize}},
  popcurve/.style={poporange,line width=1.8pt,no markers,smooth},
}
% Même style disponible dans un \draw TikZ simple (hors pgfplots)
\tikzset{popcurve/.style={poporange,line width=1.8pt}}

% ============================================================
% Pill / squiggle
% ============================================================
\newcommand{\poppill}[3]{%
  \tikz[baseline=-0.55ex]\node[draw=popink,line width=1.2pt,rounded corners=2.6pt,%
    fill=#2,inner xsep=6.5pt,inner ysep=3pt]{%
    \popxboldls\fontsize{7}{7}\selectfont\color{#3}\MakeUppercase{#1}};%
}
\newcommand{\popsquiggle}[1]{%
  \tikz[baseline=-0.4ex]{
    \draw[#1,line width=2.6pt,line cap=round]
      plot [smooth] coordinates {(0,0.09) (0.34,0.01) (0.68,0.21) (1.02,0.01) (1.36,0.10)};
  }%
}
% Mot-clé surligné (marqueur orange sous le texte)
\newcommand{\cle}[1]{{\popxbold\color{popdark}#1}}

% ============================================================
% En-tête / pied
% ============================================================
\pagestyle{fancy}
\fancyhf{}
\renewcommand{\headrulewidth}{0pt}
\renewcommand{\footrulewidth}{0pt}
\lhead{\raisebox{-0.15ex}{\poplogo\fontsize{13}{13}\selectfont\color{popink}OnBuch\color{poporange}+}}
\rhead{\poppill{\DOCMATIERE\ \textbullet\ \DOCNIVEAU\ \textbullet\ Leçon \DOCLECON}{white}{popink}}
\lfoot{\popxbold\fontsize{7.6}{7.6}\selectfont\color{popmuted}\MakeUppercase{Cours signé OnBuch+}\ \textbullet\ {\popsemi\DOCTITRE}}
\rfoot{\tikz[baseline=-0.6ex]\node[rounded corners=8.5pt,fill=popink,minimum height=17pt,minimum width=17pt,inner xsep=5pt,inner sep=0]{\popxbold\fontsize{8}{8}\selectfont\color{popcream}\thepage\,/\,\pageref*{LastPage}};}

% ============================================================
% Titres
% ============================================================
\newcommand{\chaptitre}[2]{%
  \begin{tcolorbox}[enhanced,colback=poporange,colframe=popink,boxrule=2.6pt,arc=11pt,
      drop shadow=popink,left=15pt,right=15pt,top=12pt,bottom=11pt,before skip=3pt,after skip=18pt]
    \poppill{#2}{popink}{popcream}\par\vspace{9pt}
    {\raggedright\hyphenpenalty=10000\exhyphenpenalty=10000\popdisplay\fontsize{22}{24}\selectfont\color{popcream}\MakeUppercase{#1}\par}
  \end{tcolorbox}
}
% Section numérotée : pastille orange avec le numéro + titre Archivo
\newcounter{popsec}
\newcommand{\coursec}[1]{%
  \stepcounter{popsec}\setcounter{popsub}{0}%
  \par\vspace{16pt}\noindent
  \begin{minipage}{\linewidth}
  \tikz[baseline=-0.7ex]\node[draw=popink,line width=1.6pt,fill=poporange,rounded corners=4pt,minimum size=22pt,inner sep=0]{\popdisplay\fontsize{12}{12}\selectfont\color{popcream}\thepopsec};%
  \hspace{8pt}{\popdisplay\fontsize{15}{17}\selectfont\color{popink}\MakeUppercase{#1}}\par
  \vspace{4pt}\noindent{\color{poporange}\rule{36pt}{3.6pt}}
  \end{minipage}\par\nopagebreak\vspace{8pt}
}
\newcounter{popsub}
\newcommand{\courssub}[1]{%
  \stepcounter{popsub}%
  \par\vspace{9pt}\noindent{\popxbold\fontsize{11.5}{13}\selectfont\color{popink}{\color{poporange}\thepopsec.\thepopsub}\hspace{6pt}#1}\par\nopagebreak\vspace{4pt}
}

% ============================================================
% Boîtes pédagogiques — même recette : fond blanc, bordure épaisse colorée,
% ombre dure encre, badge plein. Titre optionnel [#1].
% ============================================================
\tcbset{popbase/.style={breakable,enhanced,colback=white,boxrule=2.3pt,arc=9pt,
  shadow={3pt}{-3pt}{0pt}{popink},left=14pt,right=14pt,top=11pt,bottom=11pt,before skip=11pt,after skip=15pt,
  fontupper=\popsemi\fontsize{10.6}{14.5}\selectfont\color{popink},
  fontlower=\popsemi\fontsize{10.6}{14.5}\selectfont\color{popink}}}
% Coche dessinée (le glyphe ✓ n'existe pas dans les polices du document)
\newcommand{\popcheck}[1]{\tikz[baseline=-0.5ex]\draw[#1,line width=1.6pt,line cap=round,line join=round](-0.13,0.0)--(-0.04,-0.09)--(0.13,0.1);}
\newcommand{\popboxhead}[3]{\poppill{#1}{#2}{white}\ifx\relax#3\relax\else\hspace{6pt}{\popxbold\fontsize{10.6}{12}\selectfont\color{popink}#3}\fi\par\vspace{7pt}}

\newtcolorbox{aretenir}[1][]{popbase,colframe=popgreen,before upper={\popboxhead{À retenir}{popgreen}{#1}}}
\newtcolorbox{exemplebox}[1][]{popbase,colframe=poporange,before upper={\popboxhead{Exemple}{poporange}{#1}}}
\newtcolorbox{definition}[1][]{popbase,colframe=popblue,before upper={\popboxhead{Définition}{popblue}{#1}}}
\newtcolorbox{propriete}[1][]{popbase,colframe=poppurple,before upper={\popboxhead{Propriété}{poppurple}{#1}}}
\newtcolorbox{methode}[1][]{popbase,colframe=popgold,before upper={\popboxhead{Méthode}{popgold}{#1}}}
\newtcolorbox{attention}[1][]{popbase,colframe=poppink,before upper={\popboxhead{Attention}{poppink}{#1}}}
\newtcolorbox{experience}[1][]{popbase,colframe=popblue,colback=popblueL,before upper={\popboxhead{Expérience}{popblue}{#1}}}
% « Le savais-tu ? » : carte violette pleine (contexte, culture, Cameroun)
\newtcolorbox{savaistu}[1][]{popbase,colback=poppurple,colframe=popink,
  fontupper=\popsemi\fontsize{10.4}{14.2}\selectfont\color{white},
  before upper={\poppill{Le savais-tu\,?}{white}{poppurple}\ifx\relax#1\relax\else\hspace{6pt}{\popxbold\color{white}#1}\fi\par\vspace{7pt}}}
% Exercice résolu : énoncé en haut, \tcblower puis la solution sur fond crème
\newcounter{popexo}\newcounter{popexr}
\newtcolorbox{exoresolu}[1][]{popbase,colframe=poporange,colbacklower=popcream,
  segmentation style={popink,line width=1.2pt,dashed},
  before upper={\stepcounter{popexr}\popboxhead{Exercice résolu \thepopexr}{poporange}{#1}},
  before lower={\poppill{Solution}{popink}{popcream}\par\vspace{6pt}}}
% Exercice (non résolu en place) — la correction est dans la partie Corrigés
\newtcolorbox{exercice}[1][]{popbase,colframe=popink,
  before upper={\stepcounter{popexo}\popboxhead{Exercice \thepopexo}{popink}{#1}}}
% Corrigé
\newtcolorbox{corrige}[1][]{popbase,colframe=popgreen,colback=popgreenL,
  before upper={\popboxhead{Corrigé}{popgreen}{#1}}}
% Objectifs / prérequis / bilan
\newtcolorbox{objectifs}{popbase,colframe=popink,colback=poporangeL,
  before upper={\popboxhead{Objectifs de la leçon}{popink}{}}}
\newtcolorbox{prerequis}{popbase,colframe=popink,colback=popblueL,
  before upper={\popboxhead{Prérequis}{popblue}{}}}
\newtcolorbox{bilan}{popbase,colframe=popink,colback=white,boxrule=2.8pt,
  before upper={\popboxhead{Fiche bilan}{poporange}{L'essentiel à connaître pour l'examen}}}
% Figure encadrée avec légende
\newtcolorbox{popfigure}[1][]{enhanced,breakable=false,colback=white,colframe=popline,boxrule=1.4pt,arc=8pt,
  left=10pt,right=10pt,top=10pt,bottom=8pt,before skip=10pt,after skip=12pt,halign=center,
  lower separated=false,halign lower=center,
  fontlower=\popsemi\fontsize{9}{11.5}\selectfont\color{popmuted},
  #1}
\newcommand{\legende}[1]{\par\vspace{6pt}{\popsemi\fontsize{9}{11.5}\selectfont\color{popmuted}#1}}

% ============================================================
% Signature OnBuch+
% ============================================================
\newcommand{\poplogoinline}[1]{{\poplogo\fontsize{#1}{#1}\selectfont\color{popink}OnBuch\color{poporange}+}}
\newcommand{\signatureonbuch}{%
  \begin{tcolorbox}[enhanced,colback=popink,colframe=popink,boxrule=2.2pt,arc=10pt,
      drop shadow=poporange,left=14pt,right=14pt,top=10pt,bottom=10pt,before skip=0pt,after skip=0pt]
    \tikz[baseline=-0.7ex]\node[circle,fill=popgreen,draw=popcream,line width=1.3pt,minimum size=22pt,inner sep=0]{\popcheck{white}};%
    \hspace{9pt}\begin{minipage}[c]{0.62\linewidth}
      {\popxbold\fontsize{12}{13}\selectfont\color{popcream}Signé\ \textbullet\ {\poplogo OnBuch\color{poporange}+}}\par\vspace{2pt}
      {\raggedright\popsemi\fontsize{8}{10.5}\selectfont\color{popline}\MakeUppercase{Équipe pédagogique OnBuch+}\\\MakeUppercase{Conforme au programme officiel MINESEC}\par}
    \end{minipage}\hfill
    {\poplogo\fontsize{22}{22}\selectfont\color{popcream}OnBuch\color{poporange}+}
  \end{tcolorbox}%
}

% ============================================================
% Page de garde
% #1 titre · #2 matière · #3 niveau · #4 module · #5 n° leçon · #6 durée ·
% #7 description / objectifs (texte ou itemize)
% ============================================================
\newcommand{\pagegarde}[7]{%
  \thispagestyle{empty}
  \newgeometry{a4paper,margin=1.7cm,top=1.6cm,bottom=1.4cm}%
  \begin{tikzpicture}[remember picture,overlay]
    \fill[poporange!18] ([xshift=-3.2cm,yshift=-3.6cm]current page.north east) circle (4.6cm);
    \fill[poppurple!14] ([xshift=2.4cm,yshift=7.6cm]current page.south west) circle (3.8cm);
  \end{tikzpicture}%
  {\poplogo\fontsize{30}{30}\selectfont\color{popink}OnBuch\color{poporange}+}\hfill
  \raisebox{0.6ex}{\poppill{Cours signé \textbullet\ Premium}{popink}{popcream}}\par
  \vspace{1pt}\popsquiggle{poppurple}\par
  \vspace{8pt}
  \begin{tcolorbox}[enhanced,colback=poporange,colframe=popink,boxrule=2.8pt,arc=13pt,
      drop shadow=popink,left=18pt,right=18pt,top=12pt,bottom=13pt,after skip=10pt]
    \poppill{#2 \textbullet\ #3}{popink}{popcream}\hspace{5pt}\poppill{#4}{white}{popink}\par\vspace{8pt}
    {\popdisplay\fontsize{14}{14}\selectfont\color{popink}LEÇON #5}\par\vspace{4pt}
    {\raggedright\hyphenpenalty=10000\exhyphenpenalty=10000\popdisplay\fontsize{25}{27}\selectfont\color{popcream}\MakeUppercase{#1}\par}
    \vspace{8pt}
    {\popxbold\fontsize{9.5}{9.5}\selectfont\color{popink}\MakeUppercase{Durée conseillée : #6}}
  \end{tcolorbox}
  \begin{tcolorbox}[enhanced,colback=white,colframe=popink,boxrule=2.2pt,arc=10pt,
      drop shadow=popink,left=16pt,right=16pt,top=10pt,bottom=10pt,after skip=8pt]
    \poppill{Dans cette leçon}{poppurple}{white}\par\vspace{5pt}
    \popsemi\fontsize{9.3}{12.6}\selectfont\color{popink}#7
  \end{tcolorbox}
  \noindent\begin{minipage}[t]{0.487\linewidth}
    \begin{tcolorbox}[enhanced,colback=popgreen,colframe=popink,boxrule=2.2pt,arc=10pt,
        drop shadow=popink,left=11pt,right=11pt,top=10pt,bottom=10pt,height=3.1cm,valign=center]
      \tcbox[colback=white,colframe=popink,boxrule=1.3pt,arc=5pt,boxsep=0pt,nobeforeafter,
        left=3pt,right=3pt,top=3pt,bottom=3pt,valign=center]{\qrcode[height=1.9cm]{https://play.google.com/store/apps/details?id=com.luvvix.onbuch&hl=fr}}%
      \hspace{8pt}\begin{minipage}[c]{0.5\linewidth}
        {\popdisplay\fontsize{11}{12.5}\selectfont\color{white}\MakeUppercase{Télécharge OnBuch}}\par\vspace{4pt}
        {\raggedright\popsemi\fontsize{8.4}{11}\selectfont\color{white}Gratuit \textbullet\ Google Play\\Tuteur IA, annales, concours, résultats\par}
      \end{minipage}
    \end{tcolorbox}
  \end{minipage}\hfill
  \begin{minipage}[t]{0.487\linewidth}
    \begin{tcolorbox}[enhanced,colback=poppurple,colframe=popink,boxrule=2.2pt,arc=10pt,
        drop shadow=popink,left=11pt,right=11pt,top=10pt,bottom=10pt,height=3.1cm,valign=center]
      \tikz[baseline=-0.7ex]\node[circle,fill=white,draw=popink,line width=1.3pt,minimum size=1.6cm,inner sep=0]{\popdisplay\fontsize{24}{24}\selectfont\color{poppurple}+};%
      \hspace{8pt}\begin{minipage}[c]{0.6\linewidth}
        {\popdisplay\fontsize{11}{12.5}\selectfont\color{white}\MakeUppercase{Passe à OnBuch+}}\par\vspace{4pt}
        {\raggedright\popsemi\fontsize{8.4}{11}\selectfont\color{white}Tous les cours signés, Tuteur IA illimité, fiches PDF — onglet OnBuch+ de l'app\par}
      \end{minipage}
    \end{tcolorbox}
  \end{minipage}
  \vspace{8pt}
  \popcontenu
  \vfill
  \signatureonbuch
  \restoregeometry
  \newpage
}

% Rangée « Ce cours contient » (page de garde)
\newcommand{\poptile}[3]{% #1 couleur #2 gros texte #3 légende
  \begin{tcolorbox}[enhanced,colback=#1,colframe=popink,boxrule=2pt,arc=9pt,shadow={2.5pt}{-2.5pt}{0pt}{popink},
      left=6pt,right=6pt,top=9pt,bottom=9pt,halign=center,valign=center,height=1.75cm,width=\linewidth,nobeforeafter]
    {\popdisplay\fontsize{12.5}{13}\selectfont\color{white}\MakeUppercase{#2}}\par\vspace{3pt}
    {\centering\popsemi\fontsize{7.8}{9.5}\selectfont\color{white}#3\par}
  \end{tcolorbox}}
\newcommand{\popcontenu}{%
  \par\noindent{\popxboldls\fontsize{7.5}{7.5}\selectfont\color{popmuted}CE COURS SIGNÉ CONTIENT}\par\vspace{7pt}
  \noindent\begin{minipage}[t]{0.235\linewidth}\poptile{popblue}{Cours}{complet, illustré, pas à pas}\end{minipage}\hfill
  \begin{minipage}[t]{0.235\linewidth}\poptile{poporange}{Méthodes}{et exercices résolus}\end{minipage}\hfill
  \begin{minipage}[t]{0.235\linewidth}\poptile{poppink}{Exercices}{type examen + corrigés}\end{minipage}\hfill
  \begin{minipage}[t]{0.235\linewidth}\poptile{popgreen}{Fiche bilan}{l'essentiel à retenir}\end{minipage}\par}

% Sommaire
\newtcolorbox{sommaire}{enhanced,colback=white,colframe=popink,boxrule=2.2pt,arc=10pt,
  drop shadow=popink,left=18pt,right=18pt,top=13pt,bottom=13pt,before skip=6pt,after skip=20pt,
  before upper={\poppill{Sommaire}{poppurple}{white}\par\vspace{9pt}\popsemi\fontsize{10.6}{14.5}\selectfont\color{popink}}}

% Signature de fin de cours — poussée tout en bas de la dernière page
\newcommand{\findecours}{%
  \par\vfill\nopagebreak
  \begin{center}\popsquiggle{poporange}\end{center}\vspace{4pt}
  \signatureonbuch
}
\AtBeginDocument{\def\llbracket{\mathopen{[\mkern-3.5mu[}}\def\rrbracket{\mathclose{]\mkern-3.5mu]}}} % intervalles d'entiers (glyphe ⟦ absent de la police)
% Glyphes absents de Fira Math : redéfinis à partir de glyphes disponibles
\AtBeginDocument{%
  \def\setminus{\mathbin{\backslash}}\let\smallsetminus\setminus
  \def\vdots{\mathord{\vbox{\baselineskip3.2pt\lineskiplimit0pt\kern4pt\hbox{.}\hbox{.}\hbox{.}}}}%
  \def\ddots{\mathinner{\mkern1mu\raise6.5pt\hbox{.}\mkern2mu\raise3.6pt\hbox{.}\mkern2mu\raise0.7pt\hbox{.}\mkern1mu}}%
  \def\triangle{\mathord{\scalebox{0.9}{$\symup{\Delta}$}}}%
  \def\nabla{\mathord{\raisebox{\depth}{\scalebox{1}[-1]{$\symup{\Delta}$}}}}%
  \def\checkmark{\popcheck{popdark}}%
  \def\star{\mathbin{\ast}}\def\bigstar{\mathord{\ast}}%
  \def\diamond{\mathbin{\scalebox{0.6}{$\Diamond$}}}%
}

%% FIGFIX-BEGIN v5 — correctifs globaux des figures (docs/onbuch-generation/tools/figfix)
\usepackage{adjustbox}\usepackage{etoolbox}\tcbuselibrary{raster}
% toute figure TikZ/pgfplots plus large que la ligne est réduite à la largeur disponible
\BeforeBeginEnvironment{tikzpicture}{\begin{adjustbox}{max width=\linewidth}}
\AfterEndEnvironment{tikzpicture}{\end{adjustbox}}
% légendes pgfplots lisibles par-dessus les courbes
\pgfplotsset{every axis/.append style={legend style={fill=white,fill opacity=0.92,text opacity=1,draw=black!15,inner sep=3pt}}}
% étiquettes d'arêtes (syntaxe edge["..."]) avec fond blanc
\tikzset{every edge quotes/.append style={fill=white,inner sep=1.2pt}}
%% FIGFIX-END
\begin{document}

\pagegarde{Suites numériques}{Mathématiques}{Tle C}{Relations et opérations fondamentales dans l'ensemble des nombres réels}{04}{8 h}{Cette leçon introduit le raisonnement par récurrence et les outils d'étude des suites : monotonie, bornitude, convergence des suites monotones bornées. Elle aboutit à l'étude complète des suites récurrentes uₙ₊₁ = f(uₙ) et à l'encadrement numérique de leur limite par l'inégalité des accroissements finis, compétence phare du Bac C, D, E et TI.
\vspace{4pt}
\begin{multicols}{2}\small
\begin{itemize}
\item À la fin de cette leçon, je sais démontrer une propriété dépendant de n par récurrence, en rédigeant rigoureusement l'initialisation, l'hérédité et la conclusion.
\item À la fin de cette leçon, je sais déterminer si une suite est majorée, minorée ou bornée.
\item À la fin de cette leçon, je sais étudier la monotonie d'une suite par différence, par quotient ou par étude de fonction.
\item À la fin de cette leçon, je sais appliquer le théorème de convergence des suites monotones bornées pour justifier qu'une suite converge.
\item À la fin de cette leçon, je sais étudier une suite définie par uₙ₊₁ = f(uₙ) : sens de variation, bornitude, convergence et calcul de la limite.
\item À la fin de cette leçon, je sais utiliser l'inégalité des accroissements finis pour majorer |uₙ − ℓ| et obtenir une valeur approchée de la limite.
\end{itemize}\end{multicols}}

\begin{sommaire}
\begin{multicols}{2}
\begin{enumerate}
\item Raisonnement par récurrence
\item Suites majorées, minorées, bornées
\item Suites monotones et théorème de convergence monotone
\item Suites récurrentes uₙ₊₁ = f(uₙ) : étude de la convergence
\item Valeur approchée de la limite par l'inégalité des accroissements finis
\item Synthèse : méthode complète d'étude d'une suite au Bac
\item Définition rigoureuse de la limite d'une suite (limite par epsilon)
\item Activité d'intégration
\item Méthodes et exercices résolus
\item Exercices
\item Corrigés des exercices
\item Fiche bilan
\end{enumerate}
\end{multicols}
\end{sommaire}

\begin{objectifs}
\begin{itemize}
\item À la fin de cette leçon, je sais démontrer une propriété dépendant de n par récurrence, en rédigeant rigoureusement l'initialisation, l'hérédité et la conclusion.
\item À la fin de cette leçon, je sais déterminer si une suite est majorée, minorée ou bornée.
\item À la fin de cette leçon, je sais étudier la monotonie d'une suite par différence, par quotient ou par étude de fonction.
\item À la fin de cette leçon, je sais appliquer le théorème de convergence des suites monotones bornées pour justifier qu'une suite converge.
\item À la fin de cette leçon, je sais étudier une suite définie par uₙ₊₁ = f(uₙ) : sens de variation, bornitude, convergence et calcul de la limite.
\item À la fin de cette leçon, je sais utiliser l'inégalité des accroissements finis pour majorer |uₙ − ℓ| et obtenir une valeur approchée de la limite.
\item À la fin de cette leçon, je sais déterminer un rang N à partir duquel uₙ approche la limite à une précision donnée.
\item À la fin de cette leçon, je sais rédiger une étude complète de suite récurrente dans le format attendu au Baccalauréat.
\end{itemize}
\end{objectifs}

\begin{prerequis}
\begin{itemize}
\item Notion de suite, suites arithmétiques et géométriques (Première)
\item Dérivation et sens de variation d'une fonction (début de Terminale)
\item Inégalités, valeur absolue et encadrements dans ℝ
\item Limites simples de suites (qⁿ, 1/n) vues en Première
\item Inégalité des accroissements finis pour les fonctions dérivables
\end{itemize}
\end{prerequis}

\newpage

\coursec{Raisonnement par récurrence}

\courssub{Une situation concrète : la tontine numérique de Maman Aïcha}

Maman Aïcha, commerçante au marché Mokolo à Yaoundé, dépose \SI{100000}{} FCFA dans une tontine numérique qui rapporte $2\,\%$ par mois. À la fin de chaque mois, elle ajoute \SI{5000}{} FCFA. Deux questions la taraudent : combien aura-t-elle après un an ? Sa somme dépassera-t-elle \SI{200000}{} FCFA ?

Notons $u_n$ le montant, en FCFA, au bout de $n$ mois. On a $u_0 = 100\,000$ et, chaque mois, le capital est multiplié par $1{,}02$ puis augmenté de $5\,000$ :
\[ u_{n+1} = 1{,}02\,u_n + 5\,000. \]
Calculer mois par mois est fastidieux. Surtout : comment \emph{prouver} que, par exemple, $u_n \leq 350\,000$ pour tout $n \leq 12$, sans vérifier les douze cas un par un ? Comment être certain qu'une propriété vraie au mois 1, au mois 2, au mois 3\ldots{} le reste \emph{toujours} ?

C'est précisément l'objet du \cle{raisonnement par récurrence} : un mode de démonstration qui permet d'établir qu'une propriété dépendant d'un entier naturel $n$ est vraie pour \emph{tous} les entiers à partir d'un certain rang. C'est l'outil fondamental de tout ce chapitre : il servira à étudier les suites définies par une relation de récurrence, à prouver des encadrements, des monotonies, des divisibilités.

\textbf{Problématique.} Comment démontrer rigoureusement qu'une propriété $P(n)$ est vraie pour une infinité d'entiers $n$, en ne faisant qu'un nombre fini de vérifications ?

\courssub{L'intuition : l'effet domino}

Imagine une rangée infinie de dominos debout, comme ceux qu'alignent les enfants sur le perron. Deux conditions suffisent pour que \emph{tous} les dominos tombent :

\begin{itemize}
\item le \textbf{premier} domino tombe ;
\item chaque domino, en tombant, \textbf{fait tomber le suivant}.
\end{itemize}

Si l'une des deux conditions manque, la chute s'arrête : si le premier domino ne tombe pas, rien ne démarre ; si un domino est trop éloigné du suivant, la chaîne se brise. Le raisonnement par récurrence, c'est exactement cela, traduit en langage mathématique.

\begin{popfigure}
\begin{tikzpicture}[scale=0.9, every node/.style={transform shape}]
% Dominos debout puis tombés
\foreach \i/\ang in {0/0, 1/0, 2/25, 3/55, 4/80}{
  \begin{scope}[shift={(\i*1.6,0)}, rotate=\ang]
    \draw[fill=popblueL, draw=popblue, thick, rounded corners=1pt] (0,0) rectangle (0.35,1.6);
  \end{scope}
}
% Premier domino : étiquette initialisation
\node[below, popdark, align=center] at (0.17,-0.15) {\small domino $n_0$};
\node[above, popdark, align=center] at (0.17,1.75) {\small \textbf{Initialisation}};
% Flèche de transmission
\draw[-{Stealth[length=3mm]}, poporange, very thick] (2.1,1.9) .. controls (3.2,2.5) .. (4.4,2.0);
\node[poporange, align=center] at (3.3,2.75) {\small \textbf{Hérédité} : chaque domino\\ \small fait tomber le suivant};
% Points de suspension
\node[popmuted] at (8.6,0.8) {\Large $\cdots$};
\node[popmuted, align=center] at (8.6,-0.4) {\small et ainsi de suite\\ \small jusqu'à l'infini};
\end{tikzpicture}
\legende{L'effet domino : si le premier domino tombe (initialisation) et si chaque domino entraîne le suivant (hérédité), alors tous tombent.}
\end{popfigure}

\courssub{Le principe de récurrence}

\begin{definition}[Axiome de récurrence]
Soit $P(n)$ une propriété dépendant d'un entier naturel $n$, et soit $n_0$ un entier naturel. Si les deux conditions suivantes sont satisfaites :
\begin{enumerate}
\item \textbf{Initialisation} : $P(n_0)$ est vraie ;
\item \textbf{Hérédité} : pour tout entier $n \geq n_0$, si $P(n)$ est vraie alors $P(n+1)$ est vraie ;
\end{enumerate}
alors la propriété $P(n)$ est vraie pour \textbf{tout} entier $n \geq n_0$.
\end{definition}

Ce principe n'est pas un théorème que l'on démontre : c'est un \cle{axiome}, c'est-à-dire une propriété admise, qui fait partie de la définition même de l'ensemble $\mathbb{N}$ des entiers naturels. On l'accepte parce qu'il traduit fidèlement notre intuition de la « file infinie » des entiers : $0, 1, 2, 3, \ldots$

\begin{attention}[Vocabulaire et quantificateurs]
Dans l'hérédité, il faut prouver l'implication $P(n) \Rightarrow P(n+1)$ pour \emph{tout} entier $n \geq n_0$ : on raisonne donc sur un entier $n$ \emph{quelconque} (fixé mais arbitraire). Il est interdit de choisir une valeur particulière comme $n = 3$ : cela ne prouverait l'implication que pour ce seul rang.
\end{attention}

\courssub{La rédaction type en trois étapes}

\begin{methode}[Rédiger une démonstration par récurrence]
Pour démontrer que $P(n)$ est vraie pour tout $n \geq n_0$ :
\begin{enumerate}
\item \textbf{Initialisation.} On vérifie directement que $P(n_0)$ est vraie (calcul explicite au rang $n_0$).
\item \textbf{Hérédité.} On fixe un entier $n \geq n_0$ quelconque. On \emph{suppose} que $P(n)$ est vraie : c'est l'\cle{hypothèse de récurrence} (HR). Sous cette hypothèse, on \emph{démontre} que $P(n+1)$ est vraie. L'hypothèse de récurrence doit être \emph{utilisée} dans la preuve.
\item \textbf{Conclusion.} On invoque le principe de récurrence : « Par récurrence, $P(n)$ est vraie pour tout entier $n \geq n_0$. »
\end{enumerate}
\end{methode}

\begin{popfigure}
\begin{tikzpicture}[
  etape/.style={rectangle, rounded corners=4pt, draw=popblue, fill=popblueL, thick, text width=4.2cm, align=center, minimum height=1.5cm},
  fleche/.style={-{Stealth[length=3mm]}, very thick, poporange}]
\node[etape, align=center] (e1) at (0,0) {\textbf{Étape 1 : Initialisation}\\ Vérifier $P(n_0)$\\ \emph{(le premier domino tombe)}};
\node[etape, align=center] (e2) at (6,0) {\textbf{Étape 2 : Hérédité}\\ Supposer $P(n)$ (HR),\\ prouver $P(n+1)$};
\node[etape, draw=popgreen, fill=popgreenL, align=center] (e3) at (12,0) {\textbf{Étape 3 : Conclusion}\\ $P(n)$ vraie pour tout\\ $n \geq n_0$};
\draw[fleche] (e1) -- (e2);
\draw[fleche] (e2) -- (e3);
\end{tikzpicture}
\legende{Les trois étapes obligatoires d'une démonstration par récurrence.}
\end{popfigure}

\courssub{Exemple 1 : la somme des $n$ premiers entiers}

\begin{exemplebox}[Démonstration complète : $1 + 2 + \cdots + n = \dfrac{n(n+1)}{2}$]
Montrons par récurrence que pour tout entier $n \geq 1$ :
\[ P(n) : \quad 1 + 2 + \cdots + n = \frac{n(n+1)}{2}. \]

\textbf{Initialisation.} Pour $n = 1$ : le membre de gauche vaut $1$, et le membre de droite vaut $\dfrac{1 \times 2}{2} = 1$. Donc $P(1)$ est vraie.

\textbf{Hérédité.} Soit $n \geq 1$ un entier quelconque. Supposons $P(n)$ vraie, c'est-à-dire :
\[ 1 + 2 + \cdots + n = \frac{n(n+1)}{2} \quad \text{(HR)}. \]
Montrons $P(n+1)$, c'est-à-dire $1 + 2 + \cdots + n + (n+1) = \dfrac{(n+1)(n+2)}{2}$. On a :
\begin{align*}
1 + 2 + \cdots + n + (n+1) &= \frac{n(n+1)}{2} + (n+1) \quad \text{(par HR)}\\
&= (n+1)\left(\frac{n}{2} + 1\right)\\
&= (n+1)\,\frac{n+2}{2} = \frac{(n+1)(n+2)}{2}.
\end{align*}
Donc $P(n+1)$ est vraie.

\textbf{Conclusion.} Par récurrence, pour tout entier $n \geq 1$ : $1 + 2 + \cdots + n = \dfrac{n(n+1)}{2}$.
\end{exemplebox}

Vérifions numériquement : pour $n = 100$, la formule donne $\dfrac{100 \times 101}{2} = 5\,050$. C'est la célèbre réponse que le jeune Gauss aurait trouvée instantanément à l'école primaire.

\courssub{Exemple 2 : l'inégalité de Bernoulli}

\begin{propriete}[Inégalité de Bernoulli]
Pour tout réel $a > 0$ et tout entier $n \geq 1$ :
\[ (1 + a)^n \geq 1 + na. \]
\end{propriete}

\begin{experience}[Démonstration guidée de l'inégalité de Bernoulli]
Notons $P(n) : (1+a)^n \geq 1 + na$, pour $a > 0$ fixé.

\textbf{Étape 1 — Initialisation.} Pour $n = 1$ : $(1+a)^1 = 1 + a = 1 + 1 \times a$. L'inégalité est même une égalité : $P(1)$ est vraie.

\textbf{Étape 2 — Hérédité.} Soit $n \geq 1$ quelconque. Supposons $(1+a)^n \geq 1 + na$ (HR). Comme $1 + a > 0$, on peut multiplier les deux membres par $1 + a$ sans changer le sens de l'inégalité :
\begin{align*}
(1+a)^{n+1} &= (1+a)^n (1+a) \geq (1 + na)(1 + a) \quad \text{(par HR)}\\
&= 1 + a + na + na^2 = 1 + (n+1)a + na^2.
\end{align*}
Or $na^2 \geq 0$, donc $(1+a)^{n+1} \geq 1 + (n+1)a$ : c'est $P(n+1)$.

\textbf{Étape 3 — Conclusion.} Par récurrence, $(1+a)^n \geq 1 + na$ pour tout $n \geq 1$.
\end{experience}

\begin{exemplebox}[Application numérique]
Avec $a = 0{,}02$ (les $2\,\%$ de la tontine de Maman Aïcha) et $n = 12$ : $(1{,}02)^{12} \geq 1 + 12 \times 0{,}02 = 1{,}24$. Son capital initial, sans les versements, serait donc multiplié par au moins $1{,}24$ en un an. (En réalité $(1{,}02)^{12} \approx 1{,}268$ : l'inégalité donne une minoration rapide et sûre.)
\end{exemplebox}

\courssub{Exemple 3 : prouver un encadrement pour une suite récurrente}

C'est l'usage le plus fréquent au Baccalauréat : montrer qu'une suite définie par $u_{n+1} = f(u_n)$ reste dans un intervalle.

\begin{exoresolu}[Encadrement d'une suite récurrente]
Soit $(u_n)$ la suite définie par $u_0 = 1$ et, pour tout $n \in \mathbb{N}$, $u_{n+1} = \sqrt{u_n + 2}$. Démontrer que pour tout entier naturel $n$ : $0 \leq u_n \leq 2$.
\tcblower
Notons $P(n) : 0 \leq u_n \leq 2$.

\textbf{Initialisation.} $u_0 = 1$, et $0 \leq 1 \leq 2$ : $P(0)$ est vraie.

\textbf{Hérédité.} Soit $n \in \mathbb{N}$ quelconque. Supposons $0 \leq u_n \leq 2$ (HR). Alors :
\begin{align*}
0 \leq u_n \leq 2 &\implies 2 \leq u_n + 2 \leq 4\\
&\implies \sqrt{2} \leq \sqrt{u_n + 2} \leq \sqrt{4} \quad \text{(la fonction racine est croissante sur } \mathbb{R}_+\text{)}\\
&\implies 0 \leq \sqrt{2} \leq u_{n+1} \leq 2.
\end{align*}
Donc $0 \leq u_{n+1} \leq 2$ : $P(n+1)$ est vraie.

\textbf{Conclusion.} Par récurrence, pour tout $n \in \mathbb{N}$, $0 \leq u_n \leq 2$ : la suite $(u_n)$ est \cle{bornée}.
\end{exoresolu}

Remarque la stratégie : pour passer de $u_n$ à $u_{n+1}$, on applique à l'encadrement de l'hypothèse de récurrence les mêmes opérations que celles qui définissent $f$, en veillant au sens des inégalités (croissance des fonctions utilisées).

\courssub{Complément Bac : récurrence et divisibilité}

\begin{exoresolu}[Divisibilité : $3$ divise $4^n - 1$]
Démontrer que pour tout entier naturel $n$, l'entier $4^n - 1$ est divisible par $3$.
\tcblower
Notons $P(n)$ : « il existe un entier $k \in \mathbb{Z}$ tel que $4^n - 1 = 3k$ ».

\textbf{Initialisation.} Pour $n = 0$ : $4^0 - 1 = 0 = 3 \times 0$. Donc $P(0)$ est vraie (avec $k = 0$).

\textbf{Hérédité.} Soit $n \in \mathbb{N}$ quelconque. Supposons qu'il existe $k \in \mathbb{Z}$ tel que $4^n - 1 = 3k$ (HR), c'est-à-dire $4^n = 3k + 1$. Alors :
\begin{align*}
4^{n+1} - 1 &= 4 \times 4^n - 1 = 4(3k + 1) - 1 \quad \text{(par HR)}\\
&= 12k + 3 = 3(4k + 1).
\end{align*}
Comme $4k + 1 \in \mathbb{Z}$, on a bien $4^{n+1} - 1$ divisible par $3$ : $P(n+1)$ est vraie.

\textbf{Conclusion.} Par récurrence, $4^n - 1$ est divisible par $3$ pour tout $n \in \mathbb{N}$.
\end{exoresolu}

L'astuce à retenir : traduire « $a$ divise $b$ » par « il existe $k \in \mathbb{Z}$ tel que $b = ak$ », puis faire apparaître l'hypothèse de récurrence dans l'expression au rang $n+1$.

\courssub{Les erreurs classiques à bannir}

\begin{attention}[Pièges fréquents du raisonnement par récurrence]
\begin{enumerate}
\item \textbf{Oublier l'initialisation.} L'hérédité seule ne suffit pas ! Contre-exemple : soit $P(n) : n = n + 1$. Si $P(n)$ était vraie, alors $n + 1 = n + 2$, donc $P(n+1)$ serait vraie : l'hérédité « fonctionne ». Pourtant $P(n)$ est fausse pour tout $n$, car aucun rang initial n'est vérifié. Sans premier domino, rien ne tombe.
\item \textbf{Raisonnement circulaire.} Il est interdit d'utiliser $P(n+1)$ pour prouver $P(n+1)$. On part de $P(n)$ (l'hypothèse) et on \emph{construit} $P(n+1)$. Écrire « supposons $P(n+1)$ vraie » est une faute grave de logique.
\item \textbf{Ne jamais utiliser l'hypothèse de récurrence.} Si ta preuve de $P(n+1)$ ne fait pas appel à $P(n)$, ce n'est pas une récurrence (c'est peut-être une preuve directe, à rédiger alors sans récurrence).
\item \textbf{Initialiser au mauvais rang.} Si la propriété n'est demandée que pour $n \geq 3$, l'initialisation doit se faire en $n_0 = 3$, pas en $0$. Vérifier $P(0)$ ne dispense pas de vérifier $P(3)$.
\item \textbf{Rédaction floue.} Annoncer explicitement la propriété $P(n)$, écrire « supposons $P(n)$ vraie pour un entier $n \geq n_0$ quelconque », et conclure en citant le principe de récurrence. Au Bac, la structure de la rédaction est notée autant que le calcul.
\end{enumerate}
\end{attention}

\begin{savaistu}[Les axiomes de Peano]
Le principe de récurrence n'est pas une simple astuce de calcul : il est l'un des \textbf{axiomes de Peano} (Giuseppe Peano, mathématicien italien, 1889), qui définissent formellement l'ensemble $\mathbb{N}$ des entiers naturels. Autrement dit, la récurrence est gravée dans la définition même des nombres $0, 1, 2, 3, \ldots$ Bien avant Peano, on en trouve des traces chez Pascal (\emph{Traité du triangle arithmétique}, 1654) et même chez le mathématicien perse Al-Karaji vers l'an 1000. Quant à la formule $1 + 2 + \cdots + n$, la légende raconte que Gauss, âgé d'une dizaine d'années, aurait ébloui son instituteur en regroupant les termes deux à deux : $(1 + 100) + (2 + 99) + \cdots = 50 \times 101 = 5\,050$. La récurrence, elle, garantit que cette formule reste vraie pour \emph{tous} les entiers, même ceux que personne ne calculera jamais.
\end{savaistu}

\begin{aretenir}[L'essentiel de la section]
\begin{itemize}
\item Le raisonnement par récurrence prouve qu'une propriété $P(n)$ est vraie pour tout $n \geq n_0$ grâce à deux ingrédients : \textbf{initialisation} ($P(n_0)$ vraie) et \textbf{hérédité} ($P(n) \Rightarrow P(n+1)$ pour tout $n \geq n_0$).
\item La rédaction comporte \textbf{trois étapes obligatoires}, et l'hypothèse de récurrence doit être réellement utilisée dans la preuve de l'hérédité.
\item Au Bac, la récurrence sert à établir : des formules de sommes, des inégalités (Bernoulli), des encadrements de suites récurrentes ($0 \leq u_n \leq 2$), des monotonies et des divisibilités.
\item Image mentale : l'\textbf{effet domino} — le premier tombe, chacun entraîne le suivant, donc tous tombent.
\end{itemize}
\end{aretenir}

\begin{exercice}[Pour s'entraîner]
\begin{enumerate}
\item Démontrer par récurrence que pour tout $n \geq 1$ : $1^2 + 2^2 + \cdots + n^2 = \dfrac{n(n+1)(2n+1)}{6}$.
\item Soit $(u_n)$ définie par $u_0 = 2$ et $u_{n+1} = \dfrac{1}{2}u_n + 1$. Démontrer que pour tout $n \in \mathbb{N}$ : $1 \leq u_n \leq 2$.
\item Démontrer que pour tout entier naturel $n$, $7^n - 1$ est divisible par $6$.
\item Retour à Maman Aïcha : on pose $u_0 = 100\,000$ et $u_{n+1} = 1{,}02\,u_n + 5\,000$. Démontrer par récurrence que pour tout $n \in \mathbb{N}$ : $u_n = 100\,000 \times (1{,}02)^n + 250\,000\big((1{,}02)^n - 1\big)$, puis calculer la valeur exacte du capital après 12 mois. Ce capital dépasse-t-il 200 000 FCFA ?
\end{enumerate}
\end{exercice}

\coursec{Suites majorées, minorées, bornées}

Dans la section précédente, tu as appris à démontrer par récurrence des propriétés portant sur les termes d'une suite. Mais à quoi servent concrètement ces démonstrations ? Très souvent, elles servent à établir des \emph{encadrements} : montrer que tous les termes d'une suite restent entre deux valeurs fixes. C'est précisément l'objet de cette section : donner un nom rigoureux à cette idée intuitive de suite qui « ne s'emballe pas », et apprendre à le prouver proprement.

\courssub{Intuition : des suites qui restent sous contrôle}

Imagine le solde d'un compte mobile money MTN d'un commerçant de Douala, relevé chaque soir pendant un mois. Ces trente nombres forment une suite. Deux questions naturelles se posent : le solde peut-il devenir arbitrairement grand (pas de plafond) ? Peut-il descendre en dessous de zéro (pas de plancher) ? Dire qu'il existe un plafond que le solde ne dépasse jamais, c'est dire que la suite est \cle{majorée}. Dire qu'il existe un plancher en dessous duquel elle ne descend jamais, c'est dire qu'elle est \cle{minorée}. Si les deux existent, la suite est \cle{bornée} : tous ses termes vivent dans un intervalle $[m\,;\,M]$.

Toute l'affaire consiste maintenant à transformer cette intuition en énoncés avec quantificateurs, manipulables dans une copie de Bac.

\courssub{Définitions fondamentales}

\begin{definition}[Suite majorée, suite minorée, suite bornée]
Soit $(u_n)_{n\in\mathbb{N}}$ une suite de nombres réels.
\begin{itemize}
\item On dit que $(u_n)$ est \cle{majorée} s'il existe un réel $M$ tel que :
\[ \forall n\in\mathbb{N},\quad u_n \le M. \]
Tout réel $M$ vérifiant cette propriété est appelé un \cle{majorant} de la suite.
\item On dit que $(u_n)$ est \cle{minorée} s'il existe un réel $m$ tel que :
\[ \forall n\in\mathbb{N},\quad m \le u_n. \]
Tout réel $m$ vérifiant cette propriété est appelé un \cle{minorant} de la suite.
\item On dit que $(u_n)$ est \cle{bornée} si elle est à la fois majorée et minorée, c'est-à-dire s'il existe deux réels $m$ et $M$ tels que :
\[ \forall n\in\mathbb{N},\quad m \le u_n \le M. \]
\end{itemize}
\end{definition}

Lisons ces définitions avec attention. Pour prouver qu'une suite est majorée, il suffit d'\emph{exhiber un seul} nombre $M$ qui convient : on n'a pas à trouver « le meilleur » majorant. Si $M$ est un majorant, alors tout réel plus grand que $M$ en est un aussi : une suite majorée possède donc une infinité de majorants. De même, tout réel plus petit qu'un minorant est un minorant.

\begin{exemplebox}[Lecture graphique d'une suite bornée]
Considérons la suite définie pour $n \ge 1$ par $u_n = 1 + \dfrac{\sin n}{n}$. Comme $-1 \le \sin n \le 1$ et $n \ge 1$, on a :
\[ 1 - \frac{1}{n} \le u_n \le 1 + \frac{1}{n}, \quad \text{donc} \quad 0 \le u_n \le 2. \]
La suite $(u_n)$ est bornée : minorée par $0$ et majorée par $2$. Graphiquement, tous ses points restent dans la bande horizontale comprise entre les droites d'équations $y = 0$ et $y = 2$.
\end{exemplebox}

\begin{popfigure}
\begin{center}
\begin{tikzpicture}
\begin{axis}[
  width=0.95\linewidth, height=5.2cm,
  axis lines=middle,
  xlabel={$n$}, ylabel={$u_n$},
  xmin=0, xmax=21, ymin=-0.4, ymax=2.6,
  xtick={1,5,10,15,20},
  ytick={0,1,2},
  tick label style={font=\small},
  label style={font=\small},
  clip=false
]
\addplot[domain=0:21,samples=2,color=popgreen,thick]{2} node[pos=0.97,above,font=\small,color=popgreen]{$y=M=2$};
\addplot[domain=0:21,samples=2,color=popgreen,thick]{0} node[pos=0.97,below,font=\small,color=popgreen]{$y=m=0$};
\addplot[domain=0:21,samples=2,color=popmuted,dashed]{1};
\addplot[only marks,mark=*,mark size=1.6pt,color=popblue,samples at={1,2,...,20}]{1+sin(deg(x))/x};
\end{axis}
\end{tikzpicture}
\legende{Nuage de points de la suite $u_n = 1 + \dfrac{\sin n}{n}$ : tous les termes restent dans la bande $[m\,;\,M] = [0\,;\,2]$. La suite est bornée.}
\end{center}
\end{popfigure}

\begin{popfigure}
\begin{center}
\begin{tikzpicture}[scale=1.05]
\draw[->,thick,popdark] (-0.5,0) -- (10.5,0);
\foreach \x/\lab in {1/{$m$},9/{$M$}}{
  \draw[popgreen,very thick] (\x,-0.18) -- (\x,0.18);
  \node[below,popgreen,font=\small] at (\x,-0.2) {\lab};
}
\foreach \x/\lab in {1.8/{$u_0$},2.9/{$u_1$},4.1/{$u_2$},5.3/{$u_3$},6.4/{$u_4$},7.5/{$u_5$}}{
  \fill[popblue] (\x,0) circle (1.6pt);
  \node[above,popblue,font=\small] at (\x,0.08) {\lab};
}
\foreach \x in {8.1,8.5,8.75}{
  \fill[popblue] (\x,0) circle (1.1pt);
}
\draw[poporange, thick] (1,0.9) -- (9,0.9) node[midway, above=5pt, poporange, font=\small] {tous les termes de la suite vivent ici};
\end{tikzpicture}
\legende{Sur la droite réelle, une suite bornée voit tous ses termes compris entre un minorant $m$ et un majorant $M$.}
\end{center}
\end{popfigure}

\courssub{Caractérisation par la valeur absolue}

Il existe une formulation très commode de la bornitude, qui condense les deux inégalités $m \le u_n \le M$ en une seule.

\begin{propriete}[Caractérisation des suites bornées]
La suite $(u_n)$ est bornée si et seulement s'il existe un réel $A > 0$ tel que :
\[ \forall n\in\mathbb{N},\quad |u_n| \le A. \]
\end{propriete}

Cette équivalence est formatrice à démontrer : elle entraîne proprement les quantificateurs et la valeur absolue. Suivons la démonstration guidée pas à pas.

\begin{experience}[Démonstration guidée de la caractérisation]
\textbf{Sens direct.} Supposons $(u_n)$ bornée : il existe $m$ et $M$ tels que $m \le u_n \le M$ pour tout $n$.
\begin{enumerate}
\item Rappel : pour tout réel $x$, $|x| = \max(x, -x)$, donc $|x| \le \max(|m|, |M|)$ dès que $m \le x \le M$. Justifie ce point : si $x \ge 0$ alors $|x| = x \le M \le |M|$ ; si $x < 0$ alors $|x| = -x \le -m \le |m|$.
\item Posons $A = \max(|m|, |M|)$. Si $A = 0$, alors $m = M = 0$ et la suite est nulle : on peut remplacer $A$ par $1$. Sinon $A > 0$.
\item Conclus : pour tout $n$, $|u_n| \le A$.
\end{enumerate}
\textbf{Sens réciproque.} Supposons qu'il existe $A > 0$ tel que $|u_n| \le A$ pour tout $n$.
\begin{enumerate}
\item Traduis la valeur absolue : $|u_n| \le A$ équivaut à $-A \le u_n \le A$.
\item Pose $m = -A$ et $M = A$.
\item Conclus : $(u_n)$ est minorée par $-A$ et majorée par $A$, donc bornée.
\end{enumerate}
L'équivalence est établie.
\end{experience}

\begin{exemplebox}[Application immédiate]
La suite définie par $u_n = \dfrac{(-1)^n}{n}$ pour $n \ge 1$ vérifie $|u_n| = \dfrac{1}{n} \le 1$ pour tout $n \ge 1$. Elle est donc bornée (on peut choisir $m = -1$ et $M = 1$). De même, la suite $v_n = \sin(n)$ vérifie $|v_n| \le 1$ : elle est bornée, sans que l'on sache quoi que ce soit sur son comportement oscillant.
\end{exemplebox}

\courssub{Exemples et contre-exemples de référence}

Retenons trois situations types, à connaître par c\oe ur car elles reviennent constamment dans les sujets.

\begin{itemize}
\item \textbf{Bornée :} $u_n = \sin(n)$, $u_n = \dfrac{(-1)^n}{n}$, $u_n = 1 + \dfrac{\sin n}{n}$. Toutes leurs valeurs absolues sont majorées par une constante.
\item \textbf{Minorée mais non majorée :} $u_n = n^2$. On a $u_n \ge 0$ pour tout $n$, donc $0$ est un minorant. Mais aucun réel $M$ ne peut majorer la suite : quel que soit $M$, en choisissant $n > \sqrt{M}$ (pour $M > 0$), on obtient $u_n = n^2 > M$. La suite « s'échappe » vers $+\infty$.
\item \textbf{Majorée mais non minorée :} $u_n = -n^2$, majorée par $0$, non minorée (raisonnement symétrique).
\item \textbf{Ni majorée ni minorée :} $u_n = (-1)^n n$, qui prend des valeurs arbitrairement grandes et arbitrairement petites.
\end{itemize}

\begin{attention}[Majorée ne veut pas dire « qui atteint son majorant »]
C'est l'erreur la plus fréquente. Considère la suite définie pour $n \ge 1$ par $u_n = 1 - \dfrac{1}{n}$. Pour tout $n \ge 1$ :
\[ u_n = 1 - \frac{1}{n} < 1. \]
La suite est majorée par $1$, et pourtant \emph{aucun} terme ne vaut $1$ : l'équation $1 - \dfrac{1}{n} = 1$ n'a pas de solution. Un majorant n'est donc pas nécessairement un terme de la suite, et l'inégalité large $u_n \le M$ de la définition autorise les deux situations. Autre piège de rédaction : pour prouver qu'une suite est majorée, il faut une inégalité valable \textbf{pour tout} $n$ ; vérifier $u_0 \le M$, $u_1 \le M$, $u_2 \le M$ sur quelques termes ne prouve rien.
\end{attention}

\courssub{Trois techniques pour établir un encadrement}

\begin{methode}[Prouver qu'une suite est bornée (ou majorée, minorée)]
Selon la façon dont la suite est définie, trois techniques dominent.
\begin{enumerate}
\item \textbf{Calcul direct.} Si $u_n$ est donné par une formule explicite, on majore ou minore directement l'expression en utilisant des inégalités connues : $-1 \le \sin x \le 1$, $n^2 \ge 0$, $\dfrac{1}{n} \le 1$ pour $n \ge 1$, etc. Exemple : $u_n = \dfrac{n}{n+1} = 1 - \dfrac{1}{n+1}$, donc $0 \le u_n \le 1$ pour tout $n$.
\item \textbf{Étude de fonction.} Si $u_n = f(n)$, on étudie les variations et les bornes de la fonction $f$ sur $[0\,;\,+\infty[$ : un tableau de variations donne immédiatement un majorant ou un minorant. Exemple : $f(x) = x e^{-x}$ admet un maximum en $x = 1$ égal à $e^{-1}$, donc la suite $u_n = n e^{-n}$ est majorée par $e^{-1}$ et minorée par $0$.
\item \textbf{Récurrence.} Si la suite est définie par récurrence ($u_{n+1} = f(u_n)$), on ne connaît pas de formule explicite : on prouve l'encadrement $m \le u_n \le M$ par récurrence, en utilisant la stabilité de l'intervalle $[m\,;\,M]$ par $f$, c'est-à-dire le fait que $x \in [m\,;\,M] \Rightarrow f(x) \in [m\,;\,M]$.
\end{enumerate}
\end{methode}

La troisième technique est la plus délicate et la plus rentable au Bac : elle combine le raisonnement par récurrence de la section précédente avec les notions de cette section. Détaillons-la sur un exemple classique.

\begin{exoresolu}[Encadrement d'une suite récurrente par récurrence]
On considère la suite $(u_n)$ définie par $u_0 = 1$ et, pour tout $n \in \mathbb{N}$, $u_{n+1} = \sqrt{u_n + 2}$. Démontrer que pour tout entier naturel $n$ : $0 \le u_n \le 2$.
\tcblower
Notons $P(n)$ la propriété : « $0 \le u_n \le 2$ ».

\textbf{Initialisation.} $u_0 = 1$, et $0 \le 1 \le 2$. Donc $P(0)$ est vraie.

\textbf{Hérédité.} Soit $n \in \mathbb{N}$ tel que $P(n)$ soit vraie : $0 \le u_n \le 2$ (hypothèse de récurrence). Alors :
\begin{align*}
0 \le u_n \le 2 &\implies 2 \le u_n + 2 \le 4 \\
&\implies \sqrt{2} \le \sqrt{u_n + 2} \le \sqrt{4} \quad \text{(la fonction racine carrée est croissante sur } \mathbb{R}_+\text{)} \\
&\implies \sqrt{2} \le u_{n+1} \le 2.
\end{align*}
Or $\sqrt{2} \ge 0$, donc $0 \le u_{n+1} \le 2$ : la propriété $P(n+1)$ est vraie.

\textbf{Conclusion.} Par le principe de récurrence, $P(n)$ est vraie pour tout $n \in \mathbb{N}$ : la suite $(u_n)$ est bornée, minorée par $0$ et majorée par $2$.

\emph{Remarque.} L'étape clé est le passage par la croissance de la fonction $f(x) = \sqrt{x+2}$ : c'est elle qui garantit que l'intervalle $[0\,;\,2]$ est \cle{stable} par $f$, puisque $f([0\,;\,2]) = [\sqrt{2}\,;\,2] \subset [0\,;\,2]$. Cette idée de stabilité sera le moteur de toute la section sur les suites récurrentes.
\end{exoresolu}

\begin{attention}[Conditions d'application à ne pas oublier]
Dans l'hérédité ci-dessus, deux points sont souvent bâclés en copie. D'abord, il faut justifier que $u_{n+1}$ est \emph{bien défini} : la racine carrée exige $u_n + 2 \ge 0$, ce qui découle de l'hypothèse $u_n \ge 0$. Ensuite, l'application de la fonction racine carrée à une inégalité n'est licite que parce que cette fonction est \emph{croissante} : appliquer une fonction décroissante renverserait les inégalités. Cite toujours la monotonie de la fonction utilisée.
\end{attention}

\begin{savaistu}[Pourquoi « borné » est un mot si important ?]
La notion de suite bornée est la porte d'entrée de l'analyse moderne. Le théorème de la section suivante — toute suite croissante majorée converge — montre que la bornitude, combinée à la monotonie, force la convergence. C'est aussi la clé du célèbre théorème de Bolzano–Weierstrass (hors programme, mais culturellement précieux) : de toute suite bornée, on peut extraire une sous-suite convergente. Bernard Bolzano était un mathématicien et philosophe tchèque du début du XIX\textsuperscript{e} siècle ; Karl Weierstrass, professeur de lycée devenu l'un des pères de la rigueur moderne, a fait de ces idées le socle de l'analyse que tu étudies aujourd'hui.
\end{savaistu}

\begin{aretenir}[L'essentiel de la section]
\begin{itemize}
\item $(u_n)$ est \cle{majorée} : $\exists M \in \mathbb{R},\ \forall n \in \mathbb{N},\ u_n \le M$. \cle{Minorée} : $\exists m \in \mathbb{R},\ \forall n \in \mathbb{N},\ m \le u_n$. \cle{Bornée} : les deux à la fois.
\item $(u_n)$ bornée $\iff \exists A > 0,\ \forall n \in \mathbb{N},\ |u_n| \le A$.
\item Un majorant n'est pas forcément atteint : $u_n = 1 - \dfrac{1}{n} < 1$ pour tout $n \ge 1$.
\item Trois techniques d'encadrement : calcul direct sur la formule explicite, étude de la fonction associée, récurrence avec stabilité d'un intervalle pour les suites récurrentes.
\item Pour nier la majoration, il suffit de montrer que tout $M$ est dépassé par au moins un terme (exemple : $u_n = n^2$).
\end{itemize}
\end{aretenir}

\begin{exercice}
Pour chacune des suites suivantes, déterminer si elle est majorée, minorée, bornée (justifier) :
\begin{enumerate}
\item $u_n = \dfrac{3n + 1}{n + 2}$ pour $n \in \mathbb{N}$ ;
\item $v_n = n + (-1)^n$ pour $n \in \mathbb{N}$ ;
\item $w_n = \cos(n) - 2$ pour $n \in \mathbb{N}$.
\end{enumerate}
\end{exercice}

\begin{corrige}[Exercices proposés]
\textbf{1.} On écrit $u_n = \dfrac{3(n+2) - 5}{n+2} = 3 - \dfrac{5}{n+2}$. Pour tout $n \in \mathbb{N}$, $n+2 \ge 2$, donc $0 < \dfrac{5}{n+2} \le \dfrac{5}{2}$. Par conséquent :
\[ 3 - \frac{5}{2} \le u_n < 3 \quad \text{soit} \quad \frac{1}{2} \le u_n < 3. \]
La suite est bornée (minorée par $\frac{1}{2}$, majorée par $3$).

\textbf{2.} Pour tout $n \in \mathbb{N}$, $(-1)^n \ge -1$, donc $v_n = n + (-1)^n \ge n - 1 \ge -1$. La suite est minorée (par $-1$, et même par $0$ car $v_1=0$ est le minimum). En revanche $v_n \ge n - 1$, et $n - 1$ dépasse tout réel $M$ dès que $n > M + 1$ : la suite n'est pas majorée.

\textbf{3.} Comme $-1 \le \cos(n) \le 1$, on a $-3 \le w_n \le -1$ pour tout $n$ : la suite est bornée (et même $|w_n| \le 3$).
\end{corrige}

Tu sais désormais reconnaître et prouver qu'une suite reste entre deux bornes. Dans la section suivante, nous ajouterons une deuxième information — la monotonie — et nous découvrirons que ces deux ingrédients réunis suffisent à garantir la convergence : c'est le puissant théorème de la convergence monotone.

\coursec{Suites monotones et théorème de convergence monotone}

Dans la section précédente, nous avons appris à reconnaître les suites \cle{majorées}, \cle{minorées} et \cle{bornées} : nous savons désormais dire si les termes d'une suite restent « enfermés » dans un intervalle. Mais une information essentielle nous manque encore : \emph{comment la suite évolue-t-elle d'un rang au suivant ?} Monte-t-elle ? Descend-elle ? Oscille-t-elle ? C'est la question de la \cle{monotonie}, et elle va nous offrir le premier grand théorème de convergence de l'année : une suite qui monte sans cesse tout en restant plafonnée ne peut pas faire autre chose que de converger. Ce résultat, d'une intuition lumineuse, est l'un des piliers du programme de Terminale.

\courssub{Suites monotones : définitions}

\textbf{Intuition.}\quad Imagine un élève de Yaoundé qui note chaque soir le solde de son compte Mobile Money. Trois comportements sont possibles : le solde augmente chaque jour (ou reste stable), il diminue chaque jour (ou reste stable), ou bien il monte et descend sans régularité. Dans les deux premiers cas, on dira que la suite des soldes est \emph{monotone} ; dans le troisième, elle ne l'est pas.

\begin{definition}[Suite croissante, suite décroissante, suite monotone]
Soit $(u_n)$ une suite numérique définie pour $n \geq n_0$.
\begin{itemize}
\item $(u_n)$ est \cle{croissante} si, pour tout entier $n \geq n_0$ : $u_{n+1} \geq u_n$.
\item $(u_n)$ est \cle{décroissante} si, pour tout entier $n \geq n_0$ : $u_{n+1} \leq u_n$.
\item $(u_n)$ est \cle{strictement croissante} (resp. \cle{strictement décroissante}) si l'inégalité est stricte : $u_{n+1} > u_n$ (resp. $u_{n+1} < u_n$) pour tout $n \geq n_0$.
\item $(u_n)$ est \cle{monotone} si elle est croissante ou décroissante ; \cle{strictement monotone} si elle est strictement croissante ou strictement décroissante.
\end{itemize}
Une suite \cle{constante} ($u_{n+1} = u_n$ pour tout $n$) est à la fois croissante et décroissante.
\end{definition}

\begin{attention}[Monotonie « à partir d'un certain rang »]
Il arrive qu'une suite ne devienne monotone qu'à partir d'un certain rang $n_0$ (par exemple à partir de $n = 3$). Cela suffit pour appliquer les théorèmes de convergence : ce qui compte, c'est le comportement \emph{à long terme}. En revanche, ne confonds jamais « $u_1 \geq u_0$ et $u_2 \geq u_1$ » avec « $(u_n)$ est croissante » : vérifier quelques rangs ne prouve rien pour \emph{tous} les rangs.
\end{attention}

\courssub{Trois méthodes pour étudier la monotonie}

Selon la forme de l'expression de $u_n$, trois techniques sont à ta disposition. Savoir choisir la bonne est un savoir-faire exigible au Baccalauréat.

\textbf{Méthode 1 : le signe de la différence $u_{n+1} - u_n$.}

C'est la méthode universelle : elle fonctionne toujours, sans condition sur les termes.
\begin{itemize}
\item Si $u_{n+1} - u_n \geq 0$ pour tout $n$, la suite est croissante.
\item Si $u_{n+1} - u_n \leq 0$ pour tout $n$, la suite est décroissante.
\end{itemize}

\textbf{Méthode 2 : la comparaison du quotient $\dfrac{u_{n+1}}{u_n}$ à $1$.}

Valable \emph{uniquement} si tous les termes sont strictement positifs ($u_n > 0$ pour tout $n$) :
\begin{itemize}
\item Si $\dfrac{u_{n+1}}{u_n} \geq 1$ pour tout $n$, la suite est croissante.
\item Si $\dfrac{u_{n+1}}{u_n} \leq 1$ pour tout $n$, la suite est décroissante.
\end{itemize}
Cette méthode est redoutablement efficace pour les suites faisant intervenir des puissances, des produits ou des factorielles (par exemple $u_n = \dfrac{2^n}{n!}$).

\textbf{Méthode 3 : le sens de variation de la fonction $f$ lorsque $u_n = f(n)$.}

Si la suite est définie explicitement par $u_n = f(n)$, où $f$ est une fonction définie sur $[0\,;\,+\infty[$, alors :
\begin{itemize}
\item si $f$ est croissante sur $[0\,;\,+\infty[$, la suite $(u_n)$ est croissante ;
\item si $f$ est décroissante sur $[0\,;\,+\infty[$, la suite $(u_n)$ est décroissante.
\end{itemize}
En effet, $n+1 > n$ et la croissance de $f$ conserve l'ordre : $f(n+1) \geq f(n)$.

\begin{methode}[Choisir la bonne technique de monotonie]
Face à une suite $(u_n)$, procède dans cet ordre :
\begin{enumerate}
\item \textbf{La suite est-elle explicite, $u_n = f(n)$ ?} Si oui, étudie les variations de $f$ sur $[0\,;\,+\infty[$ (dérivée ou fonction de référence). C'est souvent le plus rapide.
\item \textbf{L'expression contient-elle des puissances, produits, quotients de termes positifs ?} Si oui, et si $u_n > 0$ est évident ou prouvé, calcule $\dfrac{u_{n+1}}{u_n}$ et compare à $1$.
\item \textbf{Sinon (suite récurrente, somme, termes de signe quelconque) :} calcule et factorise $u_{n+1} - u_n$, puis étudie son signe.
\item Dans tous les cas, conclus par une phrase complète : « pour tout entier $n$, \dots donc $(u_n)$ est croissante (resp. décroissante) ».
\end{enumerate}
\end{methode}

\begin{attention}[Le piège du quotient avec des termes de signe non constant]
La méthode du quotient exige $u_n > 0$ pour \emph{tout} $n$. Contre-exemple classique : soit $u_n = -2 + n$, c'est-à-dire $-2, -1, 0, 1, 2, \dots$ Cette suite est visiblement croissante. Pourtant $\dfrac{u_1}{u_0} = \dfrac{-1}{-2} = \dfrac{1}{2} < 1$, ce qui suggérerait à tort une décroissance ! Pire : diviser par un terme négatif \emph{renverse} les inégalités. Avant d'écrire le moindre quotient, justifie la positivité des termes (souvent par une récurrence rapide, ou par évidence de l'expression).
\end{attention}

\begin{exemplebox}[Trois suites, trois méthodes]
\begin{itemize}
\item $u_n = n^2 - 3n$ : méthode 1. $u_{n+1} - u_n = (n+1)^2 - 3(n+1) - n^2 + 3n = 2n - 2 \geq 0$ dès que $n \geq 1$. La suite est croissante à partir du rang $1$.
\item $u_n = \dfrac{3^n}{n+1}$ : termes strictement positifs, méthode 2. $\dfrac{u_{n+1}}{u_n} = \dfrac{3^{n+1}}{n+2} \times \dfrac{n+1}{3^n} = \dfrac{3(n+1)}{n+2} \geq 1$ car $3n+3 \geq n+2$. Suite croissante.
\item $u_n = \dfrac{1}{n+1}$ : méthode 3 avec $f(x) = \dfrac{1}{x+1}$, décroissante sur $[0\,;\,+\infty[$. Suite décroissante.
\end{itemize}
\end{exemplebox}

\courssub{Le théorème de convergence monotone}

\textbf{Intuition.}\quad Imagine que tu montes un escalier dont toutes les marches restent sous un plafond horizontal. Tu montes, tu montes encore\ldots{} mais le plafond t'empêche de t'échapper. Que peut-il se passer ? Tu es obligé de te rapprocher de plus en plus d'une « hauteur limite » située sous le plafond. C'est exactement le contenu du théorème fondamental de cette leçon.

\begin{propriete}[Théorème de convergence monotone — admis]
\begin{itemize}
\item Toute suite \cle{croissante} et \cle{majorée} est convergente. Sa limite $\ell$ vérifie $u_n \leq \ell$ pour tout $n$.
\item Toute suite \cle{décroissante} et \cle{minorée} est convergente. Sa limite $\ell$ vérifie $u_n \geq \ell$ pour tout $n$.
\end{itemize}
Ce théorème est \textbf{admis} : sa démonstration repose sur la propriété de la borne supérieure de $\mathbb{R}$ (toute partie non vide et majorée de $\mathbb{R}$ admet une borne supérieure), qui est hors programme. Elle n'est \textbf{pas exigible} au Baccalauréat ; en revanche, savoir \emph{utiliser} le théorème est absolument exigible.
\end{propriete}

\begin{attention}[Ce que le théorème dit — et ne dit pas]
Deux pièges de rédaction reviennent chaque année au Bac :
\begin{itemize}
\item Le théorème donne l'\textbf{existence} de la limite, pas sa \textbf{valeur}. Écrire « $(u_n)$ est croissante et majorée par $5$, donc elle converge vers $5$ » est \textbf{faux} : elle converge vers un réel $\ell \leq 5$. Pour trouver $\ell$, il faut un travail supplémentaire (résolution d'une équation, encadrement, etc.).
\item Les deux hypothèses sont indispensables : une suite croissante \emph{non} majorée ne converge pas (voir ci-dessous), et une suite majorée \emph{non} monotone peut osciller sans converger, comme $u_n = (-1)^n$.
\end{itemize}
\end{attention}

\begin{propriete}[Conséquence : suite croissante non majorée]
Toute suite croissante et \cle{non majorée} diverge vers $+\infty$ :
\[ \lim_{n \to +\infty} u_n = +\infty. \]
De même, toute suite décroissante et non minorée diverge vers $-\infty$.
\end{propriete}

\textbf{Justification (formatrice).}\quad Dire que $(u_n)$ n'est pas majorée signifie : pour tout réel $A$, il existe un rang $N$ tel que $u_N > A$. Comme la suite est croissante, pour tout $n \geq N$ on a $u_n \geq u_N > A$. Ainsi, quel que soit le seuil $A$, tous les termes finissent par le dépasser définitivement : c'est exactement la définition de $\lim u_n = +\infty$.

\begin{aretenir}[Le bilan à retenir]
Pour une suite \cle{monotone}, il n'y a que deux destins possibles :
\begin{center}
\begin{tabularx}{\linewidth}{|l|X|X|}
\hline
\rowcolor{popblueL} \textbf{Comportement} & \textbf{Si majorée (cas croissant)} & \textbf{Si non majorée} \\
\hline
Suite croissante & converge vers $\ell \geq u_n$ & diverge vers $+\infty$ \\
\hline
Suite décroissante & converge vers $\ell \leq u_n$ si minorée & diverge vers $-\infty$ si non minorée \\
\hline
\end{tabularx}
\end{center}
Autrement dit : \emph{une suite monotone a toujours une limite, finie ou infinie.}
\end{aretenir}

\begin{popfigure}
\begin{center}
\begin{tikzpicture}
\begin{axis}[
width=0.85\linewidth, height=6cm,
xmin=0, xmax=12.5, ymin=0, ymax=1.6,
xlabel={$n$}, ylabel={$u_n$},
axis lines=left, grid=both, grid style={popline!40},
legend style={at={(0.97,0.35)},anchor=east, font=\small}
]
\addplot[popcurve, domain=0:12.5, samples=100] {1};
\addlegendentry{asymptote $y = \ell = 1$}
\addplot[only marks, mark=*, mark size=2pt, poporange] coordinates {
(1,0) (2,0.5) (3,0.6667) (4,0.75) (5,0.8) (6,0.8333) (7,0.8571) (8,0.875) (9,0.8889) (10,0.9) (11,0.9091) (12,0.9167)
};
\addlegendentry{$u_n = 1 - \dfrac{1}{n}$}
\end{axis}
\end{tikzpicture}
\end{center}
\legende{La suite $u_n = 1 - \dfrac{1}{n}$ est croissante et majorée par $1$ : ses points montent en se rapprochant de l'asymptote horizontale $y = 1$ sans jamais la dépasser.}
\end{popfigure}

\begin{popfigure}
\begin{center}
\begin{tikzpicture}
\begin{axis}[
width=0.85\linewidth, height=6cm,
xmin=0, xmax=9, ymin=0, ymax=20,
xlabel={$n$}, ylabel={$u_n$},
axis lines=left, grid=both, grid style={popline!40},
legend style={at={(0.03,0.97)},anchor=north west, font=\small}
]
\addplot[only marks, mark=*, mark size=2pt, popgreen] coordinates {
(0,0) (1,1.41) (2,2) (3,2.45) (4,2.83) (5,3.16) (6,3.46) (7,3.74) (8,4)
};
\addlegendentry{$u_n = \sqrt{n}$ : croissante non majorée, $\lim u_n = +\infty$}
\addplot[only marks, mark=square*, mark size=2pt, poppurple] coordinates {
(0,1) (1,0.5) (2,0.3333) (3,0.25) (4,0.2) (5,0.1667) (6,0.1429) (7,0.125) (8,0.1111)
};
\addlegendentry{$v_n = \dfrac{1}{n+1}$ : décroissante minorée par $0$, converge vers $0$}
\end{axis}
\end{tikzpicture}
\end{center}
\legende{Deux destins monotones : la suite croissante non majorée $(\sqrt{n})$ « s'échappe » vers $+\infty$, tandis que la suite décroissante minorée $\left(\dfrac{1}{n+1}\right)$ converge vers $0$.}
\end{popfigure}

\courssub{Exemples d'application}

\begin{exoresolu}[Étude complète de la suite $u_n = \dfrac{n}{n+1}$]
On considère la suite définie pour tout $n \in \mathbb{N}$ par $u_n = \dfrac{n}{n+1}$.
\begin{enumerate}
\item Étudier la monotonie de $(u_n)$.
\item Montrer que $(u_n)$ est bornée.
\item En déduire que $(u_n)$ converge et déterminer sa limite.
\end{enumerate}
\tcblower
\textbf{Solution détaillée.}

\textbf{1. Monotonie.} Remarquons que $u_n = \dfrac{n+1-1}{n+1} = 1 - \dfrac{1}{n+1}$. Calculons la différence (méthode 1) :
\begin{align*}
u_{n+1} - u_n &= \left(1 - \dfrac{1}{n+2}\right) - \left(1 - \dfrac{1}{n+1}\right) = \dfrac{1}{n+1} - \dfrac{1}{n+2} \\
&= \dfrac{(n+2)-(n+1)}{(n+1)(n+2)} = \dfrac{1}{(n+1)(n+2)}.
\end{align*}
Pour tout $n \in \mathbb{N}$, $(n+1)(n+2) > 0$, donc $u_{n+1} - u_n > 0$ : la suite $(u_n)$ est \textbf{strictement croissante}.

\emph{Autre rédaction possible (méthode 3) :} $u_n = f(n)$ avec $f(x) = 1 - \dfrac{1}{x+1}$, et $f'(x) = \dfrac{1}{(x+1)^2} > 0$ sur $[0\,;\,+\infty[$, donc $f$ est croissante et $(u_n)$ aussi.

\textbf{2. Bornitude.} Pour tout $n \in \mathbb{N}$, $\dfrac{1}{n+1} > 0$, donc $u_n = 1 - \dfrac{1}{n+1} < 1$ : la suite est \textbf{majorée par $1$}. De plus $u_n \geq u_0 = 0$ par croissance : la suite est minorée par $0$. Donc $0 \leq u_n < 1$ pour tout $n$ : $(u_n)$ est bornée.

\textbf{3. Convergence.} La suite $(u_n)$ est croissante et majorée par $1$ : d'après le théorème de convergence monotone, elle \textbf{converge} vers un réel $\ell \leq 1$. Pour déterminer $\ell$, on calcule directement :
\[ \lim_{n \to +\infty} u_n = \lim_{n \to +\infty} \left(1 - \dfrac{1}{n+1}\right) = 1 - 0 = 1. \]
\textbf{Conclusion :} $(u_n)$ converge vers $\ell = 1$. Vérifions numériquement : $u_9 = 0{,}9$ ; $u_{99} = 0{,}99$ ; $u_{999} = 0{,}999$. La suite s'approche bien de $1$ par valeurs inférieures, conformément au théorème.
\end{exoresolu}

\begin{exemplebox}[Suite géométrique de raison $0 < q < 1$]
Soit $(u_n)$ la suite géométrique de premier terme $u_0 > 0$ et de raison $q$ avec $0 < q < 1$ : $u_n = u_0 \times q^n$.
\begin{itemize}
\item \textbf{Monotonie (méthode 2) :} tous les termes sont strictement positifs et $\dfrac{u_{n+1}}{u_n} = q < 1$, donc $(u_n)$ est strictement décroissante.
\item \textbf{Bornitude :} $u_n > 0$ pour tout $n$, donc $(u_n)$ est minorée par $0$.
\item \textbf{Convergence :} décroissante et minorée, $(u_n)$ converge d'après le théorème de convergence monotone. On sait par ailleurs (limites des suites géométriques) que sa limite est $0$.
\end{itemize}
\emph{Illustration concrète :} un forfait data Orange dont la consommation quotidienne est multipliée par $0{,}8$ chaque jour (habitude qui s'essouffle) : $u_0 = 500$ Mo, $u_1 = 400$, $u_2 = 320$, $u_3 = 256$\ldots{} La consommation décroît et tend vers $0$.
\end{exemplebox}

\begin{exoresolu}[Une suite croissante non majorée]
Soit $(u_n)$ définie par $u_n = \sqrt{n}$ pour $n \in \mathbb{N}$. Montrer que $(u_n)$ est croissante, qu'elle n'est pas majorée, et en déduire sa limite.
\tcblower
\textbf{Solution.} \textbf{Monotonie :} la fonction $f(x) = \sqrt{x}$ est strictement croissante sur $[0\,;\,+\infty[$, donc $(u_n)$ est strictement croissante (méthode 3).

\textbf{Non majorée :} soit $A$ un réel quelconque. Pour tout entier $n > A^2$ (un tel entier existe toujours), on a $\sqrt{n} > A$. Aucun réel $A$ ne majore la suite : $(u_n)$ n'est pas majorée.

\textbf{Limite :} $(u_n)$ est croissante et non majorée, donc d'après la conséquence du théorème de convergence monotone, $\lim\limits_{n \to +\infty} u_n = +\infty$.
\end{exoresolu}

\begin{savaistu}[Pourquoi $\mathbb{R}$ et pas $\mathbb{Q}$ ?]
Le théorème de convergence monotone est l'une des propriétés qui distinguent profondément $\mathbb{R}$ de $\mathbb{Q}$. Considère la suite des approximations décimales de $\sqrt{2}$ :
\[ u_0 = 1, \quad u_1 = 1{,}4, \quad u_2 = 1{,}41, \quad u_3 = 1{,}414, \quad u_4 = 1{,}4142, \dots \]
Tous ses termes sont des nombres \emph{rationnels} (ce sont des décimaux). Elle est croissante et majorée par $2$ (même par $1{,}5$). Pourtant, sa limite $\sqrt{2}$ \textbf{n'appartient pas à} $\mathbb{Q}$, comme le savait déjà l'école de Pythagore, qui en fut profondément troublée ! Ainsi, \emph{dans} $\mathbb{Q}$, une suite croissante majorée peut ne pas converger. C'est la propriété de la borne supérieure — l'« absence de trous » de la droite réelle — qui garantit que dans $\mathbb{R}$, la limite existe toujours. Ce théorème, que tu utilises désormais à chaque étude de suite, est donc une fenêtre sur la construction même des nombres réels, achevée seulement au XIX\textsuperscript{e} siècle par Dedekind et Cantor.
\end{savaistu}

\begin{exercice}[Pour t'entraîner]
\begin{enumerate}
\item Soit $(u_n)$ définie par $u_n = \dfrac{2^n}{n!}$ pour $n \geq 1$. Montrer que $(u_n)$ est décroissante à partir d'un rang à préciser, puis qu'elle converge.
\item Soit $(v_n)$ définie par $v_n = 3 - \dfrac{2}{n+1}$. Étudier sa monotonie, sa bornitude, et conclure sur sa convergence en précisant sa limite.
\item Soit $(w_n)$ définie par $w_0 = 1$ et $w_{n+1} = w_n + 2n + 1$. Calculer $w_{n+1} - w_n$, en déduire la monotonie de $(w_n)$, puis montrer que $(w_n)$ n'est pas majorée (on pourra remarquer que $w_n = (n+1)^2$). Conclure sur sa limite.
\end{enumerate}
\end{exercice}

\begin{corrige}[Exercice n°1]
Pour $n \geq 1$, $u_n > 0$ et $\dfrac{u_{n+1}}{u_n} = \dfrac{2^{n+1}}{(n+1)!} \times \dfrac{n!}{2^n} = \dfrac{2}{n+1}$. Pour $n \geq 1$, $\dfrac{2}{n+1} \leq 1$, donc $(u_n)$ est décroissante (à partir du rang $1$). De plus $u_n > 0$ : la suite est minorée. Décroissante et minorée, elle converge d'après le théorème de convergence monotone.
\end{corrige}

\begin{corrige}[Exercice n°2]
$v_{n+1} - v_n = \dfrac{2}{n+1} - \dfrac{2}{n+2} = \dfrac{2}{(n+1)(n+2)} > 0$ : $(v_n)$ est strictement croissante. Comme $\dfrac{2}{n+1} > 0$, on a $v_n < 3$ : majorée par $3$ ; et $v_n \geq v_0 = 1$ : minorée. Croissante et majorée, $(v_n)$ converge, et $\lim v_n = 3 - 0 = 3$.
\end{corrige}

\begin{corrige}[Exercice n°3]
$w_{n+1} - w_n = 2n + 1 > 0$ pour tout $n \geq 0$ : $(w_n)$ est strictement croissante. En calculant : $w_0 = 1 = 1^2$, $w_1 = 1 + 1 = 4 = 2^2$, $w_2 = 4 + 3 = 9 = 3^2$, ce qui suggère $w_n = (n+1)^2$ (se prouve rigoureusement par récurrence, voir section 1). Alors $w_n = (n+1)^2$ n'est pas majorée : pour tout réel $A$, dès que $n > \sqrt{A} - 1$, on a $w_n > A$. Croissante et non majorée, $(w_n)$ diverge vers $+\infty$.
\end{corrige}

En résumé : la monotonie est le comportement, la bornitude est la contrainte, et le théorème de convergence monotone est le pont qui relie les deux à la conclusion « la suite converge ». Dans la section suivante, nous appliquerons cette stratégie complète aux suites définies par une relation de récurrence $u_{n+1} = f(u_n)$, le type de suite le plus fréquent au Baccalauréat.

\coursec{Suites récurrentes $u_{n+1} = f(u_n)$ : étude de la convergence}

Dans la section précédente, nous avons établi le \cle{théorème de convergence monotone} : toute suite croissante et majorée converge, toute suite décroissante et minorée converge. Ce théorème est une arme redoutable... à condition de savoir prouver la monotonie et la bornitude ! Or, au Baccalauréat, la grande majorité des suites proposées ne sont pas définies par une formule explicite $u_n = g(n)$, mais par une \cle{relation de récurrence} du type $u_{n+1} = f(u_n)$ : chaque terme se calcule à partir du précédent, comme une machine qui transforme un nombre en un autre. Cette section te donne la méthode complète, étape par étape, pour étudier ces suites — c'est l'un des exercices les plus fréquents du Bac C, D, E et TI.

\courssub{Définition et rôle du premier terme}

\begin{definition}[Suite récurrente]
Soit $f$ une fonction définie sur un ensemble $D$ de $\mathbb{R}$ et $u_0 \in D$. La \cle{suite récurrente} de premier terme $u_0$ associée à $f$ est la suite $(u_n)$ définie par :
\[
\begin{cases}
u_0 \text{ donné} \\
u_{n+1} = f(u_n) \quad \text{pour tout } n \in \mathbb{N}
\end{cases}
\]
à condition que tous les termes restent dans le domaine de définition de $f$.
\end{definition}

L'idée est celle d'une \emph{chaîne} : $u_1 = f(u_0)$, puis $u_2 = f(u_1)$, puis $u_3 = f(u_2)$, etc. Chaque terme est l'image du précédent par $f$.

\begin{attention}[Le premier terme change tout]
Contrairement aux suites explicites, le comportement d'une suite récurrente dépend \textbf{crucialement} de $u_0$. Pour $u_{n+1} = u_n^2$ : si $u_0 = \num{0,5}$, la suite tend vers $0$ ; si $u_0 = 2$, elle explose vers $+\infty$. Ne jamais oublier de tenir compte de $u_0$ dans l'analyse.
\end{attention}

\begin{exemplebox}[Premiers termes]
Soit $(u_n)$ définie par $u_0 = 1$ et $u_{n+1} = \sqrt{u_n + 2}$. Alors :
\[
u_1 = \sqrt{1+2} = \sqrt{3} \approx \num{1,732} ; \quad u_2 = \sqrt{\sqrt{3}+2} \approx \num{1,931} ; \quad u_3 \approx \num{1,983}.
\]
On devine une convergence vers $2$. Toute cette section va consister à le \emph{prouver} rigoureusement.
\end{exemplebox}

\courssub{Représentation graphique : la construction en escalier}

Pour visualiser une suite récurrente, on trace dans un repère la \cle{courbe} $\mathcal{C}_f$ d'équation $y = f(x)$ et la \cle{première bissectrice} $\Delta$ d'équation $y = x$. Le procédé de construction est le suivant :

\begin{enumerate}
\item On place $u_0$ sur l'axe des abscisses ; on monte verticalement jusqu'à $\mathcal{C}_f$ : l'ordonnée du point obtenu est $u_1 = f(u_0)$.
\item Pour reporter $u_1$ sur l'axe des abscisses, on se déplace \emph{horizontalement} jusqu'à la droite $\Delta$ (où ordonnée = abscisse).
\item On recommence : verticale vers $\mathcal{C}_f$ (on obtient $u_2$), horizontale vers $\Delta$, etc.
\end{enumerate}

On obtient une ligne brisée en \emph{escalier} (quand $f$ est croissante) ou en \emph{escargot} (quand $f$ est décroissante), qui s'enroule autour du point d'intersection de $\mathcal{C}_f$ et de $\Delta$.

\begin{popfigure}
\begin{center}
\begin{tikzpicture}[scale=2.2]
\draw[->,popmuted] (-0.15,0) -- (2.7,0) node[below left] {$x$};
\draw[->,popmuted] (0,-0.15) -- (0,2.7) node[below left] {$y$};
\draw[popmuted,dashed] (0,0) -- (2.55,2.55) node[pos=0.9,sloped,above] {\small $y=x$};
\draw[popblue,thick,domain=0:2.55,samples=120] plot (\x,{sqrt(max(0,\x+2))}) node[pos=0.35,sloped,above] {\small $y=\sqrt{x+2}$};
% escalier : u0=1
\draw[poporange,thick] (1,0) node[below] {\small $u_0$} -- (1,{sqrt(max(0,3))}) -- ({sqrt(max(0,3))},{sqrt(max(0,3))});
\draw[poporange,thick] ({sqrt(max(0,3))},{sqrt(max(0,3))}) -- ({sqrt(max(0,3))},{sqrt(max(0,sqrt(3)+2))}) -- ({sqrt(max(0,sqrt(3)+2))},{sqrt(max(0,sqrt(3)+2))});
\draw[poporange,thick] ({sqrt(max(0,sqrt(3)+2))},{sqrt(max(0,sqrt(3)+2))}) -- ({sqrt(max(0,sqrt(3)+2))},{sqrt(max(0,sqrt(sqrt(3)+2)+2))}) -- ({sqrt(max(0,sqrt(sqrt(3)+2)+2))},{sqrt(max(0,sqrt(sqrt(3)+2)+2))});
\fill[popink] (2,2) circle (0.7pt) node[above right] {\small $(2\,;\,2)$};
\draw[popmuted,dotted] (2,0) node[below] {\small $2$} -- (2,2);
\end{tikzpicture}
\end{center}
\legende{Construction en escalier de $u_{n+1} = \sqrt{u_n+2}$ avec $u_0 = 1$ : les marches montent vers le point fixe $(2\,;\,2)$.}
\end{popfigure}

\begin{popfigure}
\begin{center}
\begin{tikzpicture}[scale=1.5]
\clip (0,0) rectangle (4.4,4.4);
\draw[->,popmuted] (-0.15,0) -- (4.4,0) node[below left] {$x$};
\draw[->,popmuted] (0,-0.15) -- (0,4.4) node[below left] {$y$};
\draw[popmuted,dashed] (0,0) -- (4.2,4.2) node[pos=0.92,sloped,above] {\small $y=x$};
\draw[poppurple,thick,domain=1.17:4.2,samples=120] plot (\x,{1+4/\x}) node[pos=0.85,sloped,above] {\small $y=1+\dfrac{4}{x}$};
% escargot : u0 = 1.2
\draw[popgreen,thick] (1.2,0) node[below] {\small $u_0$} -- (1.2,4.3333) -- (4.3333,4.3333);
\draw[popgreen,thick] (4.3333,4.3333) -- (4.3333,1.9231) -- (1.9231,1.9231);
\draw[popgreen,thick] (1.9231,1.9231) -- (1.9231,3.0769) -- (3.0769,3.0769);
\draw[popgreen,thick] (3.0769,3.0769) -- (3.0769,2.2987) -- (2.2987,2.2987);
\fill[popink] (2.5616,2.5616) circle (0.9pt) node[below right] {\small point fixe $\ell$};
\end{tikzpicture}
\end{center}
\legende{Construction en escargot pour $f$ décroissante ($f(x) = 1 + \frac{4}{x}$) : la ligne brisée s'enroule autour du point fixe, les termes alternent de part et d'autre de $\ell$.}
\end{popfigure}

\courssub{Points fixes : la condition nécessaire $\ell = f(\ell)$}

Observe les figures : quand la suite converge, elle converge vers un point d'\emph{intersection} de la courbe de $f$ avec la droite $y = x$, c'est-à-dire vers une solution de l'équation $f(x) = x$. Ce n'est pas un hasard.

\begin{definition}[Point fixe]
Soit $f$ une fonction. On appelle \cle{point fixe} de $f$ tout réel $\ell$ tel que $f(\ell) = \ell$.
\end{definition}

\begin{propriete}[La limite est un point fixe]
Soit $(u_n)$ une suite définie par $u_{n+1} = f(u_n)$. Si $(u_n)$ converge vers un réel $\ell$ et si $f$ est \cle{continue} en $\ell$, alors :
\[
f(\ell) = \ell.
\]
\emph{Démonstration exigible au Bac.}
\end{propriete}

\begin{experience}[Démonstration par passage à la limite]
\textbf{Hypothèses :} $u_n \to \ell$ et $f$ continue en $\ell$.
\begin{enumerate}
\item Puisque $u_n \to \ell$, la suite $(u_{n+1})$, qui est la même suite décalée d'un rang, converge aussi vers $\ell$ : $\lim\limits_{n\to+\infty} u_{n+1} = \ell$.
\item Comme $u_n \to \ell$ et que $f$ est continue en $\ell$, le théorème de composition des limites donne : $\lim\limits_{n\to+\infty} f(u_n) = f(\ell)$.
\item Or la relation de récurrence affirme que $u_{n+1} = f(u_n)$ pour tout $n$ : les suites $(u_{n+1})$ et $(f(u_n))$ sont \emph{égales}, donc ont la même limite.
\item Par unicité de la limite : $\ell = f(\ell)$. \hfill $\blacksquare$
\end{enumerate}
\end{experience}

Ce résultat fournit les \emph{candidats} possibles pour la limite : il suffit de résoudre l'équation $f(x) = x$.

\begin{exemplebox}[Recherche des points fixes]
Pour $f(x) = \sqrt{x+2}$, résolvons $f(x) = x$ sur $[0\,;\,+\infty[$ :
\[
\sqrt{x+2} = x \iff x+2 = x^2 \ (\text{avec } x \geq 0) \iff x^2 - x - 2 = 0 \iff (x-2)(x+1) = 0.
\]
Seule la solution $x = 2$ est positive. Donc \textbf{si} la suite converge, sa limite ne peut être que $\ell = 2$.
\end{exemplebox}

\begin{attention}[Le piège classique du Bac]
L'équation $\ell = f(\ell)$ ne prouve \textbf{pas} que la suite converge ! Elle dit seulement : \emph{si} la suite converge, \emph{alors} sa limite est parmi les solutions de $f(x) = x$. Rédiger « $u_{n+1} = f(u_n)$ donc $\ell = f(\ell)$, d'où $\ell = 2$ » \textbf{sans avoir d'abord établi la convergence} est une erreur grave, sanctionnée au Bac. La convergence doit être prouvée \emph{avant} (typiquement par le théorème de convergence monotone), et la continuité de $f$ en $\ell$ doit être mentionnée.
\end{attention}

\courssub{Intervalle stable et bornitude}

Pour prouver que tous les termes de la suite restent dans un intervalle $I$ (par exemple pour montrer que la suite est majorée), l'outil clé est la notion d'\emph{intervalle stable}.

\begin{definition}[Intervalle stable]
Un intervalle $I$ est dit \cle{stable} par $f$ (ou $f$-stable) si $f(I) \subset I$, c'est-à-dire si pour tout $x \in I$, on a $f(x) \in I$.
\end{definition}

\begin{propriete}[Bornitude par récurrence]
Si $I$ est stable par $f$ et si $u_0 \in I$, alors $u_n \in I$ pour tout $n \in \mathbb{N}$. En particulier, si $I = [a\,;\,b]$, la suite $(u_n)$ est bornée (minorée par $a$, majorée par $b$).
\end{propriete}

La preuve est une récurrence immédiate : l'initialisation est l'hypothèse $u_0 \in I$, et l'hérédité vient de la stabilité : si $u_n \in I$, alors $u_{n+1} = f(u_n) \in f(I) \subset I$.

\begin{methode}[Prouver qu'un intervalle est stable]
Pour montrer que $I = [a\,;\,b]$ est stable par $f$ :
\begin{enumerate}
\item Étudier les variations de $f$ sur $I$ (tableau de variations).
\item Calculer les images des bornes :
      \begin{itemize}
      \item si $f$ est croissante, $f(I) = [f(a)\,;\,f(b)]$ ;
      \item si $f$ est décroissante, $f(I) = [f(b)\,;\,f(a)]$.
      \end{itemize}
\item Vérifier que $f(I) \subset [a\,;\,b]$, c'est-à-dire $f(a) \geq a$ et $f(b) \leq b$ (cas croissante) ou $f(b) \geq a$ et $f(a) \leq b$ (cas décroissante).
\end{enumerate}
\end{methode}

\courssub{Monotonie : le rôle décisif de la croissance de $f$}

\begin{propriete}[Si $f$ est croissante, la suite est monotone]
Soit $I$ un intervalle stable par $f$, avec $u_0 \in I$. Si $f$ est \cle{croissante} sur $I$, alors la suite $(u_n)$ est \cle{monotone}. Plus précisément :
\begin{itemize}
\item si $u_1 \geq u_0$, alors $(u_n)$ est croissante ;
\item si $u_1 \leq u_0$, alors $(u_n)$ est décroissante.
\end{itemize}
\emph{Démonstration guidée (formatrice, souvent demandée au Bac).}
\end{propriete}

\begin{experience}[Démonstration guidée par récurrence]
Supposons $f$ croissante sur $I$ stable, $u_0 \in I$ et $u_1 \geq u_0$. Montrons par récurrence que pour tout $n$ : $u_{n+1} \geq u_n$.
\begin{enumerate}
\item \textbf{Initialisation} ($n = 0$) : $u_1 \geq u_0$ par hypothèse. La propriété est vraie au rang $0$.
\item \textbf{Hérédité} : supposons $u_{n+1} \geq u_n$ pour un rang $n$ fixé. Les deux termes sont dans $I$ (stabilité), et $f$ est croissante sur $I$ : elle conserve l'ordre. Donc :
\[
u_{n+1} \geq u_n \implies f(u_{n+1}) \geq f(u_n) \implies u_{n+2} \geq u_{n+1}.
\]
La propriété est héréditaire.
\item \textbf{Conclusion} : par le principe de récurrence, $u_{n+1} \geq u_n$ pour tout $n$ : la suite est croissante. \hfill $\blacksquare$
\end{enumerate}
Le cas $u_1 \leq u_0$ se traite de même, la croissance de $f$ conservant l'inégalité dans l'autre sens.
\end{experience}

\begin{aretenir}[Idée directrice]
Une fonction croissante \emph{conserve les inégalités}. Donc l'inégalité initiale entre $u_0$ et $u_1$ se « propage » de proche en proche à tous les rangs : c'est exactement la monotonie. Le sens de variation de la suite se lit donc simplement sur le \textbf{signe de $u_1 - u_0$}.
\end{aretenir}

\courssub{La stratégie complète en cinq étapes}

\begin{methode}[Étude complète d'une suite $u_{n+1} = f(u_n)$]
\begin{enumerate}
\item \textbf{Intervalle stable} : trouver un intervalle $I$ (souvent donné ou suggéré par l'énoncé) tel que $f(I) \subset I$ et $u_0 \in I$. En déduire par récurrence que $u_n \in I$ pour tout $n$ : la suite est \emph{bornée}.
\item \textbf{Monotonie} : vérifier que $f$ est croissante sur $I$ (étude du signe de $f'$), puis comparer $u_1$ et $u_0$ pour conclure au sens de variation.
\item \textbf{Convergence} : appliquer le théorème de convergence monotone (croissante + majorée, ou décroissante + minorée) : la suite converge vers un réel $\ell \in I$.
\item \textbf{Points fixes} : résoudre l'équation $f(x) = x$ dans $I$.
\item \textbf{Conclusion} : $f$ étant continue (le justifier !), la limite $\ell$ est l'unique point fixe de $f$ dans $I$ compatible avec la position de la suite.
\end{enumerate}
\end{methode}

\begin{exoresolu}[L'exemple de référence : $u_{n+1} = \sqrt{u_n + 2}$, $u_0 = 1$]
Soit $(u_n)$ la suite définie par $u_0 = 1$ et $u_{n+1} = \sqrt{u_n + 2}$ pour tout $n \in \mathbb{N}$.
\begin{enumerate}
\item Montrer que pour tout $n$, $0 \leq u_n \leq 2$.
\item Étudier la monotonie de $(u_n)$.
\item En déduire que $(u_n)$ converge et déterminer sa limite.
\end{enumerate}
\tcblower
\textbf{1. Bornitude.} Posons $f(x) = \sqrt{x+2}$, définie sur $[-2\,;\,+\infty[$. Montrons que $I = [0\,;\,2]$ est stable par $f$. La fonction $f$ est croissante sur $I$ (composée de $x \mapsto x+2$ croissante et de la racine carrée croissante), donc :
\[
f(I) = [f(0)\,;\,f(2)] = [\sqrt{2}\,;\,2] \subset [0\,;\,2].
\]
Prouvons par récurrence que $u_n \in [0\,;\,2]$ pour tout $n$.
\emph{Initialisation} : $u_0 = 1 \in [0\,;\,2]$. \emph{Hérédité} : si $u_n \in [0\,;\,2]$, alors $u_{n+1} = f(u_n) \in f(I) \subset [0\,;\,2]$. \emph{Conclusion} : pour tout $n$, $0 \leq u_n \leq 2$ : la suite est bornée.

\textbf{2. Monotonie.} $f$ est croissante sur $[0\,;\,2]$ et $u_1 = \sqrt{3} \approx \num{1,732} > 1 = u_0$. D'après la propriété du cours (démontrée par récurrence), $(u_n)$ est \textbf{croissante}.

\emph{Rédaction directe possible} : montrons par récurrence que $u_{n+1} \geq u_n$. Initialisation : $u_1 = \sqrt{3} \geq u_0$. Hérédité : si $u_{n+1} \geq u_n$, alors $u_{n+1} + 2 \geq u_n + 2$, et la racine carrée étant croissante, $\sqrt{u_{n+1}+2} \geq \sqrt{u_n+2}$, soit $u_{n+2} \geq u_{n+1}$.

\textbf{3. Convergence et limite.} La suite $(u_n)$ est croissante et majorée par $2$ : par le \cle{théorème de convergence monotone}, elle converge vers un réel $\ell \in [0\,;\,2]$.

La fonction $f$ est continue sur $[0\,;\,2]$, donc en $\ell$. Par passage à la limite dans $u_{n+1} = f(u_n)$, on obtient $\ell = \sqrt{\ell + 2}$, c'est-à-dire $\ell^2 - \ell - 2 = 0$ avec $\ell \geq 0$, d'où $\ell = 2$ ou $\ell = -1$. Comme $\ell \in [0\,;\,2]$, on conclut :
\[
\lim_{n \to +\infty} u_n = 2.
\]
\end{exoresolu}

\courssub{Complément : le cas où $f$ est décroissante}

\begin{propriete}[Cas décroissant : sous-suites de rangs pairs et impairs]
Soit $I$ un intervalle stable par $f$, $u_0 \in I$, et supposons $f$ \cle{décroissante} sur $I$. Alors la suite $(u_n)$ n'est en général pas monotone : elle \emph{alterne} autour du point fixe. En revanche, la fonction $f \circ f$ est croissante, et les deux sous-suites extraites
\[
(u_{2n}) \quad \text{et} \quad (u_{2n+1})
\]
sont monotones de \textbf{sens contraires} : l'une est croissante, l'autre décroissante. Si de plus elles convergent vers la même limite $\ell$, alors $(u_n)$ converge vers $\ell$.
\end{propriete}

C'est exactement ce que montre la figure en « escargot » : les termes de rang pair grimpent vers $\ell$ par la gauche pendant que les termes de rang impair descendent vers $\ell$ par la droite.

\begin{attention}[Ne pas conclure trop vite]
Si $f$ est décroissante, il est \textbf{faux} d'affirmer que $(u_n)$ est décroissante ! Le réflexe correct est d'étudier séparément $(u_{2n})$ et $(u_{2n+1})$ via la fonction croissante $g = f \circ f$, puis de vérifier que les deux limites coïncident. Ce complément est régulièrement exploité dans les problèmes de Bac les plus difficiles.
\end{attention}

\begin{exercice}[Pour s'entraîner]
Soit $(v_n)$ définie par $v_0 = 3$ et $v_{n+1} = \dfrac{4 v_n - 3}{v_n}$ pour tout $n$.
\begin{enumerate}
\item Montrer que $[1\,;\,3]$ est stable par $f(x) = \dfrac{4x-3}{x}$.
\item Étudier la monotonie de $(v_n)$ et en déduire sa convergence.
\item Déterminer sa limite.
\end{enumerate}
\end{exercice}

\begin{corrige}[Exercice]
\textbf{1.} On a $f(x) = \dfrac{4x-3}{x} = 4 - \dfrac{3}{x}$, définie et dérivable sur $[1\,;\,3]$. Sa dérivée est $f'(x) = \dfrac{3}{x^2} > 0$ : $f$ est strictement croissante sur $[1\,;\,3]$. Donc $f([1\,;\,3]) = [f(1)\,;\,f(3)] = [1\,;\,3] \subset [1\,;\,3]$ : l'intervalle est stable. Comme $v_0 = 3 \in [1\,;\,3]$, une récurrence immédiate donne $v_n \in [1\,;\,3]$ pour tout $n$.

\textbf{2.} Calculons $v_1 = f(3) = 4 - 1 = 3 = v_0$. Comme $f$ est croissante sur l'intervalle stable $[1\,;\,3]$ et que $v_1 = v_0$, une récurrence simple montre que $v_{n+1} = v_n$ pour tout $n$ : la suite $(v_n)$ est \textbf{constante} égale à $3$. Elle est donc convergente.

\textbf{3.} La fonction $f$ est continue sur $[1\,;\,3]$, donc la limite $\ell$ vérifie $\ell = 4 - \dfrac{3}{\ell}$, soit $\ell^2 - 4\ell + 3 = 0$, d'où $\ell = 1$ ou $\ell = 3$. La suite étant constante égale à $3$, sa limite est $\ell = 3$.
\end{corrige}

\begin{savaistu}[La méthode de Héron d'Alexandrie]
Au I\textsuperscript{er} siècle de notre ère, le mathématicien grec \cle{Héron d'Alexandrie} décrivait déjà un procédé pour calculer la racine carrée d'un nombre $a > 0$ : partir d'une estimation $u_0$, puis itérer
\[
u_{n+1} = \frac{1}{2}\left(u_n + \frac{a}{u_n}\right).
\]
C'est une suite récurrente avec $f(x) = \frac{1}{2}\left(x + \frac{a}{x}\right)$ ! On vérifie que $f$ laisse stable $[\sqrt{a}\,;\,+\infty[$, que la suite est décroissante à partir du rang 1, minorée par $\sqrt{a}$, et que l'équation $f(x) = x$ donne $x^2 = a$. La convergence est si rapide que quelques itérations suffisent : pour $a = 2$ et $u_0 = 1$, on obtient $u_3 \approx \num{1,4142135}$, déjà exact à six décimales. Cette méthode, vieille de deux mille ans, est encore au cœur des algorithmes de nos calculatrices — et c'est un cas particulier de la \emph{méthode de Newton} que tu rencontreras peut-être dans tes études supérieures.
\end{savaistu}

\begin{aretenir}[L'essentiel de la section]
\begin{itemize}
\item Une suite récurrente $u_{n+1} = f(u_n)$ se visualise par une construction en escalier ($f$ croissante) ou en escargot ($f$ décroissante) entre la courbe de $f$ et la droite $y = x$.
\item Si $(u_n)$ converge vers $\ell$ et si $f$ est continue en $\ell$, alors $\ell = f(\ell)$ : la limite est un point fixe.
\item Intervalle stable $I$ contenant $u_0$ $\Rightarrow$ suite bornée (récurrence).
\item $f$ croissante sur $I$ stable $\Rightarrow$ $(u_n)$ monotone, de sens donné par le signe de $u_1 - u_0$.
\item Ordre impératif de la rédaction : bornitude $\to$ monotonie $\to$ convergence $\to$ équation aux points fixes $\to$ valeur de $\ell$.
\item $f$ décroissante : étudier $(u_{2n})$ et $(u_{2n+1})$, monotones de sens contraires.
\end{itemize}
\end{aretenir}

\coursec{Valeur approchée de la limite par l'inégalité des accroissements finis}

Dans la section précédente, nous avons appris à étudier une suite récurrente $u_{n+1} = f(u_n)$ : monotonie, bornitude, convergence, puis détermination de la limite $\ell$ comme point fixe de $f$, c'est-à-dire solution de l'équation $f(\ell) = \ell$. Mais le Baccalauréat va souvent plus loin : il ne demande pas seulement de \emph{prouver} que $(u_n)$ converge vers $\ell$, il demande de \emph{calculer} une valeur approchée de $\ell$ à une précision donnée, par exemple à $10^{-3}$ près. Autrement dit : à partir de quel rang $N$ peut-on affirmer que $u_N$ est une bonne approximation de $\ell$ ?

L'outil qui répond à cette question est une inégalité que tu as déjà rencontrée dans la leçon sur les applications de la dérivation : l'\cle{inégalité des accroissements finis}. Cette section est entièrement consacrée à son utilisation dans l'étude des suites récurrentes. C'est l'un des savoir-faire les plus rentables au Bac, car la démarche est toujours la même, quelle que soit la fonction $f$.

\courssub{Rappel : l'inégalité des accroissements finis}

Commençons par l'intuition. Si une fonction $f$ ne varie « pas trop vite » sur un intervalle $I$ — c'est-à-dire si sa dérivée reste petite en valeur absolue — alors deux images $f(x)$ et $f(y)$ ne peuvent pas être plus éloignées l'une de l'autre que $k$ fois la distance entre $x$ et $y$. La dérivée mesure la « vitesse de variation » : borner cette vitesse, c'est borner l'écart entre les images.

\begin{propriete}[Inégalité des accroissements finis (IAF)]
Soit $f$ une fonction dérivable sur un intervalle $I$. On suppose qu'il existe un réel $k \geq 0$ tel que, pour tout $t \in I$ :
\[ |f'(t)| \leq k. \]
Alors, pour tous réels $x$ et $y$ de $I$ :
\[ |f(x) - f(y)| \leq k\,|x - y|. \]
Ce résultat est \textbf{admis} ici : il est démontré dans la leçon sur les applications de la dérivation, comme conséquence du théorème des accroissements finis.
\end{propriete}

Géométriquement, dire que $|f'| \leq k$ signifie que la pente de toute tangente à la courbe de $f$ est comprise entre $-k$ et $k$. La corde joignant deux points de la courbe a donc, elle aussi, une pente dont la valeur absolue ne dépasse pas $k$.

\begin{popfigure}
\begin{center}
\begin{tikzpicture}[scale=1.6]
% axes
\draw[->, popmuted] (-0.3,0) -- (4.6,0) node[below right] {$x$};
\draw[->, popmuted] (0,-0.4) -- (0,3.1) node[above left] {$y$};
% courbe de f (pente modérée)
\draw[popblue, very thick, domain=0.2:4.2, samples=100] plot (\x, {0.45*\x + 0.35*sin(deg(1.6*\x)) + 0.4});
% points x et y
\coordinate (A) at (1, {0.45*1 + 0.35*sin(deg(1.6)) + 0.4});
\coordinate (B) at (3.4, {0.45*3.4 + 0.35*sin(deg(1.6*3.4)) + 0.4});
\fill[popink] (A) circle (1.6pt) node[above left, popink] {$(x,\ f(x))$};
\fill[popink] (B) circle (1.6pt) node[above right, popink] {$(y,\ f(y))$};
% corde
\draw[poporange, thick] (A) -- (B);
% projections
\draw[dashed, popmuted] (A) -- (1,0) node[below] {$x$};
\draw[dashed, popmuted] (B) -- (3.4,0) node[below] {$y$};
\draw[dashed, popmuted] (A) -- (0,{0.45 + 0.35*sin(deg(1.6)) + 0.4}) node[left] {$f(x)$};
\draw[dashed, popmuted] (B) -- (0,{0.45*3.4 + 0.35*sin(deg(1.6*3.4)) + 0.4}) node[left] {$f(y)$};
% droites de pente k et -k pour comparaison
\draw[popgreen, thick, dashed] (0.4,0.35) -- (2.4,2.15) node[midway, below left, popgreen] {pente $k$};
\draw[popgreen, thick, dashed] (2.6,2.6) -- (4.4,0.8) node[midway, above right, popgreen] {pente $-k$};
\end{tikzpicture}
\end{center}
\legende{Si toutes les tangentes ont une pente entre $-k$ et $k$, alors la corde (en orange) a une pente dont la valeur absolue est au plus $k$ : c'est le contenu géométrique de l'IAF.}
\end{popfigure}

\begin{attention}[Conditions d'application]
L'IAF exige que la majoration $|f'(t)| \leq k$ soit valable pour \textbf{tout} $t$ de l'intervalle $I$, et que $x$ et $y$ appartiennent tous les deux à $I$. Dans les exercices de suites, on choisit presque toujours pour $I$ un intervalle \emph{stable} par $f$ (c'est-à-dire tel que $f(I) \subset I$) qui contient à la fois le premier terme $u_0$ et la limite $\ell$.
\end{attention}

\courssub{Application aux suites récurrentes : convergence géométrique}

Voici l'idée centrale de toute cette section. Supposons que $(u_n)$ soit définie par $u_{n+1} = f(u_n)$, qu'elle converge vers $\ell$, et que $f(\ell) = \ell$. Alors l'écart entre $u_{n+1}$ et la limite s'écrit :
\[ |u_{n+1} - \ell| = |f(u_n) - f(\ell)|. \]
Or l'IAF, appliquée avec $x = u_n$ et $y = \ell$, donne :
\[ |f(u_n) - f(\ell)| \leq k\,|u_n - \ell|. \]
Autrement dit : \textbf{à chaque étape, la distance à la limite est multipliée par un facteur au plus $k$}. Si $k < 1$, cette distance diminue au moins aussi vite qu'une suite géométrique de raison $k$ : on parle de \cle{convergence géométrique} (ou convergence à vitesse géométrique). C'est exactement le principe qui fait « zoomer » la suite vers sa limite.

\begin{propriete}[Majoration géométrique de l'écart — démonstration exigible]
Soit $f$ une fonction dérivable sur un intervalle $I$ \textbf{stable} par $f$, et $(u_n)$ la suite définie par $u_0 \in I$ et $u_{n+1} = f(u_n)$. On suppose que :
\begin{itemize}
\item il existe $k \in [0\,;\,1[$ tel que $|f'(t)| \leq k$ pour tout $t \in I$ ;
\item l'équation $f(\ell) = \ell$ admet une solution $\ell \in I$.
\end{itemize}
Alors, pour tout entier naturel $n$ :
\[ |u_n - \ell| \leq k^n\,|u_0 - \ell|. \]
En particulier, comme $0 \leq k < 1$, on a $\lim\limits_{n \to +\infty} k^n = 0$, donc $(u_n)$ converge vers $\ell$.
\end{propriete}

\begin{experience}[Démonstration par récurrence (rédaction exigible au Bac)]
Démontrons par récurrence sur $n$ la propriété $\mathcal{P}(n)$ : « $|u_n - \ell| \leq k^n |u_0 - \ell|$ ».

\textbf{Initialisation.} Pour $n = 0$ : $|u_0 - \ell| = k^0 |u_0 - \ell|$ puisque $k^0 = 1$. La propriété $\mathcal{P}(0)$ est vraie (c'est même une égalité).

\textbf{Hérédité.} Soit $n \in \mathbb{N}$. Supposons $\mathcal{P}(n)$ vraie, c'est-à-dire $|u_n - \ell| \leq k^n |u_0 - \ell|$, et montrons $\mathcal{P}(n+1)$.

Comme $I$ est stable par $f$ et que $u_0 \in I$, tous les termes $u_n$ appartiennent à $I$ ; en particulier $u_n \in I$ et $\ell \in I$. On peut donc appliquer l'IAF entre $u_n$ et $\ell$ :
\begin{align*}
|u_{n+1} - \ell| &= |f(u_n) - f(\ell)| & &\text{car } u_{n+1} = f(u_n) \text{ et } \ell = f(\ell)\\
&\leq k\,|u_n - \ell| & &\text{par l'IAF, puisque } |f'| \leq k \text{ sur } I\\
&\leq k \times k^n |u_0 - \ell| & &\text{par l'hypothèse de récurrence}\\
&= k^{n+1} |u_0 - \ell|. & &
\end{align*}
La propriété $\mathcal{P}(n+1)$ est donc vraie.

\textbf{Conclusion.} Par le principe de récurrence, pour tout $n \in \mathbb{N}$ : $|u_n - \ell| \leq k^n |u_0 - \ell|$. Comme $0 \leq k < 1$, $k^n \to 0$, et le théorème des gendarmes assure que $|u_n - \ell| \to 0$, c'est-à-dire $u_n \to \ell$. $\blacksquare$
\end{experience}

\begin{attention}[Le piège du $k \geq 1$]
C'est l'erreur la plus fréquente au Bac. Si le meilleur majorant de $|f'|$ que tu obtiens est $k \geq 1$, alors l'inégalité $|u_n - \ell| \leq k^n |u_0 - \ell|$ reste vraie, mais elle \textbf{ne prouve rien} : le majorant $k^n |u_0 - \ell|$ explose, il ne tend pas vers $0$. On ne peut pas conclure à la convergence par cette méthode. Dans ce cas, il faut soit restreindre l'intervalle $I$ pour obtenir un $k < 1$, soit utiliser une autre méthode (monotonie + bornitude, par exemple). Retiens : \textbf{l'IAF ne donne la convergence que si $0 \leq k < 1$}.
\end{attention}

\courssub{Valeur approchée de la limite à une précision donnée}

L'inégalité $|u_n - \ell| \leq k^n |u_0 - \ell|$ fait mieux que prouver la convergence : elle \textbf{quantifie} l'erreur. Si l'on veut une valeur approchée de $\ell$ à $\varepsilon$ près (par exemple $\varepsilon = 10^{-3}$), il suffit de trouver un rang $N$ tel que :
\[ k^N\,|u_0 - \ell| \leq \varepsilon. \]
Alors, pour ce rang, $|u_N - \ell| \leq \varepsilon$ : le terme $u_N$, que l'on sait calculer explicitement, est une valeur approchée de l'inconnue $\ell$ à $\varepsilon$ près. C'est remarquable : on encadre une limite que l'on ne connaît pas encore !

\begin{methode}[Obtenir une valeur approchée de $\ell$ à $\varepsilon$ près — 4 étapes]
\begin{enumerate}
\item \textbf{Trouver $k$.} Montrer que $|f'(t)| \leq k$ avec $0 \leq k < 1$ sur un intervalle $I$ stable contenant $u_0$ et $\ell$.
\item \textbf{Majorer l'écart initial.} Déterminer une majoration de $|u_0 - \ell|$ (souvent donnée par un encadrement de $\ell$ obtenu précédemment).
\item \textbf{Résoudre l'inéquation.} Chercher le plus petit entier $N$ tel que $k^N |u_0 - \ell| \leq \varepsilon$, à la calculatrice (table de valeurs de $k^n$) ou par les logarithmes.
\item \textbf{Conclure.} Calculer $u_N$ et affirmer : « $u_N$ est une valeur approchée de $\ell$ à $\varepsilon$ près », car $|u_N - \ell| \leq k^N |u_0 - \ell| \leq \varepsilon$.
\end{enumerate}
\end{methode}

\textbf{Résolution par les logarithmes (complément utile).} Pour résoudre $k^N \leq \varepsilon'$ avec $0 < k < 1$, on compose par la fonction $\ln$, qui est croissante :
\[ N \ln k \leq \ln \varepsilon'. \]
Attention : $\ln k < 0$ puisque $k < 1$, donc en divisant par $\ln k$ \textbf{le sens de l'inégalité change} :
\[ N \geq \frac{\ln \varepsilon'}{\ln k}. \]
Il suffit alors de prendre pour $N$ le plus petit entier supérieur ou égal à ce quotient. Ce complément, très pratique, est signalé car il fait appel au logarithme népérien, étudié dans la leçon correspondante ; à défaut, une simple table de valeurs à la calculatrice fait parfaitement l'affaire au Bac.

\begin{exemplebox}[Recherche du rang avec $k = \dfrac{1}{5}$]
Supposons que l'on ait établi $|u_n - \ell| \leq \left(\dfrac{1}{5}\right)^n |u_0 - \ell|$ avec $|u_0 - \ell| \leq 1$. On cherche $N$ tel que $u_N$ approche $\ell$ à $10^{-3}$ près. Il suffit que :
\[ \left(\frac{1}{5}\right)^N \leq 10^{-3} \quad \Longleftrightarrow \quad 5^N \geq 1000. \]
Calculons les puissances de $5$ : $5^4 = 625 < 1000$ et $5^5 = 3125 \geq 1000$. Donc $N = 5$ convient (et $N = 4$ ne convient pas). Par les logarithmes : $N \geq \dfrac{\ln 10^{-3}}{\ln (1/5)} = \dfrac{-3\ln 10}{-\ln 5} = \dfrac{3\ln 10}{\ln 5} \approx 4{,}29$, d'où $N = 5$. Conclusion : $u_5$ est une valeur approchée de $\ell$ à $10^{-3}$ près.
\end{exemplebox}

\begin{popfigure}
\begin{center}
\begin{tikzpicture}
\begin{semilogyaxis}[
    width=0.85\linewidth, height=6.2cm,
    xlabel={rang $n$}, ylabel={$|u_n - \ell|$ (échelle logarithmique)},
    xmin=0, xmax=10, ymin=1e-7, ymax=2,
    grid=both, grid style={popline!40},
    legend style={at={(0.97,0.97)}, anchor=north east, font=\small}
]
\addplot[popblue, very thick, domain=0:10, samples=11, mark=*, mark size=1.5pt] {(0.2)^x};
\addlegendentry{majorant $k^n |u_0 - \ell|$ avec $k = \frac{1}{5}$, $|u_0-\ell|=1$}
\addplot[poporange, thick, dashed, domain=0:10, samples=2] {1e-3};
\addlegendentry{seuil de précision $\varepsilon = 10^{-3}$}
\addplot[popgreen, thick, dashed, domain=0:10, samples=2] {1e-6};
\addlegendentry{seuil $\varepsilon = 10^{-6}$}
\end{semilogyaxis}
\end{tikzpicture}
\end{center}
\legende{Décroissance géométrique de l'écart : en échelle semi-logarithmique, le majorant $k^n |u_0 - \ell|$ est une droite. Le seuil $10^{-3}$ est franchi dès $n = 5$, le seuil $10^{-6}$ dès $n = 9$ : chaque pas divise l'erreur par $5$.}
\end{popfigure}

\courssub{Exemple complet d'étude}

\begin{exoresolu}[Encadrement de la limite de $u_{n+1} = \dfrac{u_n + 3}{5}$ à $10^{-3}$ près]
On considère la suite $(u_n)$ définie par $u_0 = 0$ et, pour tout $n \in \mathbb{N}$, $u_{n+1} = \dfrac{u_n + 3}{5}$. On admet que $(u_n)$ converge vers l'unique solution $\ell$ de $\ell = \dfrac{\ell + 3}{5}$, et l'on a montré que $\ell = \dfrac{3}{4}$. On pose $f(x) = \dfrac{x+3}{5}$.

\textbf{1.} Montrer que pour tout $n \in \mathbb{N}$ : $|u_n - \ell| \leq \left(\dfrac{1}{5}\right)^n$.

\textbf{2.} En déduire un rang $N$ tel que $u_N$ soit une valeur approchée de $\ell$ à $10^{-3}$ près, puis donner cet encadrement.
\tcblower
\textbf{1.} La fonction $f$ est dérivable sur $\mathbb{R}$ et $f'(x) = \dfrac{1}{5}$ pour tout réel $x$. On a donc $|f'(x)| \leq k$ avec $k = \dfrac{1}{5} < 1$, sur $\mathbb{R}$ (qui est stable par $f$ et contient $u_0 = 0$ et $\ell = \tfrac{3}{4}$). Montrons par récurrence la propriété $\mathcal{P}(n)$ : $|u_n - \ell| \leq \left(\dfrac{1}{5}\right)^n$.

\emph{Initialisation} : pour $n = 0$, $|u_0 - \ell| = \left|0 - \dfrac{3}{4}\right| = \dfrac{3}{4} \leq 1 = \left(\dfrac{1}{5}\right)^0$. $\mathcal{P}(0)$ est vraie.

\emph{Hérédité} : supposons $\mathcal{P}(n)$ vraie pour un entier $n$ fixé. Alors, par l'IAF :
\begin{align*}
|u_{n+1} - \ell| &= |f(u_n) - f(\ell)| \leq \frac{1}{5}|u_n - \ell| \leq \frac{1}{5}\times\left(\frac{1}{5}\right)^n = \left(\frac{1}{5}\right)^{n+1}.
\end{align*}
Donc $\mathcal{P}(n+1)$ est vraie. Par récurrence, $\mathcal{P}(n)$ est vraie pour tout $n \in \mathbb{N}$.

\textbf{2.} On cherche $N$ tel que $\left(\dfrac{1}{5}\right)^N \leq 10^{-3}$, c'est-à-dire $5^N \geq 1000$. Comme $5^4 = 625 < 1000$ et $5^5 = 3125 \geq 1000$, le rang $N = 5$ convient.

Calculons les termes : $u_1 = \dfrac{3}{5} = 0{,}6$ ; $u_2 = \dfrac{0{,}6 + 3}{5} = 0{,}72$ ; $u_3 = \dfrac{0{,}72 + 3}{5} = 0{,}744$ ; $u_4 = 0{,}7488$ ; $u_5 = \dfrac{0{,}7488 + 3}{5} = 0{,}74976$.

Comme $|u_5 - \ell| \leq 10^{-3}$, on obtient l'encadrement :
\[ u_5 - 10^{-3} \leq \ell \leq u_5 + 10^{-3} \quad \Longleftrightarrow \quad 0{,}74876 \leq \ell \leq 0{,}75076. \]
Vérification : $\ell = \dfrac{3}{4} = 0{,}75$ appartient bien à cet intervalle, et $|u_5 - \ell| = 0{,}00024 \leq 10^{-3}$. La majoration par l'IAF était donc correcte (et même prudente).
\end{exoresolu}

\begin{attention}[Rédaction : ne pas confondre les deux inégalités]
Deux inégalités coexistent dans ces exercices, et les mélanger coûte des points :
\begin{itemize}
\item $|u_{n+1} - \ell| \leq k\,|u_n - \ell|$ : c'est le \textbf{passage d'un rang au suivant}, obtenu par l'IAF ;
\item $|u_n - \ell| \leq k^n |u_0 - \ell|$ : c'est la \textbf{conclusion globale}, obtenue par récurrence.
\end{itemize}
La première seule ne suffit pas à conclure : c'est la récurrence qui transforme le facteur $k$ en puissance $k^n$. Rédige toujours les deux étapes distinctement.
\end{attention}

\begin{savaistu}[Le point fixe de Banach : des maths qui compressent tes photos]
Le principe que tu viens d'étudier — « si $f$ contracte les distances d'un facteur $k < 1$, alors l'itération $u_{n+1} = f(u_n)$ converge vers l'unique point fixe de $f$ » — est un cas particulier du célèbre \textbf{théorème du point fixe de Banach} (1922), pilier de l'analyse moderne. Ce théorème garantit l'existence et l'unicité de solutions dans des situations très générales : équations différentielles, équations intégrales, problèmes d'optimisation. Une application étonnante : la \textbf{compression d'images par fractales}. L'idée est de coder une image non pas pixel par pixel, mais comme le point fixe d'une transformation contractante ; pour reconstituer l'image, on itère la transformation — exactement comme tu itères $f$ pour approcher $\ell$. Les mêmes mathématiques permettent aussi de résoudre numériquement les équations qui modélisent les réseaux de téléphonie mobile ou la diffusion de la chaleur. Quand tu rédiges « $|u_n - \ell| \leq k^n |u_0 - \ell|$ » au Bac, tu utilises un outil qui fait tourner une partie de l'industrie numérique mondiale !
\end{savaistu}

\begin{aretenir}[L'essentiel de la section]
\begin{itemize}
\item \textbf{IAF} : si $|f'| \leq k$ sur $I$, alors $|f(x) - f(y)| \leq k|x - y|$ pour tous $x, y \in I$.
\item Pour $u_{n+1} = f(u_n)$ avec $\ell = f(\ell)$ : $|u_{n+1} - \ell| \leq k|u_n - \ell|$, puis par récurrence $|u_n - \ell| \leq k^n |u_0 - \ell|$.
\item Si $0 \leq k < 1$ : convergence géométrique vers $\ell$. Si $k \geq 1$ : la méthode ne conclut pas.
\item Valeur approchée à $\varepsilon$ près : résoudre $k^N |u_0 - \ell| \leq \varepsilon$, éventuellement via $N \geq \dfrac{\ln(\varepsilon / |u_0 - \ell|)}{\ln k}$ (le sens change car $\ln k < 0$), puis calculer $u_N$.
\end{itemize}
\end{aretenir}

\begin{exercice}[À toi de jouer]
Soit $(v_n)$ la suite définie par $v_0 = 1$ et $v_{n+1} = \dfrac{v_n}{4} + 1$ pour tout $n \in \mathbb{N}$.
\begin{enumerate}
\item Résoudre l'équation $\ell = \dfrac{\ell}{4} + 1$.
\item Montrer que pour tout $n \in \mathbb{N}$ : $|v_n - \ell| \leq \left(\dfrac{1}{4}\right)^n |v_0 - \ell|$.
\item Déterminer un rang $N$ à partir duquel $v_N$ est une valeur approchée de $\ell$ à $10^{-2}$ près, puis donner l'encadrement correspondant.
\end{enumerate}
\end{exercice}

\begin{corrige}[Exercice : suite $v_{n+1} = \dfrac{v_n}{4} + 1$]
\textbf{1.} $\ell = \dfrac{\ell}{4} + 1 \iff \dfrac{3\ell}{4} = 1 \iff \ell = \dfrac{4}{3}$.

\textbf{2.} Posons $f(x) = \dfrac{x}{4} + 1$ ; alors $f'(x) = \dfrac{1}{4}$ pour tout $x \in \mathbb{R}$, donc $|f'| \leq k = \dfrac{1}{4} < 1$ sur $\mathbb{R}$. Récurrence sur $\mathcal{P}(n)$ : $|v_n - \ell| \leq \left(\dfrac14\right)^n |v_0 - \ell|$.

\emph{Initialisation} : pour $n = 0$, l'inégalité est une égalité. \emph{Hérédité} : si $\mathcal{P}(n)$ est vraie, alors par l'IAF :
\[ |v_{n+1} - \ell| = |f(v_n) - f(\ell)| \leq \frac{1}{4}|v_n - \ell| \leq \frac{1}{4}\left(\frac{1}{4}\right)^n |v_0 - \ell| = \left(\frac{1}{4}\right)^{n+1}|v_0 - \ell|. \]
Donc $\mathcal{P}(n+1)$ est vraie, et la propriété est établie pour tout $n$.

\textbf{3.} On a $|v_0 - \ell| = \left|1 - \dfrac{4}{3}\right| = \dfrac{1}{3}$. Il suffit que $\dfrac{1}{3}\left(\dfrac{1}{4}\right)^N \leq 10^{-2}$, c'est-à-dire $4^N \geq \dfrac{100}{3} \approx 33{,}3$. Or $4^2 = 16 < 33{,}3$ et $4^3 = 64 \geq 33{,}3$, donc $N = 3$ convient. On calcule : $v_1 = \dfrac{5}{4} = 1{,}25$ ; $v_2 = \dfrac{21}{16} = 1{,}3125$ ; $v_3 = \dfrac{85}{64} = 1{,}328125$. D'où l'encadrement $v_3 - 0{,}01 \leq \ell \leq v_3 + 0{,}01$, soit $1{,}318125 \leq \ell \leq 1{,}338125$. (Vérification : $\ell = \dfrac{4}{3} \approx 1{,}3333$ est bien dans cet intervalle.)
\end{corrige}

\coursec{Synthèse : méthode complète d'étude d'une suite au Bac}

Nous voici arrivés au terme de notre parcours sur les suites numériques. Dans la section précédente, tu as appris à \emph{encadrer numériquement} la limite d'une suite récurrente grâce à l'inégalité des accroissements finis : c'était la dernière pièce du puzzle. Il te reste maintenant à assembler toutes les pièces : récurrence, monotonie, bornitude, convergence, calcul de la limite, approximation. C'est exactement ce que demande un exercice de Baccalauréat : non pas une notion isolée, mais un \cle{raisonnement enchaîné}, où chaque étape justifie la suivante. Cette section est ta boîte à outils finale : lis-la, relis-la, et refais les exemples sans regarder les solutions.

\courssub{Le plan type d'un exercice de Bac sur les suites}

Un exercice de Bac sur les suites suit presque toujours la même architecture logique. La comprendre, c'est déjà gagner la moitié des points, car tu sais \emph{où va le sujet} avant même de lire toutes les questions.

\begin{methode}[Plan rédigé type d'un exercice complet de suite récurrente]
Voici les \cle{huit points de rédaction} attendus, dans l'ordre, pour l'étude d'une suite définie par $u_{n+1} = f(u_n)$ :
\begin{enumerate}
\item \textbf{Bornitude par récurrence.} Conjecturer un encadrement (par exemple $0 \le u_n \le 1$) en calculant les premiers termes, puis le démontrer par récurrence : initialisation, hérédité (en utilisant les variations de $f$), conclusion.
\item \textbf{Monotonie.} Étudier le signe de $u_{n+1} - u_n$ (souvent par récurrence, ou en factorisant à l'aide de la bornitude déjà prouvée), ou comparer $\dfrac{u_{n+1}}{u_n}$ à $1$ si tous les termes sont strictement positifs.
\item \textbf{Convergence.} Appliquer le théorème de convergence monotone : « la suite $(u_n)$ est croissante et majorée par $M$, donc elle converge vers un réel $\ell \le M$ » (ou décroissante et minorée).
\item \textbf{Calcul de la limite.} Passer à la limite dans la relation $u_{n+1} = f(u_n)$ : par continuité de $f$, la limite $\ell$ vérifie l'\cle{équation aux points fixes} $\ell = f(\ell)$. Résoudre cette équation.
\item \textbf{Choisir la bonne solution.} L'équation $\ell = f(\ell)$ peut avoir plusieurs solutions : éliminer celles qui contredisent la bornitude ou la monotonie (par exemple $\ell$ doit appartenir à l'intervalle où vivent tous les $u_n$).
\item \textbf{Vitesse de convergence (si demandé).} Introduire une suite auxiliaire ou montrer $|u_{n+1} - \ell| \le k\,|u_n - \ell|$ avec $0 \le k < 1$, puis en déduire $|u_n - \ell| \le k^n |u_0 - \ell|$.
\item \textbf{Approximation.} Déterminer un rang $n_0$ tel que $k^{n_0}|u_0 - \ell| < \varepsilon$ (par tâtonnement ou avec les puissances de $k$), et conclure que $u_{n_0}$ est une valeur approchée de $\ell$ à $\varepsilon$ près.
\item \textbf{Conclusion rédigée.} Terminer par une phrase complète : « Ainsi, pour tout $n \ge n_0$, $u_n$ approche $\ell$ à moins de $\varepsilon$ près. »
\end{enumerate}
\end{methode}

\begin{attention}[L'ordre logique est sacré]
L'erreur la plus coûteuse au Bac : \textbf{calculer $\ell$ avant d'avoir prouvé que la suite converge}. Résoudre $\ell = f(\ell)$ ne prouve \emph{rien} sur la convergence : cela donne seulement les \emph{candidats} possibles pour la limite. Une suite peut très bien vérifier $u_{n+1} = f(u_n)$ et diverger, alors que l'équation $\ell = f(\ell)$ possède des solutions !

Contre-exemple : $u_{n+1} = 2u_n$ avec $u_0 = 1$. L'équation $\ell = 2\ell$ donne $\ell = 0$, pourtant $u_n = 2^n \to +\infty$. Rédiger « la limite est $0$ » serait doublement faux.

Retiens la règle d'or : \cle{on prouve d'abord la convergence, on calcule ensuite la limite}. Au Bac, inverser cet ordre fait perdre des points même si la valeur trouvée est correcte.
\end{attention}

\courssub{La technique de la suite auxiliaire géométrique}

\textbf{Intuition.} Certaines suites récurrentes ne sont ni arithmétiques ni géométriques, et leur étude directe est délicate. L'astuce du mathématicien consiste à \emph{changer de point de vue} : au lieu d'étudier $(u_n)$, on étudie une suite $(v_n)$ construite à partir de $(u_n)$ — par exemple $v_n = u_n - a$ ou $v_n = \dfrac{u_n - a}{u_n - b}$ — choisie de sorte que $(v_n)$ soit \cle{géométrique}. Or les suites géométriques, on les connaît par cœur : $v_n = v_0 \times q^n$. On en déduit alors une expression explicite de $u_n$, et tout devient facile : monotonie, convergence, limite.

\begin{propriete}[Suite auxiliaire géométrique]
Soit $(u_n)$ une suite et $(v_n)$ définie par une relation $v_n = g(u_n)$ où $g$ est une fonction bien choisie (souvent indiquée par l'énoncé). Si l'on montre qu'il existe un réel $q$ tel que, pour tout entier $n$,
\[ v_{n+1} = q \times v_n, \]
alors $(v_n)$ est géométrique de raison $q$, donc $v_n = v_0 \, q^n$ pour tout $n$. Si de plus $|q| < 1$, alors $(v_n)$ converge vers $0$.
\end{propriete}

\textbf{Démonstration sur un exemple.} Posons $u_0 = 5$ et $u_{n+1} = \dfrac{1}{2}u_n + 3$. L'équation aux points fixes $\ell = \dfrac{1}{2}\ell + 3$ donne $\ell = 6$. Posons $v_n = u_n - 6$. Alors :
\[ v_{n+1} = u_{n+1} - 6 = \frac{1}{2}u_n + 3 - 6 = \frac{1}{2}u_n - 3 = \frac{1}{2}(u_n - 6) = \frac{1}{2}v_n. \]
Donc $(v_n)$ est géométrique de raison $\dfrac{1}{2}$, avec $v_0 = 5 - 6 = -1$. Ainsi $v_n = -\left(\dfrac{1}{2}\right)^n$, c'est-à-dire
\[ u_n = 6 - \left(\frac{1}{2}\right)^n. \]
On lit alors tout sur cette formule : $(u_n)$ est croissante, majorée par $6$, et converge vers $6$. La soustraction du point fixe a « tué » la constante $3$ et révélé la structure géométrique cachée. \hfill$\blacksquare$

\begin{propriete}[Convergence géométrique vers la limite]
Soit $(u_n)$ une suite convergente vers $\ell$. S'il existe un réel $k$ avec $0 \le k < 1$ tel que, pour tout entier $n$,
\[ |u_{n+1} - \ell| \le k\,|u_n - \ell|, \]
alors, pour tout entier $n$ :
\[ |u_n - \ell| \le k^n \, |u_0 - \ell|. \]
En particulier, comme $k^n \to 0$, l'écart $|u_n - \ell|$ tend vers $0$ au moins aussi vite qu'une suite géométrique de raison $k$.
\end{propriete}

\textbf{Démonstration (exigible, c'est une récurrence immédiate).} Notons $(P_n)$ : « $|u_n - \ell| \le k^n |u_0 - \ell|$ ».
\begin{itemize}
\item \emph{Initialisation} : pour $n = 0$, $|u_0 - \ell| \le k^0 |u_0 - \ell| = |u_0 - \ell|$ : vrai.
\item \emph{Hérédité} : supposons $(P_n)$ vraie. Alors
\[ |u_{n+1} - \ell| \le k\,|u_n - \ell| \le k \times k^n |u_0 - \ell| = k^{n+1}|u_0 - \ell|, \]
donc $(P_{n+1})$ est vraie.
\end{itemize}
Par récurrence, $(P_n)$ est vraie pour tout $n \in \mathbb{N}$. \hfill$\blacksquare$

C'est cette inégalité qui permet, comme tu l'as vu dans la section précédente, de trouver un rang à partir duquel $u_n$ approche $\ell$ avec une précision donnée : il suffit de résoudre $k^n |u_0 - \ell| < \varepsilon$.

\courssub{Exemple chiffré : une étude complète rédigée}

Voici le cœur de cette section : un exercice complet, rédigé comme au Bac, utilisant la technique de la suite auxiliaire.

\begin{exoresolu}[Étude complète d'une suite homographique]
On considère la suite $(u_n)$ définie par $u_0 = 0$ et, pour tout $n \in \mathbb{N}$ :
\[ u_{n+1} = \frac{2u_n + 3}{u_n + 4}. \]
\begin{enumerate}
\item Démontrer que pour tout $n$, $0 \le u_n \le 1$.
\item On pose $v_n = \dfrac{u_n - 1}{u_n + 3}$. Montrer que $(v_n)$ est géométrique ; exprimer $v_n$ puis $u_n$ en fonction de $n$.
\item Étudier la convergence de $(u_n)$ et donner sa limite.
\item Déterminer un rang $n_0$ tel que $|u_{n_0} - \ell| < 10^{-3}$.
\end{enumerate}
\tcblower
\textbf{1. Bornitude par récurrence.} Notons $(P_n)$ : « $0 \le u_n \le 1$ ».

\emph{Initialisation} : $u_0 = 0$, donc $0 \le u_0 \le 1$ : $(P_0)$ est vraie.

\emph{Hérédité} : supposons $0 \le u_n \le 1$. La fonction $f(x) = \dfrac{2x+3}{x+4}$ est dérivable sur $[0\,;1]$ et $f'(x) = \dfrac{2(x+4) - (2x+3)}{(x+4)^2} = \dfrac{5}{(x+4)^2} > 0$ : $f$ est strictement croissante sur $[0\,;1]$. Donc :
\[ f(0) \le f(u_n) \le f(1) \quad\Longrightarrow\quad \frac{3}{4} \le u_{n+1} \le \frac{5}{5} = 1, \]
et a fortiori $0 \le u_{n+1} \le 1$ : $(P_{n+1})$ est vraie.

Par récurrence, pour tout $n \in \mathbb{N}$, $0 \le u_n \le 1$.

\textbf{2. Suite auxiliaire.} Comme $u_n \ge 0$, on a $u_n + 3 \neq 0$ : $(v_n)$ est bien définie. Calculons :
\begin{align*}
v_{n+1} &= \frac{u_{n+1} - 1}{u_{n+1} + 3}
= \frac{\dfrac{2u_n + 3}{u_n + 4} - 1}{\dfrac{2u_n + 3}{u_n + 4} + 3}
= \frac{2u_n + 3 - (u_n + 4)}{2u_n + 3 + 3(u_n + 4)} \\
&= \frac{u_n - 1}{5u_n + 15}
= \frac{1}{5} \times \frac{u_n - 1}{u_n + 3}
= \frac{1}{5}\,v_n.
\end{align*}
Donc $(v_n)$ est géométrique de raison $q = \dfrac{1}{5}$, de premier terme $v_0 = \dfrac{0 - 1}{0 + 3} = -\dfrac{1}{3}$. Ainsi :
\[ v_n = -\frac{1}{3}\left(\frac{1}{5}\right)^n. \]
Or $v_n = \dfrac{u_n - 1}{u_n + 3}$ équivaut à $v_n(u_n + 3) = u_n - 1$, soit $u_n(v_n - 1) = -1 - 3v_n$, d'où :
\[ u_n = \frac{1 + 3v_n}{1 - v_n} = \frac{1 - \left(\frac{1}{5}\right)^n}{1 + \frac{1}{3}\left(\frac{1}{5}\right)^n}. \]

\textbf{3. Convergence.} Comme $\left|\dfrac{1}{5}\right| < 1$, on a $\left(\dfrac{1}{5}\right)^n \to 0$, donc $v_n \to 0$ et, par opérations sur les limites :
\[ \lim_{n \to +\infty} u_n = \frac{1 - 0}{1 + 0} = 1. \]
La suite $(u_n)$ converge vers $\ell = 1$ (qui est bien le point fixe de $f$ dans $[0\,;1]$, puisque $\ell = \dfrac{2\ell+3}{\ell+4}$ donne $\ell^2 + 2\ell - 3 = 0$, soit $\ell = 1$ ou $\ell = -3$, et $-3$ est exclu car $u_n \ge 0$).

\textbf{4. Approximation.} On a $|u_n - 1| = \left|\dfrac{1 + 3v_n}{1 - v_n} - 1\right| = \left|\dfrac{4v_n}{1 - v_n}\right|$. Pour $n \ge 1$, $|v_n| \le \dfrac{1}{15}$, donc $|1 - v_n| \ge \dfrac{14}{15}$ et :
\[ |u_n - 1| \le \frac{4 \times \frac{1}{3}\left(\frac{1}{5}\right)^n}{\frac{14}{15}} = \frac{10}{7}\left(\frac{1}{5}\right)^n. \]
Il suffit que $\dfrac{10}{7}\left(\dfrac{1}{5}\right)^n < 10^{-3}$, c'est-à-dire $5^n > \dfrac{10^4}{7} \approx 1428$. Or $5^4 = 625$ et $5^5 = 3125$. Donc $n_0 = 5$ convient : $u_5$ est une valeur approchée de $1$ à $10^{-3}$ près.
\end{exoresolu}

On remarque l'enchaînement : la question 1 (bornitude) garantit que $(v_n)$ est bien définie à la question 2 ; la question 2 donne la forme explicite qui rend la question 3 immédiate ; la question 3 fournit la limite dont on mesure l'écart à la question 4. \cle{Chaque question prépare la suivante} : c'est la signature d'un sujet de Bac bien construit.

\courssub{Quel outil pour quelle question ?}

Face à une question, le réflexe du bon élève est de se demander : « quel outil du chapitre répond exactement à ceci ? » Voici le tableau de correspondance à mémoriser.

\begin{center}
\begin{tabularx}{\linewidth}{|l|X|}
\hline
\rowcolor{popblueL} \textbf{Question posée} & \textbf{Outil à mobiliser} \\
\hline
Montrer que $u_n \in [a\,;b]$ pour tout $n$ & Raisonnement par récurrence (hérédité via la croissance de $f$) \\
\hline
Montrer que $u_{n+1} \ge u_n$ pour tout $n$ & Signe de $u_{n+1} - u_n$, ou récurrence si $f$ est croissante \\
\hline
Justifier que $(u_n)$ converge & Théorème de convergence monotone (croissante majorée / décroissante minorée) \\
\hline
Calculer la limite $\ell$ & Équation aux points fixes $\ell = f(\ell)$ + continuité de $f$ (après convergence prouvée) \\
\hline
Exprimer $u_n$ en fonction de $n$ & Suite auxiliaire géométrique ($v_n = u_n - \ell$ ou forme indiquée) \\
\hline
Montrer $|u_{n+1} - \ell| \le k|u_n - \ell|$ & Inégalité des accroissements finis si $|f'| \le k < 1$ \\
\hline
Trouver $n$ tel que $|u_n - \ell| < \varepsilon$ & Inégalité $|u_n - \ell| \le k^n |u_0 - \ell|$ puis tâtonnement sur $k^n$ \\
\hline
\end{tabularx}
\end{center}

Et lorsque l'énoncé ne précise pas la méthode, c'est à toi de choisir. L'organigramme suivant te guide.

\begin{popfigure}
\begin{center}
\begin{tikzpicture}[
  font=\small,
  decision/.style={diamond, draw=poporange, fill=poporangeL, text width=3.4cm, align=center, inner sep=1pt, aspect=2.2},
  bloc/.style={rectangle, rounded corners, draw=popblue, fill=popblueL, text width=3.6cm, align=center, inner sep=4pt},
  final/.style={rectangle, rounded corners, draw=popgreen, fill=popgreenL, text width=3.6cm, align=center, inner sep=4pt},
  fleche/.style={-{Stealth}, thick, popdark}
]
\node[decision] (q1) {La suite est-elle définie par $u_{n+1}=f(u_n)$ ?};
\node[bloc, below left=1.3cm and 0.2cm of q1] (ag) {Forme explicite $u_n = g(n)$ : étudier directement $g$ (variations, limites)};
\node[decision, below right=1.3cm and 0.2cm of q1] (q2) {L'énoncé donne-t-il une suite auxiliaire $v_n$ ?};
\node[final, below left=1.3cm and 0.2cm of q2] (geo) {Montrer $v_{n+1}=q\,v_n$ : $(v_n)$ géométrique, en déduire $u_n$};
\node[decision, below right=1.3cm and 0.2cm of q2] (q3) {$f$ est-elle croissante sur un intervalle stable ?};
\node[final, below left=1.3cm and 0.2cm of q3] (mono) {Récurrence pour la bornitude, monotonie, théorème de convergence monotone};
\node[final, below right=1.3cm and 0.2cm of q3] (iaf) {$|f'|\le k<1$ : IAF pour $|u_{n+1}-\ell|\le k|u_n-\ell|$, convergence et approximation};
\draw[fleche] (q1) -- node[above left]{non} (ag);
\draw[fleche] (q1) -- node[above right]{oui} (q2);
\draw[fleche] (q2) -- node[above left]{oui} (geo);
\draw[fleche] (q2) -- node[above right]{non} (q3);
\draw[fleche] (q3) -- node[above left]{oui} (mono);
\draw[fleche] (q3) -- node[above right]{non} (iaf);
\end{tikzpicture}
\end{center}
\legende{Arbre de décision : quelle méthode choisir selon la forme de la suite ?}
\end{popfigure}

\courssub{Entraînement à la rédaction : l'enchaînement logique}

Au Bac, la rédaction compte autant que le résultat. Voici les connecteurs logiques qui structurent une copie exemplaire :
\begin{itemize}
\item « \emph{Montrons par récurrence que pour tout $n \in \mathbb{N}$, ...} » — annonce claire de la méthode ;
\item « \emph{La suite $(u_n)$ est donc croissante et majorée par $M$ ; d'après le théorème de convergence monotone, elle converge vers un réel $\ell \le M$.} » — citation explicite du théorème ;
\item « \emph{La fonction $f$ étant continue sur $[a\,;b]$, en passant à la limite dans $u_{n+1} = f(u_n)$, on obtient $\ell = f(\ell)$.} » — justification du passage à la limite ;
\item « \emph{Or tous les termes appartiennent à $[a\,;b]$, donc $\ell \in [a\,;b]$ : la seule solution possible est $\ell = ...$} » — sélection argumentée de la solution.
\end{itemize}

\begin{exercice}[Pour t'entraîner]
Soit $(u_n)$ la suite définie par $u_0 = 2$ et, pour tout entier $n \ge 0$,
\[ u_{n+1} = \frac{1}{3}u_n + 2. \]
\begin{enumerate}
\item Déterminer un encadrement de la forme $a \le u_n \le b$ valable pour tout $n \in \mathbb{N}$ et le démontrer par récurrence.
\item Étudier la monotonie de la suite $(u_n)$.
\item En déduire que $(u_n)$ converge et déterminer sa limite $\ell$.
\item On pose $v_n = u_n - 3$. Montrer que $(v_n)$ est géométrique, en déduire l'expression de $u_n$ en fonction de $n$ et retrouver la limite.
\item Déterminer un entier $n_0$ tel que $|u_{n_0} - \ell| < 10^{-2}$.
\end{enumerate}
\end{exercice}

\begin{corrige}[Exercice]
\textbf{1.} L'équation $\ell = \dfrac{\ell}{3} + 2$ donne $\ell = 3$. Calculons les premiers termes : $u_0=2$, $u_1=\frac{8}{3}\approx 2,67$, $u_2=\frac{26}{9}\approx 2,89$. La suite semble croissante et majorée par $3$. Montrons par récurrence $(P_n)$ : « $2 \le u_n \le 3$ ». 

\emph{Initialisation} : $u_0 = 2$, vrai. 

\emph{Hérédité} : si $2 \le u_n \le 3$, alors $u_{n+1} = \dfrac{u_n}{3} + 2$ vérifie $\dfrac{2}{3} + 2 \le u_{n+1} \le 1 + 2$, soit $\dfrac{8}{3} \le u_{n+1} \le 3$, et a fortiori $2 \le u_{n+1} \le 3$. 

Conclusion : $(P_n)$ vraie pour tout $n$.

\textbf{2.} $u_{n+1} - u_n = 2 - \dfrac{2u_n}{3} = \dfrac{2}{3}(3 - u_n) \ge 0$ car $u_n \le 3$ (d'après 1.). Donc $(u_n)$ est croissante.

\textbf{3.} $(u_n)$ est croissante et majorée par $3$ : par le théorème de convergence monotone, elle converge vers $\ell \le 3$. La fonction $f(x) = \dfrac{x}{3} + 2$ est continue, donc $\ell = \dfrac{\ell}{3} + 2$, soit $\dfrac{2\ell}{3} = 2$, d'où $\ell = 3$.

\textbf{4.} $v_{n+1} = u_{n+1} - 3 = \dfrac{u_n}{3} - 1 = \dfrac{1}{3}(u_n - 3) = \dfrac{1}{3}v_n$. Donc $(v_n)$ est géométrique de raison $\dfrac{1}{3}$, $v_0 = -1$, et $v_n = -\left(\dfrac{1}{3}\right)^n \to 0$, donc $u_n = 3 + v_n = 3 - \left(\dfrac{1}{3}\right)^n \to 3$.

\textbf{5.} $|u_n - 3| = \left(\dfrac{1}{3}\right)^n$. On cherche $\left(\dfrac{1}{3}\right)^n < 10^{-2}$, soit $3^n > 100$. Or $3^4 = 81$ et $3^5 = 243$. Donc $n_0 = 5$ convient.
\end{corrige}

\courssub{Les suites dans les sujets récents du Bac camerounais}

Depuis 2015, les épreuves de mathématiques des séries C, D, E et TI au Cameroun proposent quasiment chaque année un exercice sur les suites, très souvent structuré ainsi :
\begin{itemize}
\item une suite récurrente $u_{n+1} = f(u_n)$ avec $f$ homographique (type $\dfrac{au_n + b}{cu_n + d}$) ou affine ;
\item une suite auxiliaire $v_n$ \emph{donnée par l'énoncé}, à montrer géométrique ou arithmétique ;
\item l'expression de $u_n$ en fonction de $n$, puis le calcul de la limite ;
\item parfois, une dernière question d'approximation ou de somme de termes ($S_n = v_0 + \cdots + v_n$).
\end{itemize}
Les barèmes y consacrent typiquement de \cle{4 à 6 points} sur 20 : c'est un investissement rentable ! Les questions « Montrer que $(v_n)$ est géométrique » et « Exprimer $v_n$ puis $u_n$ en fonction de $n$ » sont des classiques quasi annuels : maîtrise-les parfaitement.

\begin{savaistu}[Les suites, un rendez-vous incontournable du Bac]
Une analyse des sujets du Baccalauréat camerounais (séries C et D notamment) entre 2015 et aujourd'hui montre que les suites numériques apparaissent dans la grande majorité des épreuves, souvent pour 4 à 6 points. Mieux : la technique de la suite auxiliaire géométrique y revient de façon récurrente — c'est le cas de le dire ! Au-delà de l'examen, les suites modélisent des situations bien concrètes : le remboursement d'un crédit auprès d'une microfinance, la croissance d'une population, la propagation d'une information sur les réseaux mobiles MTN ou Orange (chaque jour, une fraction des abonnés informés en informe de nouveaux)... Autant de phénomènes qui suivent des lois du type $u_{n+1} = f(u_n)$.
\end{savaistu}

\begin{aretenir}[L'essentiel de la synthèse]
\begin{itemize}
\item Un exercice de Bac sur les suites s'enchaîne toujours dans le même ordre : \cle{bornitude (récurrence)} $\to$ \cle{monotonie} $\to$ \cle{convergence (théorème monotone)} $\to$ \cle{limite (point fixe)} $\to$ \cle{approximation (IAF)}.
\item On ne calcule \textbf{jamais} la limite avant d'avoir prouvé la convergence.
\item La suite auxiliaire $v_n = u_n - \ell$ (ou toute forme proposée par l'énoncé) transforme une suite récurrente compliquée en suite géométrique : $v_n = v_0 q^n$.
\item Si $|u_{n+1} - \ell| \le k|u_n - \ell|$ avec $0 \le k < 1$, alors $|u_n - \ell| \le k^n |u_0 - \ell|$ : convergence garantie et approximation contrôlée.
\item Chaque question d'un sujet prépare la suivante : si tu bloques, relis les résultats déjà établis — la clé s'y trouve presque toujours.
\end{itemize}
\end{aretenir}

Tu disposes maintenant de l'ensemble des outils du chapitre « Suites numériques ». Le secret de la réussite au Bac n'est pas de connaître les théorèmes — tu les connais — mais de savoir \emph{dans quel ordre} les invoquer et comment les rédiger. Entraîne-toi sur les sujets des années passées en chronométrant ta rédaction : c'est ainsi que la méthode deviendra un réflexe. Bon courage !

\newpage

\coursec{Définition rigoureuse de la limite d'une suite (définition \og $\varepsilon$\fg{})}

Dans la leçon commune, tu as admis l'idée intuitive de limite : une suite $(u_n)$ converge vers un réel $\ell$ lorsque ses termes se rapprochent \og indéfiniment \fg{} de $\ell$. Cette intuition suffit pour conjecturer, mais pas pour démontrer proprement. Que signifie exactement \og se rapprocher \fg{} ? Les mathématiciens du XIX\textsuperscript{e} siècle, notamment Cauchy et Weierstrass, ont transformé cette image vague en une définition précise, quantifiée, inattaquable. C'est l'objet de ce complément, indispensable pour toute rédaction de démonstration rigoureuse au Baccalauréat.

\courssub{De l'intuition à la définition formelle}

Imaginons un défi : tu affirmes que la suite de terme général $u_n = 1 + \dfrac{1}{n}$ converge vers $1$. Un contradicteur te répond : \og Prouve-le. Je te donne une marge d'erreur $\varepsilon$, aussi petite que je veux, par exemple $\varepsilon = 0{,}001$, et tu dois me garantir qu'à partir d'un certain rang, tous les termes sont à moins de $\varepsilon$ de $1$. \fg{}

Ici, $|u_n - 1| = \dfrac{1}{n} < 0{,}001$ dès que $n > 1000$, c'est-à-dire dès le rang $N = 1001$. Tu gagnes le défi. Et comme tu peux le gagner pour \emph{toute} marge $\varepsilon$ proposée, la convergence est établie. C'est exactement l'esprit de la définition.

\begin{definition}[Limite d'une suite (définition par $\varepsilon$)]
On dit que la suite $(u_n)$ \cle{converge} vers le réel $\ell$, et on note $\lim\limits_{n \to +\infty} u_n = \ell$, lorsque :
\[
\forall \varepsilon > 0,\ \exists N \in \mathbb{N},\ \forall n \in \mathbb{N},\ n \geq N \implies |u_n - \ell| < \varepsilon.
\]
Autrement dit : pour toute marge d'erreur $\varepsilon$ strictement positive, on peut trouver un rang $N$ à partir duquel tous les termes de la suite sont dans l'intervalle $]\ell - \varepsilon\,;\,\ell + \varepsilon[$.
\end{definition}

\begin{propriete}[Unicité de la limite]
Si une suite converge, sa limite est unique.
\begin{proof}[Preuve (exigible au Bac)]
Supposons $\lim u_n = \ell$ et $\lim u_n = \ell'$ avec $\ell \neq \ell'$. Posons $\varepsilon = \frac{|\ell - \ell'|}{2} > 0$. D'après la définition, il existe $N_1$ tel que $\forall n \ge N_1, |u_n - \ell| < \varepsilon$ et $N_2$ tel que $\forall n \ge N_2, |u_n - \ell'| < \varepsilon$. Pour $n \ge \max(N_1, N_2)$, l'inégalité triangulaire donne :
\[
|\ell - \ell'| \le |u_n - \ell| + |u_n - \ell'| < \varepsilon + \varepsilon = |\ell - \ell'|,
\]
ce qui est impossible. Donc $\ell = \ell'$.
\end{proof}
\end{propriete}

\begin{aretenir}[Lecture de la définition]
\begin{itemize}
\item $\varepsilon$ est \textbf{donné d'abord}, arbitrairement petit ; $N$ est \textbf{trouvé ensuite}, et dépend de $\varepsilon$ : on écrit souvent $N(\varepsilon)$.
\item L'inégalité $|u_n - \ell| < \varepsilon$ signifie géométriquement que $u_n \in ]\ell - \varepsilon\,;\,\ell + \varepsilon[$.
\item Dire que $(u_n)$ converge vers $\ell$, c'est donc dire que tout intervalle ouvert centré en $\ell$, aussi étroit soit-il, contient \textbf{tous les termes de la suite à partir d'un certain rang}.
\end{itemize}
\end{aretenir}

\begin{popfigure}
\begin{tikzpicture}[scale=0.9]
\draw[->, popline] (0,0) -- (11,0) node[right] {$n$};
\draw[->, popline] (0,-0.5) -- (0,3) node[above] {$u_n$};
\draw[popblue, thick] (0,1.5) -- (10.8,1.5) node[right] {$\ell$};
\draw[poporange, dashed] (0,2.2) -- (10.8,2.2) node[right] {$\ell+\varepsilon$};
\draw[poporange, dashed] (0,0.8) -- (10.8,0.8) node[right] {$\ell-\varepsilon$};
\fill[poporange!15] (5.5,0.8) rectangle (10.8,2.2);
\draw[popgreen, thick] (5.5,0) -- (5.5,2.5) node[above] {$N$};
% Points "avant N" : un peu dispersés
\foreach \x/\y in {0.5/2.7, 1.5/0.5, 2.5/2.4, 3.5/0.9, 4.5/2.0}
  \fill[popink] (\x,\y) circle (1.8pt);
% Points "après N" : dans la bande, convergents
\foreach \x/\y in {5.8/1.6, 6.5/1.45, 7.2/1.55, 7.9/1.48, 8.6/1.52, 9.3/1.5, 10/1.51}
  \fill[popblue] (\x,\y) circle (1.8pt);
\node[popblue, anchor=west] at (10.2, 1.55) {$\bullet\ u_n$};
\node[popink, anchor=west] at (10.2, 1.35) {$\bullet\ u_n$ (avant $N$)};
\end{tikzpicture}
\legende{À partir du rang $N$, tous les points $(n\,;\,u_n)$ tombent dans la bande $]\ell-\varepsilon\,;\,\ell+\varepsilon[$.}
\end{popfigure}

\courssub{Démontrer une convergence avec la définition}

\begin{methode}[Rédaction type d'une preuve par $\varepsilon$]
Pour montrer que $(u_n)$ converge vers $\ell$ :
\begin{enumerate}
\item \textbf{Fixer} $\varepsilon > 0$ quelconque (écrire : \og Soit $\varepsilon > 0$ \fg).
\item \textbf{Majorer} la distance : exprimer ou majorer $|u_n - \ell|$ en fonction de $n$.
\item \textbf{Résoudre} l'inéquation $|u_n - \ell| < \varepsilon$ pour trouver un seuil en fonction de $\varepsilon$.
\item \textbf{Exhiber} un entier $N$ convenable (par exemple $N = \left\lfloor \dfrac{1}{\varepsilon} \right\rfloor + 1$) et \textbf{conclure} en citant la définition.
\end{enumerate}
\end{methode}

\begin{exoresolu}[Convergence de $u_n = \dfrac{2n+1}{n+3}$ vers $2$]
Soit $(u_n)$ définie pour $n \in \mathbb{N}$ par $u_n = \dfrac{2n+1}{n+3}$. Démontrer, à l'aide de la définition par $\varepsilon$, que $(u_n)$ converge vers $2$. Déterminer ensuite un rang $N$ pour $\varepsilon = 10^{-2}$.
\tcblower
\textbf{Étape 1.} Soit $\varepsilon > 0$.

\textbf{Étape 2.} Calculons la distance :
\[
|u_n - 2| = \left|\dfrac{2n+1}{n+3} - 2\right| = \left|\dfrac{2n+1 - 2n - 6}{n+3}\right| = \dfrac{5}{n+3}.
\]

\textbf{Étape 3.} On veut $\dfrac{5}{n+3} < \varepsilon$, c'est-à-dire $n + 3 > \dfrac{5}{\varepsilon}$, soit $n > \dfrac{5}{\varepsilon} - 3$.

\textbf{Étape 4.} Posons $N = \left\lfloor \dfrac{5}{\varepsilon} \right\rfloor + 1$. Alors pour tout $n \geq N$, on a $n > \dfrac{5}{\varepsilon} \ge \dfrac{5}{\varepsilon} - 3$, donc $|u_n - 2| < \varepsilon$. Par définition, $(u_n)$ converge vers $2$. $\blacksquare$

\textbf{Application numérique.} Pour $\varepsilon = 10^{-2}$ : $\dfrac{5}{\varepsilon} - 3 = 500 - 3 = 497$. On peut donc prendre $N = 498$ : à partir du rang $498$, tous les termes sont à moins de $0{,}01$ de $2$. Vérification : $|u_{498} - 2| = \dfrac{5}{501} \approx 0{,}00998 < 0{,}01$.
\end{exoresolu}

\begin{attention}[Pièges fréquents]
\begin{itemize}
\item \textbf{Ne pas inverser l'ordre des quantificateurs} : $N$ dépend de $\varepsilon$, jamais l'inverse. Écrire \og il existe $N$ tel que pour tout $\varepsilon > 0$... \fg{} est une erreur logique grave.
\item \textbf{Quelques termes proches ne suffisent pas} : il faut que \emph{tous} les termes vérifient la condition \emph{à partir} du rang $N$. La suite $u_n = (-1)^n$ passe infiniment souvent près de $1$, mais ne converge pas vers $1$.
\item \textbf{L'inégalité est stricte} : $|u_n - \ell| < \varepsilon$. En pratique, si tu obtiens $|u_n - \ell| \leq \varepsilon$, prends $N$ un peu plus grand (par exemple $N = \lfloor \dots \rfloor + 2$) pour retrouver une inégalité stricte.
\item La définition sert à \textbf{démontrer} une limite déjà conjecturée ; elle ne donne pas directement la valeur de $\ell$.
\end{itemize}
\end{attention}

\begin{savaistu}[Pourquoi cette définition a changé les mathématiques]
Avant Weierstrass (vers 1860), les notions de limite et d'infiniment petit reposaient sur des images floues, ce qui conduisait à des paradoxes. La définition par $\varepsilon$ a donné à l'analyse une fondation solide : c'est sur elle que reposent la continuité, la dérivation et l'intégration que tu étudies cette année. Au Cameroun comme ailleurs, c'est le passage de la \og recette de calcul \fg{} à la \og preuve rigoureuse \fg{} qui distingue le mathématicien du simple calculateur.
\end{savaistu}

\begin{exercice}[À toi de jouer]
Soit $(v_n)$ la suite définie par $v_n = 3 - \dfrac{2}{n^2}$ pour $n \geq 1$.
\begin{enumerate}
\item Conjecturer la limite de $(v_n)$.
\item Démontrer cette limite à l'aide de la définition par $\varepsilon$.
\item Déterminer un rang $N$ garantissant $|v_n - 3| < 10^{-4}$.
\end{enumerate}
\end{exercice}

\begin{corrige}[Exercice]
\begin{enumerate}
\item Puisque $\dfrac{2}{n^2} \to 0$ quand $n \to +\infty$, on conjecture $\lim\limits_{n \to +\infty} v_n = 3$.
\item Soit $\varepsilon > 0$. On a $|v_n - 3| = \left|-\dfrac{2}{n^2}\right| = \dfrac{2}{n^2}$.
On veut $\dfrac{2}{n^2} < \varepsilon \iff n^2 > \dfrac{2}{\varepsilon} \iff n > \sqrt{\dfrac{2}{\varepsilon}}$.
Posons $N = \left\lfloor \sqrt{\dfrac{2}{\varepsilon}} \right\rfloor + 1$. Pour tout $n \geq N$, on a $n > \sqrt{\dfrac{2}{\varepsilon}}$, donc $|v_n - 3| < \varepsilon$. Par définition, $(v_n)$ converge vers $3$.
\item Pour $\varepsilon = 10^{-4}$ : $\sqrt{\dfrac{2}{10^{-4}}} = \sqrt{20\,000} = 100\sqrt{2} \approx 141{,}42$. On prend $N = 142$.
\end{enumerate}
\end{corrige}

\newpage

\coursec{Activité d'intégration}

\coursec{Suites numériques : modélisation, convergence et encadrement de la limite}

\courssub{Objectifs de l'activité}

À l'issue de cette activité d'intégration, tu seras capable de :

\begin{itemize}
\item modéliser une situation réelle d'épargne par une \cle{suite récurrente} du type $u_{n+1}=f(u_n)$ ;
\item calculer les premiers termes d'une suite et émettre une \cle{conjecture} sur son comportement ;
\item démontrer par \cle{récurrence} la monotonie et la bornitude d'une suite ;
\item appliquer le \cle{théorème de convergence monotone} pour justifier qu'une suite converge ;
\item déterminer la limite d'une suite récurrente par l'\cle{équation du point fixe} ;
\item encadrer la limite à l'aide de l'\cle{inégalité des accroissements finis} (IAF) et donner un rang à partir duquel une précision souhaitée est atteinte ;
\item résoudre une inéquation faisant intervenir une suite géométrique pour répondre à une question concrète.
\end{itemize}

\courssub{Situation}

Dans le quartier Nkolndongo à Yaoundé, Maman Aïcha, commerçante au marché Mokolo, fait partie d'une \cle{tontine} (association d'épargne rotative très répandue au Cameroun). Soucieuse de préparer les études universitaires de sa fille, elle ouvre aussi une \og cagnotte numérique \fg{} sur son compte mobile money, qu'elle alimente selon une règle précise.

\textbf{Première règle d'épargne (cagnotte libre).} Chaque mois, le solde de sa cagnotte augmente de $2\,\%$ (grâce aux intérêts et aux petits bénéfices qu'elle y reverse), puis elle y ajoute un dépôt fixe de \num{5000} FCFA. Elle démarre avec un solde initial de \num{100000} FCFA.

On note $u_n$ le solde de la cagnotte, en FCFA, au bout de $n$ mois. On a donc :
\[ u_0 = \num{100000} \qquad \text{et} \qquad u_{n+1} = 1{,}02\,u_n + \num{5000}. \]

\textbf{Seconde règle d'épargne (tontine plafonnée).} Les membres de la tontine proposent à Maman Aïcha un dispositif solidaire : chaque mois, $10\,\%$ du solde est reversé au fonds social de la tontine (qui aide les membres en difficulté), puis Maman Aïcha dépose \num{50000} FCFA. Elle démarre là encore avec \num{100000} FCFA.

On note $v_n$ le solde, en FCFA, au bout de $n$ mois dans ce second dispositif. On a :
\[ v_0 = \num{100000} \qquad \text{et} \qquad v_{n+1} = 0{,}9\,v_n + \num{50000}. \]

\textbf{Problème posé.} Dans chaque dispositif, l'épargne de Maman Aïcha peut-elle croître sans limite, ou se stabilise-t-elle autour d'un montant ? Dans le second dispositif, peut-on estimer ce montant stabilisé à $100$ FCFA près ? Au bout de combien de mois le solde dépassera-t-il \num{300000} FCFA, montant nécessaire pour payer la première année d'université de sa fille ?

\courssub{Consignes}

Le travail se fait en groupes de quatre élèves, avec calculatrice. Chaque groupe rédige ses réponses sur une feuille, puis un rapporteur restitue les résultats au tableau.

\textbf{Partie A : la cagnotte libre $(u_n)$.}

\begin{enumerate}
\item Calculer $u_1$, $u_2$ et $u_3$ (arrondir à l'unité). Quelle conjecture peut-on émettre sur le sens de variation de $(u_n)$ ?
\item Démontrer par récurrence que, pour tout entier naturel $n$, $u_{n+1} \geq u_n$.
\item Démontrer par récurrence que, pour tout entier naturel $n$, $u_n \geq \num{100000}\times (1{,}02)^n$.
\item En déduire que la suite $(u_n)$ n'est pas majorée. Que peut-on conclure sur l'épargne de Maman Aïcha dans ce dispositif ?
\end{enumerate}

\textbf{Partie B : la tontine plafonnée $(v_n)$.}

\begin{enumerate}
\item Calculer $v_1$, $v_2$, $v_3$. Conjecturer le sens de variation et un éventuel plafond de $(v_n)$.
\item Démontrer par récurrence que, pour tout entier naturel $n$, $\num{100000} \leq v_n \leq \num{500000}$.
\item Démontrer par récurrence que la suite $(v_n)$ est croissante.
\item Justifier que $(v_n)$ converge et déterminer sa limite $\ell$ en résolvant l'équation du point fixe.
\item On pose $f(x)=0{,}9x+\num{50000}$. Montrer que pour tous réels $a$ et $b$, $|f(a)-f(b)| \leq 0{,}9\,|a-b|$, puis démontrer par récurrence que :
\[ |v_n - \num{500000}| \leq \num{400000}\times (0{,}9)^n. \]
\item En déduire le plus petit entier $n$ tel que $v_n$ soit une valeur approchée de la limite à $100$ FCFA près.
\item Résoudre l'inéquation $v_n \geq \num{300000}$ et répondre à la question de Maman Aïcha : au bout de combien de mois pourra-t-elle payer la première année d'université ?
\end{enumerate}

\courssub{Ressources à mobiliser}

\begin{aretenir}[Boîte à outils de la leçon]
\begin{itemize}
\item \cle{Raisonnement par récurrence} : pour prouver qu'une propriété $P(n)$ est vraie pour tout $n\in\mathbb{N}$, on vérifie l'\textbf{initialisation} ($P(0)$ vraie) et l'\textbf{hérédité} (si $P(n)$ est vraie, alors $P(n+1)$ est vraie).
\item \cle{Théorème de convergence monotone} : toute suite croissante et majorée converge ; toute suite décroissante et minorée converge. (Démonstration exigible au Bac).
\item \cle{Point fixe} : si $u_{n+1}=f(u_n)$ avec $f$ continue sur un intervalle $I$ stable par $f$, et si $(u_n)$ converge vers $\ell \in I$, alors $f(\ell)=\ell$.
\item \cle{Inégalité des accroissements finis (IAF)} : si $f$ est dérivable sur un intervalle $I$ et $|f'(x)|\leq k$ sur $I$, alors pour tous $a,b\in I$ : $|f(a)-f(b)|\leq k\,|a-b|$.
\item Pour $a>0$ et $q>0$ ($q\neq 1$) : $a\,q^n \leq M \iff n \geq \dfrac{\ln(M/a)}{\ln q}$ (attention au sens de l'inégalité si $\ln q < 0$, c'est-à-dire si $0<q<1$).
\end{itemize}
\end{aretenir}

\courssub{Démarche et corrigé commenté}

\begin{exoresolu}[La tontine numérique de Maman Aïcha — corrigé complet]
\textbf{Partie A.} 1) Calculer les premiers termes de $(u_n)$ et conjecturer. 2) Prouver la croissance par récurrence. 3) Prouver le minorant géométrique. 4) Conclure sur la bornitude. \textbf{Partie B.} 1) Conjecturer. 2) et 3) Prouver bornitude et croissance. 4) Conclure à la convergence et calculer $\ell$. 5) Établir l'encadrement par l'IAF. 6) Trouver le rang pour une précision de $100$ FCFA. 7) Déterminer le mois où le solde atteint \num{300000} FCFA.
\tcblower
\textbf{Partie A : la cagnotte libre.}

\textbf{1.} On calcule :
\begin{align*}
u_1 &= 1{,}02\times \num{100000} + \num{5000} = \num{107000},\\
u_2 &= 1{,}02\times \num{107000} + \num{5000} = \num{114140},\\
u_3 &= 1{,}02\times \num{114140} + \num{5000} = \num{121422{,}8} \approx \num{121423}.
\end{align*}
\emph{Conjecture :} la suite $(u_n)$ semble strictement croissante, et les écarts entre termes successifs augmentent.

\textbf{2.} Montrons par récurrence que $u_{n+1}\geq u_n$ pour tout $n\in\mathbb{N}$.
\emph{Initialisation ($n=0$) :} $u_1 = \num{107000} \geq \num{100000} = u_0$ : la propriété est vraie au rang $0$.
\emph{Hérédité :} supposons la propriété vraie au rang $n$, c'est-à-dire $u_{n+1}\geq u_n$.
Comme $1{,}02>0$, en multipliant les deux membres par $1{,}02$ on conserve l'inégalité :
$1{,}02\,u_{n+1} \geq 1{,}02\,u_n$.
En ajoutant \num{5000} des deux côtés :
$1{,}02\,u_{n+1} + \num{5000} \geq 1{,}02\,u_n + \num{5000}$, soit $u_{n+2} \geq u_{n+1}$.
La propriété est héréditaire. Par récurrence, $(u_n)$ est \cle{croissante}.

\textbf{3.} Montrons par récurrence que $u_n \geq \num{100000}\,(1{,}02)^n$ pour tout $n\in\mathbb{N}$.
\emph{Initialisation ($n=0$) :} $u_0 = \num{100000} = \num{100000}\,(1{,}02)^0$ : vrai.
\emph{Hérédité :} supposons $u_n \geq \num{100000}\,(1{,}02)^n$.
Alors :
\[ u_{n+1} = 1{,}02\,u_n + \num{5000} \geq 1{,}02\times \num{100000}\,(1{,}02)^n + \num{5000} = \num{100000}\,(1{,}02)^{n+1} + \num{5000} \geq \num{100000}\,(1{,}02)^{n+1}. \]
La propriété est héréditaire, donc vraie pour tout $n$.

\textbf{4.} Comme $1{,}02>1$, la suite géométrique $\bigl((1{,}02)^n\bigr)$ diverge vers $+\infty$ : pour tout réel $M$, il existe un rang $n$ tel que $\num{100000}\,(1{,}02)^n > M$, donc $u_n > M$.
La suite $(u_n)$ n'est \textbf{pas majorée} : elle diverge vers $+\infty$.
\emph{Interprétation :} dans ce dispositif, l'épargne de Maman Aïcha croît sans limite (théoriquement !). Remarquons que l'équation du point fixe $\ell = 1{,}02\,\ell + \num{5000}$ donne $\ell = -\num{250000}$, impossible pour un solde : cela confirme qu'il n'y a pas de limite finie.

\begin{popfigure}
\begin{tikzpicture}[scale=0.8, >=Stealth]
% Axes
\draw[->] (0,0) -- (6,0) node[right] {$n$};
\draw[->] (0,0) -- (0,5) node[above] {$u_n$ (en $10^5$ FCFA)};
% Points
\foreach \n/\val in {0/1, 1/1.07, 2/1.1414, 3/1.2142, 4/1.2885, 5/1.3643} {
    \fill[popblue] (\n, \val) circle (2pt);
    \draw[popblue, dashed] (\n, 0) node[below, font=\tiny] {\n} -- (\n, \val);
}
% Tendance
\draw[popblue, thick, domain=0:5.5, samples=50] plot (\x, {1*1.02^\x + 250*(1.02^\x - 1)}); % Approx explicit formula scaled
\node[popblue, right] at (5.5, 3.5) {Croissance exponentielle};
\draw[poporange, dashed] (0,0) -- (5.5, 0); % baseline
\end{tikzpicture}
\legende{Évolution du solde $u_n$ (en $10^5$ FCFA) : divergence vers $+\infty$.}
\end{popfigure}

\textbf{Partie B : la tontine plafonnée.}

\textbf{1.} 
\begin{align*}
v_1 &= 0{,}9\times\num{100000}+\num{50000} = \num{140000},\\
v_2 &= 0{,}9\times\num{140000}+\num{50000} = \num{176000},\\
v_3 &= 0{,}9\times\num{176000}+\num{50000} = \num{208400}.
\end{align*}
\emph{Conjecture :} $(v_n)$ est croissante et semble se stabiliser vers \num{500000} FCFA.

\textbf{2.} Montrons par récurrence que $\num{100000}\leq v_n \leq \num{500000}$ pour tout $n\in\mathbb{N}$.
\emph{Initialisation ($n=0$) :} $v_0=\num{100000}$, donc $\num{100000}\leq v_0 \leq \num{500000}$ : vrai.
\emph{Hérédité :} supposons $\num{100000}\leq v_n \leq \num{500000}$.
Comme $0{,}9>0$, en multipliant par $0{,}9$ : $\num{90000}\leq 0{,}9\,v_n \leq \num{450000}$.
En ajoutant \num{50000} : $\num{140000}\leq 0{,}9\,v_n + \num{50000} \leq \num{500000}$, c'est-à-dire $\num{140000}\leq v_{n+1}\leq \num{500000}$.
Or $\num{140000} \geq \num{100000}$, donc a fortiori $\num{100000}\leq v_{n+1}\leq \num{500000}$.
La propriété est héréditaire. La suite $(v_n)$ est donc \cle{bornée} (minorée par \num{100000} et majorée par \num{500000}).

\textbf{3.} Montrons par récurrence que $v_{n+1}\geq v_n$ pour tout $n\in\mathbb{N}$ (suite croissante).
\emph{Initialisation ($n=0$) :} $v_1=\num{140000}\geq\num{100000}=v_0$ : vrai.
\emph{Hérédité :} supposons $v_{n+1}\geq v_n$.
Alors $v_{n+2}=0{,}9\,v_{n+1}+\num{50000}\geq 0{,}9\,v_n+\num{50000}=v_{n+1}$.
La propriété est héréditaire. Donc $(v_n)$ est \cle{croissante}.

\begin{center}
\begin{tabular}{|c|cccccc|}
\hline
$n$ & $0$ & $1$ & $2$ & $3$ & $\cdots$ & $+\infty$ \\
\hline
$v_n$ & \num{100000} & \num{140000} & \num{176000} & \num{208400} & $\nearrow$ & \num{500000} \\
\hline
\end{tabular}
\end{center}
\emph{Tableau de variation de la suite $(v_n)$.}

\textbf{4.} La suite $(v_n)$ est croissante et majorée par \num{500000} : d'après le \cle{théorème de convergence monotone}, elle converge vers un réel $\ell$.
La fonction $f(x)=0{,}9x+\num{50000}$ étant continue (affine) sur $\mathbb{R}$, la limite $\ell$ vérifie l'équation du point fixe :
\[ \ell = 0{,}9\,\ell + \num{50000} \iff 0{,}1\,\ell = \num{50000} \iff \ell = \num{500000}. \]
L'épargne de Maman Aïcha se stabilise donc vers \textbf{\num{500000} FCFA}.

\begin{popfigure}
\begin{tikzpicture}[scale=1.2, >=Stealth, domain=1:5]
% Axes : x = v_n / 100000, y = v_{n+1} / 100000
\draw[->] (0.5,0.5) -- (5.5,0.5) node[right] {$v_n$ ($\times 10^5$ FCFA)};
\draw[->] (0.5,0.5) -- (0.5,5.5) node[above] {$v_{n+1}$ ($\times 10^5$ FCFA)};
% y = x
\draw[popmuted, thick] (0.5,0.5) -- (5.5,5.5) node[above right] {$y=x$};
% y = f(x) = 0.9x + 0.5
\draw[popblue, thick] (0.5, 0.95) -- (5.5, 5.45) node[above right] {$y=f(x)$};
% Point fixe
\fill[poporange] (5,5) circle (2pt) node[above left] {$\ell=5$};
% Toile d'araignée (Cobweb) starting from v0=1
\draw[popgreen, thick] (1, 0.5) -- (1, 1.4); % vertical to curve
\draw[popgreen, thick] (1, 1.4) -- (1.4, 1.4); % horizontal to diagonal
\draw[popgreen, thick] (1.4, 1.4) -- (1.4, 1.76); % vertical
\draw[popgreen, thick] (1.4, 1.76) -- (1.76, 1.76); % horizontal
\draw[popgreen, thick] (1.76, 1.76) -- (1.76, 2.084); % vertical
\draw[popgreen, thick] (2.084, 2.084) -- (2.084, 2.3756); % horizontal approx
% Labels
\node[popgreen, below] at (1, 0.5) {$v_0$};
\node[popgreen, left] at (0.5, 1.4) {$v_1$};
\node[popgreen, below] at (1.4, 1.4) {$v_1$};
\node[popgreen, left] at (0.5, 1.76) {$v_2$};
\end{tikzpicture}
\legende{Représentation graphique (toile d'araignée) de la convergence de $(v_n)$ vers $\ell = \num{500000}$ FCFA.}
\end{popfigure}

\textbf{5.} Pour tous réels $a$ et $b$ :
\[ |f(a)-f(b)| = |0{,}9a + \num{50000} - (0{,}9b + \num{50000})| = |0{,}9a - 0{,}9b| = 0{,}9\,|a-b|. \]
(C'est le cas d'égalité de l'IAF, puisque $f'(x)=0{,}9$ constant sur $\mathbb{R}$.)
Comme $f(\num{500000})=\num{500000}$, on a pour tout $n$ :
\[ |v_{n+1}-\num{500000}| = |f(v_n)-f(\num{500000})| \leq 0{,}9\,|v_n-\num{500000}|. \]
Montrons alors par récurrence que $|v_n-\num{500000}|\leq \num{400000}\,(0{,}9)^n$ pour tout $n\in\mathbb{N}$.
\emph{Initialisation ($n=0$) :} $|v_0-\num{500000}| = \num{400000} = \num{400000}\,(0{,}9)^0$ : vrai.
\emph{Hérédité :} si $|v_n-\num{500000}|\leq \num{400000}\,(0{,}9)^n$, alors
\[ |v_{n+1}-\num{500000}| \leq 0{,}9\,|v_n-\num{500000}| \leq 0{,}9 \times \num{400000}\,(0{,}9)^n = \num{400000}\,(0{,}9)^{n+1}. \]
L'encadrement est établi pour tout $n$.

\textbf{6.} On cherche le plus petit $n\in\mathbb{N}$ tel que $\num{400000}\,(0{,}9)^n \leq 100$, c'est-à-dire $(0{,}9)^n \leq \dfrac{1}{\num{4000}}$.
Comme $0<0{,}9<1$, $\ln(0{,}9)<0$. En appliquant le logarithme népérien (fonction croissante) :
\[ n \ln(0{,}9) \leq \ln(1/\num{4000}) = -\ln \num{4000}. \]
En divisant par $\ln(0{,}9)<0$, \textbf{le sens de l'inégalité change} :
\[ n \geq \frac{-\ln \num{4000}}{\ln 0{,}9}. \]
Numériquement : $\ln \num{4000} \approx 8{,}294$, $\ln 0{,}9 \approx -0{,}10536$, donc $n \geq \dfrac{8{,}294}{0{,}10536} \approx 78{,}72$.
Le plus petit entier convenant est $n=79$.
Vérification : $\num{400000}\times(0{,}9)^{79}\approx 97{,}3 \leq 100$, tandis que pour $n=78$ on obtient $\approx 108{,}1 > 100$.
À partir du \textbf{79\textsuperscript{ème} mois}, $v_n$ est une valeur approchée de la limite à $100$ FCFA près.

\textbf{7.} On détermine l'expression explicite de $v_n$. La suite $(v_n - \num{500000})$ est géométrique de raison $0{,}9$ et de premier terme $v_0 - \num{500000} = -\num{400000}$.
Donc $v_n - \num{500000} = -\num{400000}\,(0{,}9)^n$, soit :
\[ v_n = \num{500000} - \num{400000}\,(0{,}9)^n. \]
L'inéquation $v_n \geq \num{300000}$ équivaut à :
\[ \num{500000} - \num{400000}\,(0{,}9)^n \geq \num{300000} \iff \num{400000}\,(0{,}9)^n \leq \num{200000} \iff (0{,}9)^n \leq 0{,}5. \]
Comme $\ln(0{,}9)<0$ :
\[ n \geq \frac{\ln 0{,}5}{\ln 0{,}9} \approx \frac{-0{,}6931}{-0{,}10536} \approx 6{,}58. \]
Le plus petit entier convenant est $n=7$.
Vérification : $v_6 = \num{500000} - \num{400000}\,(0{,}9)^6 \approx \num{287424} < \num{300000}$ et $v_7 \approx \num{308681} \geq \num{300000}$.
\emph{Conclusion :} dès le \textbf{7\textsuperscript{ème} mois}, Maman Aïcha dispose de plus de \num{300000} FCFA et peut payer la première année d'université de sa fille.
\end{exoresolu}

\begin{attention}[Pièges rencontrés pendant l'activité]
\begin{itemize}
\item Ne pas conclure à la convergence d'une suite à partir des seuls premiers termes : une conjecture doit être \textbf{prouvée} (ici par récurrence).
\item L'équation du point fixe ne donne la limite que si l'on a \textbf{d'abord prouvé} que la suite converge. Pour $(u_n)$, elle donne $\ell=-\num{250000}$, ce qui est absurde : c'est le signe que la suite diverge.
\item Dans la division par $\ln(0{,}9)$, qui est \textbf{négatif}, le sens de l'inégalité change : c'est l'erreur la plus fréquente aux questions B.6 et B.7.
\item Ne pas confondre \og majorée \fg{} et \og convergente \fg{} : c'est la conjonction \og croissante \textbf{et} majorée \fg{} qui assure la convergence (théorème de convergence monotone).
\item Pour l'expression explicite de $v_n$, bien vérifier que $v_n = \ell + (v_0-\ell)\times 0{,}9^n$ et non $v_n = \ell - (v_0-\ell)\times 0{,}9^n$ (signe du premier terme de la suite géométrique associée).
\end{itemize}
\end{attention}

\courssub{Bilan de l'activité}

Cette activité a permis de mettre en évidence les points essentiels de la leçon :

\begin{itemize}
\item Une même démarche d'épargne, modifiée par un simple prélèvement proportionnel, peut transformer une suite \textbf{divergente} en une suite \textbf{convergente} : le coefficient multiplicateur ($1{,}02$ contre $0{,}9$) joue un rôle décisif (raison de la suite géométrique associée).
\item Le \cle{raisonnement par récurrence} est l'outil privilégié pour prouver rigoureusement la monotonie et la bornitude d'une suite définie par $u_{n+1}=f(u_n)$.
\item Le \cle{théorème de convergence monotone} garantit l'existence de la limite, et l'\cle{équation du point fixe} permet de la calculer (à condition que $f$ soit continue).
\item L'\cle{inégalité des accroissements finis} fournit un contrôle quantitatif de l'écart $|v_n-\ell|$, ce qui permet de donner une \textbf{valeur approchée} de la limite avec une précision choisie à l'avance — une compétence directement exigible au Baccalauréat.
\item Enfin, la résolution d'une inéquation du type $a\,q^n \leq M$ répond à une question concrète de planification financière : les mathématiques de Terminale éclairent des décisions de la vie quotidienne au Cameroun.
\end{itemize}

\begin{savaistu}[Les tontines, une mathématique sociale]
Les tontines, pratiquées depuis des générations dans tout le Cameroun et l'Afrique de l'Ouest, reposent sur des mécanismes d'accumulation et de redistribution que les mathématiciens modélisent aujourd'hui par des suites et des graphes. Avec l'essor du mobile money (MTN MoMo, Orange Money), ces épargnes traditionnelles deviennent numériques : comprendre la convergence des suites récurrentes, c'est aussi comprendre comment une cagnotte se stabilise... ou s'emballe ! Le coefficient directeur de la fonction $f$ (ici $0{,}9$ ou $1{,}02$) détermine si le système est \emph{contractant} (convergence) ou \emph{dilatant} (divergence).
\end{savaistu}

\newpage

\coursec{Méthodes et exercices résolus}

\courssub{Méthodes et exercices résolus}

\begin{methode}[Démontrer une propriété par récurrence]
\textbf{Objectif :} Prouver qu'une propriété $\mathcal{P}(n)$ est vraie pour tout entier naturel $n$ (ou $n \ge n_0$).
\begin{enumerate}
    \item \textbf{Initialisation :} Vérifier que $\mathcal{P}(n_0)$ est vraie (souvent $n_0=0$ ou $1$). Calculer explicitement les deux membres.
    \item \textbf{Hérédité :} Supposer $\mathcal{P}(k)$ vraie pour un $k \ge n_0$ arbitraire (\emph{hypothèse de récurrence}). En déduire rigoureusement que $\mathcal{P}(k+1)$ est vraie. \textbf{Astuce :} Partez de l'expression pour $k+1$ et faites apparaître l'expression pour $k$ pour utiliser l'hypothèse.
    \item \textbf{Conclusion :} Écrire la phrase rituelle : « Par le principe de récurrence, la propriété $\mathcal{P}(n)$ est vraie pour tout entier $n \ge n_0$. »
\end{enumerate}
\cle{Réflexe :} Pour une suite définie par récurrence $u_{n+1}=f(u_n)$, on prouve souvent par récurrence un encadrement $m \le u_n \le M$ ou une formule explicite $u_n = \dots$
\end{methode}

\begin{exoresolu}[Récurrence sur une suite définie par récurrence]
Soit la suite $(u_n)$ définie par $u_0 = 3$ et $u_{n+1} = 2u_n + 1$ pour tout $n \in \mathbb{N}$.
Prouver par récurrence que pour tout $n \in \mathbb{N}$, $u_n = 2^{n+2} - 1$.
\tcblower
\textbf{Solution.}
\emph{Propriété $\mathcal{P}(n)$ :} « $u_n = 2^{n+2} - 1$ ».
\begin{enumerate}
    \item \textbf{Initialisation ($n=0$) :} $u_0 = 3$ (donnée). D'autre part $2^{0+2} - 1 = 2^2 - 1 = 4 - 1 = 3$. Donc $u_0 = 2^{2} - 1$. $\mathcal{P}(0)$ est vraie.
    \item \textbf{Hérédité :} Soit $k \in \mathbb{N}$. Supposons $\mathcal{P}(k)$ vraie, c'est-à-dire $u_k = 2^{k+2} - 1$.
    Calculons $u_{k+1}$ à l'aide de la relation de récurrence :
    \begin{align*}
        u_{k+1} &= 2u_k + 1 \\
                &= 2(2^{k+2} - 1) + 1 \quad \text{(par hypothèse de récurrence)} \\
                &= 2^{k+3} - 2 + 1 \\
                &= 2^{k+3} - 1 \\
                &= 2^{(k+1)+2} - 1.
    \end{align*}
    Ainsi $\mathcal{P}(k+1)$ est vraie.
    \item \textbf{Conclusion :} Par le principe de récurrence, pour tout $n \in \mathbb{N}$, $u_n = 2^{n+2} - 1$.
\end{enumerate}
\end{exoresolu}

\begin{methode}[Étudier la monotonie et la bornitude d'une suite]
\textbf{Objectif :} Déterminer le sens de variation (croissante/décroissante) et trouver des bornes (majorant/minorant).
\begin{enumerate}
    \item \textbf{Conjecture (calcul des premiers termes) :} Calculez $u_0, u_1, u_2, u_3$ pour deviner la monotonie et des bornes plausibles.
    \item \textbf{Monotonie :} Étudiez le signe de $u_{n+1} - u_n$ (différence) ou le rapport $\frac{u_{n+1}}{u_n}$ (si termes de signe constant).
        \begin{itemize}
            \item Si $u_{n+1} - u_n \ge 0$ : suite croissante.
            \item Si $u_{n+1} - u_n \le 0$ : suite décroissante.
        \end{itemize}
        \cle{Astuce :} Exprimez $u_{n+1}-u_n$ en fonction de $u_n$ grâce à la relation $u_{n+1}=f(u_n)$ : $u_{n+1}-u_n = f(u_n)-u_n$. Signez cette expression sur l'intervalle où se trouve $u_n$.
    \item \textbf{Bornitude (souvent par récurrence) :}
        \begin{itemize}
            \item Pour majorer : Trouver $M$ tel que $u_n \le M$. Prouver $u_0 \le M$ et $u_n \le M \implies u_{n+1} = f(u_n) \le M$ (vérifier que $f(x) \le M$ sur $[m, M]$).
            \item Pour minorer : Trouver $m$ tel que $m \le u_n$. Même démarche.
        \end{itemize}
\end{enumerate}
\end{methode}

\begin{exoresolu}[Monotonie et bornitude — Suite racine]
Soit la suite $(u_n)$ définie par $u_0 = 1$ et $u_{n+1} = \sqrt{2 + u_n}$ pour tout $n \in \mathbb{N}$.
\begin{enumerate}
    \item Calculer $u_1, u_2, u_3$. Conjecturer la monotonie et un encadrement.
    \item Prouver par récurrence que pour tout $n \in \mathbb{N}$, $1 \le u_n \le 2$.
    \item Prouver que la suite est croissante.
\end{enumerate}
\tcblower
\textbf{Solution.}
\begin{enumerate}
    \item $u_1 = \sqrt{3} \approx 1{,}73$ ; $u_2 = \sqrt{2+\sqrt{3}} \approx 1{,}93$ ; $u_3 \approx 1{,}98$.
    Conjecture : Suite croissante, majorée par $2$, minorée par $1$.
    \item \textbf{Encadrement par récurrence.}
    \emph{Propriété $\mathcal{P}(n)$ :} $1 \le u_n \le 2$.
    \begin{itemize}
        \item \textbf{Initialisation ($n=0$) :} $u_0=1$, donc $1 \le 1 \le 2$. Vrai.
        \item \textbf{Hérédité :} Supposons $1 \le u_k \le 2$.
        Alors $3 \le 2+u_k \le 4 \implies \sqrt{3} \le \sqrt{2+u_k} \le 2$.
        Or $u_{k+1} = \sqrt{2+u_k}$. Donc $1 < \sqrt{3} \le u_{k+1} \le 2$.
        $\mathcal{P}(k+1)$ est vraie.
    \end{itemize}
    Par récurrence, $\forall n \in \mathbb{N},\ 1 \le u_n \le 2$. La suite est \textbf{bornée}.
    \item \textbf{Monotonie.}
    Pour tout $n$, $u_{n+1} - u_n = \sqrt{2+u_n} - u_n$.
    Posons $f(x) = \sqrt{2+x} - x$ sur $[1, 2]$.
    $f'(x) = \frac{1}{2\sqrt{2+x}} - 1$. Sur $[1,2]$, $\sqrt{2+x} \ge \sqrt{3} > 1 \implies \frac{1}{2\sqrt{2+x}} < \frac{1}{2} \implies f'(x) < 0$.
    $f$ est décroissante. $f(2) = \sqrt{4}-2 = 0$. Donc $\forall x \in [1,2], f(x) \ge f(2) = 0$.
    Comme $u_n \in [1,2]$, $u_{n+1}-u_n = f(u_n) \ge 0$.
    La suite $(u_n)$ est \textbf{croissante}.
\end{enumerate}
\end{exoresolu}

\begin{methode}[Démontrer la convergence d'une suite monotone bornée (Théorème de la convergence monotone)]
\textbf{Théorème (Exigible au Bac) :} Toute suite croissante et majorée converge vers sa borne supérieure (son majorant le plus petit). Toute suite décroissante et minorée converge vers sa borne inférieure.
\begin{enumerate}
    \item \textbf{Montrer la monotonie} (croissante ou décroissante) — cf. méthode précédente.
    \item \textbf{Montrer la bornitude} (majorée si croissante, minorée si décroissante) — cf. méthode précédente.
    \item \textbf{Conclure la convergence :} « La suite $(u_n)$ est croissante et majorée (ou décroissante et minorée), donc elle converge d'après le théorème de la convergence monotone. »
    \item \textbf{Trouver la limite $\ell$ :} Passer à la limite dans la relation de récurrence $u_{n+1} = f(u_n)$.
        Comme $f$ est continue sur l'intervalle des valeurs de $u_n$, on a $\ell = f(\ell)$.
        Résoudre l'équation $\ell = f(\ell)$ et sélectionner la solution compatible avec l'encadrement de la suite.
\end{enumerate}
\end{methode}

\begin{exoresolu}[Convergence et limite — Suite racine (suite)]
On reprend la suite $u_0=1$, $u_{n+1}=\sqrt{2+u_n}$.
\begin{enumerate}
    \item Démontrer que la suite converge.
    \item Déterminer sa limite.
\end{enumerate}
\tcblower
\textbf{Solution.}
\begin{enumerate}
    \item D'après l'exercice précédent, $(u_n)$ est croissante et majorée par $2$.
    Par le \textbf{théorème de la convergence monotone}, la suite $(u_n)$ converge. Notons $\ell$ sa limite.
    \item La fonction $f(x) = \sqrt{2+x}$ est continue sur $[1, 2]$ (intervalle contenant tous les $u_n$).
    On passe à la limite dans $u_{n+1} = f(u_n)$ :
    $\ell = f(\ell) = \sqrt{2+\ell}$.
    $\ell \ge 0$ (car $u_n \ge 1$). En élevant au carré : $\ell^2 = 2 + \ell \iff \ell^2 - \ell - 2 = 0 \iff (\ell-2)(\ell+1)=0$.
    Solutions : $\ell = 2$ ou $\ell = -1$ (rejetée car $\ell \ge 1$).
    \textbf{Conclusion :} $\lim_{n\to+\infty} u_n = 2$.
\end{enumerate}
\end{exoresolu}

\begin{methode}[Étudier une suite $u_{n+1} = f(u_n)$ : variation, convergence, limite]
\textbf{Démarche type « Bac » (fonction $f$ dérivable sur un intervalle $I$) :}
\begin{enumerate}
    \item \textbf{Étude de $f$ :} Définir $f$ sur $I$. Calculer $f'(x)$. Dresser le tableau de variations de $f$ sur $I$. Trouver les points fixes (solutions de $f(x)=x$ dans $I$).
    \item \textbf{Invariance de $I$ (Récurrence) :} Prouver que $\forall n,\ u_n \in I$.
        Souvent : $u_0 \in I$ et $f(I) \subset I$ (i.e. $\forall x \in I, f(x) \in I$).
    \item \textbf{Monotonie :} Étudier le signe de $f(x)-x$ sur $I$.
        \begin{itemize}
            \item Si $f(x)-x \ge 0$ sur $I$ : suite croissante.
            \item Si $f(x)-x \le 0$ sur $I$ : suite décroissante.
        \end{itemize}
    \item \textbf{Convergence :} Suite monotone + bornée (car dans $I$) $\implies$ Convergence (Théorème convergence monotone).
    \item \textbf{Limite $\ell$ :} $\ell$ est un point fixe de $f$ dans $I$ ($\ell = f(\ell)$). Choisir celui compatible avec la monotonie et $u_0$.
\end{enumerate}
\cle{Point clé :} Le signe de $f(x)-x$ détermine la monotonie. Le signe de $f'(x)$ au voisinage du point fixe détermine la vitesse (attractif si $|f'(\ell)|<1$).
\end{methode}

\begin{exoresolu}[Étude complète $u_{n+1}=f(u_n)$]
Soit $f(x) = \frac{1}{2-x}$ définie sur $I = [0, 1]$. Suite $(u_n)$ : $u_0 = 0$, $u_{n+1} = f(u_n)$.
\begin{enumerate}
    \item Étudier les variations de $f$ sur $I$. Résoudre $f(x)=x$ sur $I$.
    \item Montrer que $\forall n \in \mathbb{N},\ u_n \in [0, 1]$.
    \item Montrer que $(u_n)$ est croissante. En déduire sa convergence.
    \item Calculer la limite $\ell$.
\end{enumerate}
\tcblower
\textbf{Solution.}
\begin{enumerate}
    \item $f'(x) = \frac{1}{(2-x)^2} > 0$ sur $[0,1]$. $f$ est \textbf{strictement croissante}.
    $f(0)=\frac{1}{2}$, $f(1)=1$. Donc $f([0,1]) = [\frac{1}{2}, 1] \subset [0,1]$.
    $f(x)=x \iff \frac{1}{2-x}=x \iff 1 = 2x - x^2 \iff x^2-2x+1=0 \iff (x-1)^2=0$.
    Unique point fixe : $x=1$.
    \item \textbf{Récurrence :} $\mathcal{P}(n) : u_n \in [0,1]$.
    $u_0=0 \in [0,1]$. Supposons $u_k \in [0,1]$. Alors $u_{k+1}=f(u_k) \in f([0,1]) = [\frac{1}{2}, 1] \subset [0,1]$. Prouvé.
    \item \textbf{Monotonie :} Pour $x \in [0,1]$, $f(x)-x = \frac{1}{2-x}-x = \frac{1 - x(2-x)}{2-x} = \frac{(x-1)^2}{2-x} \ge 0$.
    Donc $u_{n+1}-u_n = f(u_n)-u_n \ge 0$. Suite \textbf{croissante}.
    Suite croissante majorée par $1$ $\implies$ \textbf{Convergente}.
    \item \textbf{Limite :} $\ell = f(\ell) \implies \ell = 1$ (unique point fixe dans $[0,1]$).
    $\lim u_n = 1$.
\end{enumerate}
\end{exoresolu}

\begin{popfigure}
\begin{tikzpicture}[scale=1.2, >=Stealth]
    \draw[->] (-0.2,0) -- (2.2,0) node[right] {$x$};
    \draw[->] (0,-0.2) -- (0,2.2) node[above] {$y$};
    \draw[popblue, thick, domain=0:1, samples=100] plot (\x, {1/(2-\x)}) node[right, pos=0.9] {$y=f(x)$};
    \draw[poporange, thick, domain=0:1] plot (\x, {\x}) node[right] {$y=x$};
    \draw[dashed, popmuted] (0,1) node[left] {$1$} -- (1,1) -- (1,0) node[below] {$1$};
    \draw[dashed, popmuted] (0,0.5) node[left] {$0{,}5$} -- (0.5,0.5) -- (0.5,0) node[below] {$0{,}5$};
    % Escalier (toiles d'araignée)
    \draw[popgreen, thick] (0,0) -- (0,0.5);
    \draw[popgreen, thick] (0,0.5) -- (0.5,0.5);
    \draw[popgreen, thick] (0.5,0.5) -- (0.5, 0.666);
    \draw[popgreen, thick] (0.5, 0.666) -- (0.666, 0.666);
    \draw[popgreen, thick] (0.666, 0.666) -- (0.666, 0.75);
    \node[popgreen] at (1.5, 0.3) {Construction de $(u_n)$};
\end{tikzpicture}
\legende{Représentation graphique (toiles d'araignée) de la suite $u_{n+1}=\frac{1}{2-u_n}$. L'escalier vert monte vers le point fixe $(1,1)$.}
\end{popfigure}

\begin{methode}[Encadrer la limite par l'inégalité des accroissements finis (IAF)]
\textbf{Contexte :} Suite convergente $u_{n+1}=f(u_n)$, limite $\ell$, $f$ de classe $\mathcal{C}^1$ sur un intervalle $I$ contenant $\ell$ et tous les $u_n$.
\textbf{Théorème (IAF) :} $\forall n,\ \exists c_n$ entre $u_n$ et $\ell$ tel que $|u_{n+1}-\ell| = |f'(c_n)|\cdot|u_n-\ell|$.
\begin{enumerate}
    \item \textbf{Majoration de $|f'|$ :} Trouver $k \in [0, 1[$ tel que $\forall x \in I,\ |f'(x)| \le k$. (Souvent $|f'(\ell)| < 1$ et par continuité de $f'$, on trouve un voisinage où $|f'| \le k < 1$).
    \item \textbf{Inégalité de contraction :} $|u_{n+1}-\ell| \le k |u_n-\ell|$.
    \item \textbf{Encadrement géométrique :} Par récurrence, $|u_n-\ell| \le k^n |u_0-\ell|$.
    \item \textbf{Précision $\varepsilon$ :} Pour garantir $|u_n-\ell| < \varepsilon$, il suffit que $k^n |u_0-\ell| < \varepsilon \iff n > \frac{\ln(\varepsilon/|u_0-\ell|)}{\ln k}$.
    \item \textbf{Encadrement numérique :} $\ell \in [u_n - k^n|u_0-\ell|,\ u_n + k^n|u_0-\ell|]$.
\end{enumerate}
\cle{Attention :} L'IAF donne un encadrement \emph{théorique}. En pratique, on utilise souvent $|u_{n+1}-u_n| \approx (1-k)|u_n-\ell|$ pour estimer l'erreur résiduelle.
\end{methode}

\begin{exoresolu}[Encadrement par IAF — Méthode de Héron pour $\sqrt{5}$]
Soit $u_0 = 2$ et $u_{n+1} = \frac{1}{2}\left(u_n + \frac{5}{u_n}\right)$.
On admet que $(u_n)$ converge vers $\ell = \sqrt{5} \approx 2{,}236$.
\begin{enumerate}
    \item Montrer que $f(x)=\frac{1}{2}(x+\frac{5}{x})$ satisfait $|f'(x)| \le \frac{1}{2}$ sur $[2, 3]$.
    \item En déduire un encadrement de $|\sqrt{5} - u_n|$ en fonction de $n$.
    \item Déterminer le plus petit $n$ tel que $|u_n - \sqrt{5}| < 10^{-4}$.
\end{enumerate}
\tcblower
\textbf{Solution.}
\begin{enumerate}
    \item $f'(x) = \frac{1}{2}\left(1 - \frac{5}{x^2}\right) = \frac{x^2-5}{2x^2}$.
    Sur $[2, 3]$, $x^2 \in [4, 9]$, donc $x^2-5 \in [-1, 4]$.
    $|f'(x)| = \frac{|x^2-5|}{2x^2}$.
    Le numérateur est $\le 4$, le dénominateur $\ge 2\times 4 = 8$.
    Donc $|f'(x)| \le \frac{4}{8} = \frac{1}{2}$. On peut prendre $k = \frac{1}{2}$.
    \item Par IAF, $\exists c_n$ entre $u_n$ et $\ell$ : $|u_{n+1}-\ell| = |f'(c_n)||u_n-\ell| \le \frac{1}{2}|u_n-\ell|$.
    Par récurrence : $|u_n-\ell| \le \left(\frac{1}{2}\right)^n |u_0-\ell|$.
    $|u_0-\ell| = |2-\sqrt{5}| = \sqrt{5}-2 \approx 0{,}236 < 0{,}24$.
    Donc $|u_n - \sqrt{5}| < 0{,}24 \times 2^{-n}$.
    \item On veut $0{,}24 \times 2^{-n} < 10^{-4} \iff 2^n > 0{,}24 \times 10^4 = 2400$.
    $2^{11} = 2048 < 2400$ ; $2^{12} = 4096 > 2400$.
    Le plus petit $n$ est \textbf{$n=12$}.
    (Vérification : $u_{12}$ donne $\sqrt{5}$ à $10^{-4}$ près).
\end{enumerate}
\end{exoresolu}

\begin{methode}[Synthèse : Méthode complète d'étude d'une suite au Bac]
\textbf{Check-list « Zéro faute » pour un problème type :}
\begin{enumerate}
    \item \textbf{Lire l'énoncé :} Identifier $u_0$, $f$, l'intervalle $I$ suggéré.
    \item \textbf{Étudier $f$ sur $I$ :} Dérivée, variations, image $f(I)$, points fixes ($f(x)=x$).
    \item \textbf{Invariance (Récurrence) :} Prouver $\forall n,\ u_n \in I$ ($u_0 \in I$ et $f(I) \subset I$).
    \item \textbf{Monotonie :} Signe de $f(x)-x$ sur $I$ $\implies$ Croissante / Décroissante.
    \item \textbf{Bornitude :} Conséquence immédiate de $u_n \in I$ (majorée par $\max I$, minorée par $\min I$).
    \item \textbf{Convergence :} Invoquer \textbf{explicitement} le théorème de la convergence monotone.
    \item \textbf{Limite $\ell$ :} Résoudre $\ell = f(\ell)$ dans $I$. Justifier le choix par la monotonie/encadrement.
    \item \textbf{Approximation (si demandé) :} Utiliser l'IAF pour trouver $n$ tel que $|u_n-\ell| < \varepsilon$.
\end{enumerate}
\cle{Rédaction :} Une phrase par étape. « On pose $f(x)=...$ », « $f$ est continue sur $I$ », « Par récurrence, $\forall n, u_n \in I$ », « Donc la suite converge », « Sa limite $\ell$ vérifie... ».
\end{methode}

\begin{exoresolu}[Sujet type Bac — Synthèse]
Soit $f(x) = \frac{x^2+2}{2x+1}$ définie sur $I = [1, 2]$. Suite $(u_n)$ : $u_0 = 2$, $u_{n+1} = f(u_n)$.
\begin{enumerate}
    \item Étudier les variations de $f$ sur $I$. Montrer que $f(I) \subset I$.
    \item Montrer que $\forall n,\ u_n \in [1, 2]$. En déduire que $(u_n)$ est bien définie.
    \item Étudier la monotonie de $(u_n)$.
    \item Démontrer que $(u_n)$ converge et calculer sa limite $\ell$.
    \item On admet $|f'(x)| \le \frac{1}{3}$ sur $[1, 2]$. Déterminer un entier $n$ tel que $|u_n - \ell| < 10^{-3}$.
\end{enumerate}
\tcblower
\textbf{Solution détaillée.}
\begin{enumerate}
    \item $f'(x) = \frac{2x(2x+1) - 2(x^2+2)}{(2x+1)^2} = \frac{4x^2+2x-2x^2-4}{(2x+1)^2} = \frac{2x^2+2x-4}{(2x+1)^2} = \frac{2(x^2+x-2)}{(2x+1)^2} = \frac{2(x-1)(x+2)}{(2x+1)^2}$.
    Sur $[1, 2]$, $x-1 \ge 0$, $x+2 > 0$, dénominateur $> 0$. Donc $f'(x) \ge 0$.
    $f$ est \textbf{croissante} sur $[1, 2]$.
    $f(1) = \frac{3}{3} = 1$ ; $f(2) = \frac{6}{5} = 1{,}2$.
    Donc $f([1, 2]) = [1, 1{,}2] \subset [1, 2]$.
    \item \textbf{Récurrence :} $\mathcal{P}(n) : u_n \in [1, 2]$.
    $u_0 = 2 \in [1, 2]$. Supposons $u_k \in [1, 2]$. Alors $u_{k+1} = f(u_k) \in f([1, 2]) = [1, 1{,}2] \subset [1, 2]$.
    $\mathcal{P}(k+1)$ vraie. Par récurrence, $\forall n,\ u_n \in [1, 2]$. Suite bien définie (dénominateur $2u_n+1 \ge 3 \neq 0$).
    \item \textbf{Monotonie :} $f(x)-x = \frac{x^2+2}{2x+1} - x = \frac{x^2+2 - 2x^2 - x}{2x+1} = \frac{-x^2 - x + 2}{2x+1} = \frac{-(x-1)(x+2)}{2x+1}$.
    Sur $[1, 2]$, $x-1 \ge 0$, $x+2 > 0$, dénominateur $> 0$. Donc $f(x)-x \le 0$.
    $u_{n+1}-u_n = f(u_n)-u_n \le 0$. Suite \textbf{décroissante}.
    \item Suite décroissante minorée par $1$ $\implies$ \textbf{Convergente} (Théorème convergence monotone).
    Limite $\ell \in [1, 2]$. $\ell = f(\ell) \iff \ell = \frac{\ell^2+2}{2\ell+1} \iff \ell(2\ell+1) = \ell^2+2 \iff 2\ell^2+\ell = \ell^2+2 \iff \ell^2+\ell-2=0 \iff (\ell-1)(\ell+2)=0$.
    $\ell = 1$ ou $\ell = -2$ (hors $[1,2]$). Donc \textbf{$\ell = 1$}.
    \item \textbf{IAF :} $|u_{n+1}-1| \le \frac{1}{3}|u_n-1| \implies |u_n-1| \le \left(\frac{1}{3}\right)^n |u_0-1| = \left(\frac{1}{3}\right)^n$.
    On veut $\left(\frac{1}{3}\right)^n < 10^{-3} \iff 3^n > 1000$.
    $3^6 = 729 < 1000$ ; $3^7 = 2187 > 1000$.
    Le plus petit $n$ est \textbf{$n=7$}.
\end{enumerate}
\end{exoresolu}

\begin{attention}[Les 5 erreurs qui coûtent des points au Bac]
\begin{enumerate}
    \item \textbf{Oublier l'initialisation ou la conclusion en récurrence.} « Par récurrence, c'est vrai » sans avoir écrit $\mathcal{P}(0)$ vraie et l'implication $\mathcal{P}(k) \implies \mathcal{P}(k+1)$ explicitement. \cle{Sanction :} 0 point à la question récurrence.
    \item \textbf{Utiliser le théorème de convergence monotone sans vérifier les DEUX conditions.} Dire « la suite converge car elle est croissante » (oubli : majorée) ou « car elle est majorée » (oubli : monotone). \cle{Sanction :} Perte de 2 à 3 points sur la convergence.
    \item \textbf{Passer à la limite sans justifier la continuité de $f$.}$\ell = f(\ell)$ n'est valable que si $f$ est continue en $\ell$ (ou au moins sur l'adhérence des valeurs de $u_n$). Écrire : « $f$ est continue sur $I$ donc on peut passer à la limite. »
    \item \textbf{Mauvaise gestion du signe pour la monotonie.} Étudier le signe de $u_{n+1}-u_n$ sans le relier à $f(x)-x$, ou signer $f(x)-x$ sur $\mathbb{R}$ au lieu de l'intervalle $I$ où vit la suite. Ex: $f(x)-x \ge 0$ sur $\mathbb{R}$ mais $u_n$ sort de l'intervalle où c'est vrai.
    \item \textbf{Confondre $|f'(x)| \le k$ et $|f'(\ell)| < 1$ pour l'IAF.} L'IAF demande une majoration \textbf{uniforme} $|f'(x)| \le k < 1$ sur tout l'intervalle $I$ (ou un voisinage stable), pas seulement en $\ell$. Dire « $f'(\ell)=0,3 < 1$ donc $|u_{n+1}-\ell| \le 0,3 |u_n-\ell|$ » est \textbf{faux} sans justification de la majoration uniforme.
\end{enumerate}
\end{attention}

\newpage

\coursec{Exercices}

\coursec{Suites numériques}

\courssub{Raisonnement par récurrence}

\begin{definition}[Principe de récurrence]
Soit $\mathcal{P}(n)$ une propriété dépendant d'un entier naturel $n$. Si :
\begin{enumerate}
\item \textbf{Initialisation :} $\mathcal{P}(n_0)$ est vraie pour un certain entier $n_0$ ;
\item \textbf{Hérédité :} pour tout entier $n \geq n_0$, $\mathcal{P}(n)$ vraie $\implies$ $\mathcal{P}(n+1)$ vraie ;
\end{enumerate}
alors $\mathcal{P}(n)$ est vraie pour tout entier $n \geq n_0$.
\end{definition}

\begin{methode}[Structure d'une démonstration par récurrence]
\begin{enumerate}
\item \textbf{Annoncer} la propriété $\mathcal{P}(n)$ et l'entier initial $n_0$.
\item \textbf{Initialisation :} Vérifier $\mathcal{P}(n_0)$ par un calcul direct.
\item \textbf{Hérédité :} Supposer $\mathcal{P}(n)$ vraie pour un $n \geq n_0$ arbitraire (hypothèse de récurrence), et en déduire $\mathcal{P}(n+1)$ par un raisonnement algébrique ou logique rigoureux.
\item \textbf{Conclusion :} Par le principe de récurrence, $\mathcal{P}(n)$ est vraie pour tout $n \geq n_0$.
\end{enumerate}
\end{methode}

\begin{attention}[Pièges fréquents]
\begin{itemize}
\item Ne pas oublier l'initialisation (même si elle semble évidente).
\item Ne pas confondre « supposer $\mathcal{P}(n)$ vraie » et « supposer $\mathcal{P}(n)$ vraie \textbf{pour tout $n$} » (ce qui serait ce qu'on veut prouver).
\item L'hérédité doit être montrée \textbf{pour tout $n \geq n_0$}, pas seulement pour quelques valeurs.
\end{itemize}
\end{attention}

\begin{exemplebox}[Somme des $n$ premiers carrés]
On veut prouver $\forall n \geq 1,\; \sum_{k=1}^n k^2 = \frac{n(n+1)(2n+1)}{6}$.
\begin{itemize}
\item \textbf{Initialisation ($n=1$) :} $1^2 = 1$ et $\frac{1\times2\times3}{6}=1$. OK.
\item \textbf{Hérédité :} On suppose $\sum_{k=1}^n k^2 = \frac{n(n+1)(2n+1)}{6}$.
Alors $\sum_{k=1}^{n+1} k^2 = \frac{n(n+1)(2n+1)}{6} + (n+1)^2 = (n+1)\left[\frac{n(2n+1)}{6} + (n+1)\right]$
$= (n+1)\frac{2n^2+n+6n+6}{6} = \frac{(n+1)(2n^2+7n+6)}{6} = \frac{(n+1)(n+2)(2n+3)}{6}$.
C'est bien la formule pour $n+1$.
\end{itemize}
\end{exemplebox}

\begin{exercice}[QCM : récurrence et monotonie]
Pour chaque question, une seule réponse est exacte. Indique-la en justifiant brièvement.
\begin{enumerate}
\item Pour démontrer par récurrence qu'une propriété $\mathcal{P}(n)$ est vraie pour tout entier $n \geq n_0$, il faut :
\begin{enumerate}
\item vérifier $\mathcal{P}(n_0)$ et montrer que $\mathcal{P}(n)$ entraîne $\mathcal{P}(n+1)$ pour tout $n \geq n_0$ ;
\item vérifier $\mathcal{P}(n_0)$ et $\mathcal{P}(n_0 + 1)$ ;
\item montrer que $\mathcal{P}(n+1)$ est vraie pour tout $n$.
\end{enumerate}
\item La suite $(u_n)$ définie par $u_n = 3 + \dfrac{1}{n+1}$ est :
\begin{enumerate}
\item croissante et majorée par $4$ ;
\item décroissante et minorée par $3$ ;
\item ni monotone ni bornée.
\end{enumerate}
\item Une suite croissante et majorée par $5$ :
\begin{enumerate}
\item converge vers $5$ ;
\item converge vers une limite $\ell \leq 5$ ;
\item peut diverger.
\end{enumerate}
\item Soit $(u_n)$ définie par $u_0 = 1$ et $u_{n+1} = f(u_n)$ où $f(x) = \dfrac{x+4}{3}$. Si $(u_n)$ converge vers $\ell$, alors $\ell$ vérifie :
\begin{enumerate}
\item $\ell = \dfrac{1+4}{3}$ ;
\item $\ell = \dfrac{\ell + 4}{3}$ ;
\item $\ell = f(0)$.
\end{enumerate}
\end{enumerate}
\end{exercice}

\begin{exercice}[Démonstrations par récurrence]
\begin{enumerate}
\item Démontre par récurrence que pour tout entier $n \geq 1$ :
\[ 1^2 + 2^2 + \cdots + n^2 = \frac{n(n+1)(2n+1)}{6}. \]
\item Démontre par récurrence que pour tout entier $n \geq 1$, l'entier $3^{2n} - 2^n$ est divisible par $7$.
\item Démontre par récurrence que pour tout entier $n \geq 4$, $2^n \geq n^2$.
\end{enumerate}
\end{exercice}

\courssub{Suites majorées, minorées, bornées}

\begin{definition}[Majorante, minorante, bornée]
Soit $(u_n)$ une suite réelle.
\begin{itemize}
\item $(u_n)$ est \textbf{majorée} s'il existe $M \in \mathbb{R}$ tel que $\forall n,\; u_n \leq M$. $M$ est une majorante.
\item $(u_n)$ est \textbf{minorée} s'il existe $m \in \mathbb{R}$ tel que $\forall n,\; u_n \geq m$. $m$ est une minorante.
\item $(u_n)$ est \textbf{bornée} si elle est à la fois majorée et minorée, c'est-à-dire $\exists M>0,\; \forall n,\; |u_n| \leq M$.
\end{itemize}
\end{definition}

\begin{propriete}[Suite convergente $\implies$ bornée]
Toute suite convergente est bornée. La réciproque est fausse (ex: $u_n = (-1)^n$).
\end{propriete}

\begin{exemplebox}[Étude de la bornitude]
\begin{enumerate}
\item $u_n = \dfrac{2n+1}{n+1} = 2 - \dfrac{1}{n+1}$. Pour $n \geq 0$, $1 \leq u_n < 2$. Suite \textbf{bornée} (majorée par 2, minorée par 1).
\item $v_n = (-1)^n n$. $|v_n| = n \to +\infty$. Suite \textbf{non bornée} (ni majorée, ni minorée).
\item $w_n = 3 - n^2$. $w_n \leq 3$ (majorée par 3) mais $w_n \to -\infty$ (non minorée). Suite \textbf{non bornée}.
\end{enumerate}
\end{exemplebox}

\begin{exercice}[Vrai ou faux : justifier]
Réponds par Vrai ou Faux à chacune des affirmations suivantes, en justifiant par une démonstration ou un contre-exemple.
\begin{enumerate}
\item Toute suite croissante converge.
\item Toute suite décroissante et minorée par $0$ converge.
\item Toute suite bornée est monotone.
\item Si $u_{n+1} = f(u_n)$ avec $f$ continue et si $f(\ell) = \ell$, alors $(u_n)$ converge vers $\ell$.
\end{enumerate}
\end{exercice}

\begin{exercice}[Définitions et bornitude]
\begin{enumerate}
\item Donne la définition d'une suite majorée, d'une suite minorée, d'une suite bornée.
\item Pour chacune des suites suivantes, précise si elle est majorée, minorée, bornée :
\begin{enumerate}
\item $u_n = \dfrac{2n+1}{n+1}$ pour $n \in \mathbb{N}$ ;
\item $v_n = (-1)^n\, n$ pour $n \in \mathbb{N}$ ;
\item $w_n = 3 - n^2$ pour $n \in \mathbb{N}$.
\end{enumerate}
\end{enumerate}
\end{exercice}

\courssub{Suites monotones et théorème de convergence monotone}

\begin{definition}[Monotonie]
Soit $(u_n)$ une suite réelle.
\begin{itemize}
\item $(u_n)$ est \textbf{croissante} si $\forall n,\; u_{n+1} \geq u_n$ (équivalent à $u_{n+1}-u_n \geq 0$ ou $\frac{u_{n+1}}{u_n} \geq 1$ si $u_n>0$).
\item $(u_n)$ est \textbf{décroissante} si $\forall n,\; u_{n+1} \leq u_n$.
\item $(u_n)$ est \textbf{monotone} si elle est croissante ou décroissante.
\item On dit \textbf{strictement} croissante/décroissante si les inégalités sont strictes.
\end{itemize}
\end{definition}

\begin{propriete}[Théorème de convergence monotone -- Exigible au Bac]
\begin{itemize}
\item Toute suite \textbf{croissante et majorée} converge vers sa borne supérieure (limite $\ell \leq$ toute majorante).
\item Toute suite \textbf{décroissante et minorée} converge vers sa borne inférieure (limite $\ell \geq$ toute minorante).
\end{itemize}
\end{propriete}

\begin{attention}[Conditions d'application]
Il faut \textbf{les deux} conditions : monotonie \textbf{ET} bornitude (dans le bon sens). Une suite croissante non majorée diverge vers $+\infty$. Une suite décroissante non minorée diverge vers $-\infty$.
\end{attention}

\begin{exemplebox}[Application du théorème]
La suite $u_n = 3 + \frac{1}{n+1}$ est décroissante (car $n \mapsto n+1$ croissante $\implies \frac{1}{n+1}$ décroissante) et minorée par $3$. Elle converge vers $3$.
\end{exemplebox}

\begin{exercice}[Calculs directs]
Soit $(u_n)$ la suite définie par $u_0 = 1$ et, pour tout $n \in \mathbb{N}$, $u_{n+1} = \dfrac{u_n + 4}{3}$.
\begin{enumerate}
\item Calcule $u_1$, $u_2$ et $u_3$.
\item Conjecture le sens de variation de $(u_n)$ et sa limite éventuelle.
\item Soit $(a_n)$ définie par $a_n = \dfrac{n^2}{n+1}$. Étudie le sens de variation de $(a_n)$ en calculant $a_{n+1} - a_n$.
\end{enumerate}
\end{exercice}

\begin{exercice}[Monotonie et bornitude]
Soit $(u_n)$ la suite définie par $u_0 = 1$ et, pour tout $n \in \mathbb{N}$, $u_{n+1} = \dfrac{u_n + 4}{3}$.
\begin{enumerate}
\item Démontre par récurrence que pour tout $n \in \mathbb{N}$ : $1 \leq u_n \leq 2$.
\item Étudie le signe de $u_{n+1} - u_n$ et déduis-en le sens de variation de $(u_n)$.
\item Justifie que $(u_n)$ converge et détermine sa limite.
\end{enumerate}
\end{exercice}

\courssub{Suites récurrentes $u_{n+1} = f(u_n)$ : étude de la convergence}

\begin{methode}[Étude complète d'une suite $u_{n+1}=f(u_n)$]
\begin{enumerate}
\item \textbf{Domaine :} Déterminer un intervalle $I$ tel que $u_0 \in I$ et $f(I) \subset I$ (stabilité).
\item \textbf{Variations de $f$ :} Étudier $f$ sur $I$ (dérivée, signe, tableau de variations).
\item \textbf{Point(s) fixe(s) :} Résoudre $f(x)=x$ dans $I$. Les solutions sont les limites candidates $\ell$.
\item \textbf{Monotonie de $(u_n)$ :} Étudier le signe de $u_{n+1}-u_n = f(u_n)-u_n$ sur $I$ (souvent via le signe de $f(x)-x$).
\item \textbf{Convergence :} Appliquer le théorème de convergence monotone (si monotone et bornée par $I$).
\item \textbf{Limite :} La limite $\ell$ est le point fixe dans $I$ compatible avec la monotonie et $u_0$.
\end{enumerate}
\end{methode}

\begin{propriete}[Suite auxiliaire pour fraction rationnelle homographique]
Si $u_{n+1} = \frac{a u_n + b}{c u_n + d}$ avec $ad-bc \neq 0$, on cherche les points fixes $\alpha, \beta$ (racines de $cx^2+(d-a)x-b=0$).
Si $\alpha \neq \beta$, la suite $v_n = \frac{u_n-\alpha}{u_n-\beta}$ est \textbf{géométrique}.
Si $\alpha = \beta$ (double racine), la suite $v_n = \frac{1}{u_n-\alpha}$ est \textbf{arithmétique}.
\end{propriete}

\begin{exemplebox}[Suite auxiliaire géométrique]
Soit $u_0=0$, $u_{n+1} = \frac{3u_n+2}{u_n+4}$.
Points fixes : $x = \frac{3x+2}{x+4} \iff x^2+x-2=0 \iff x=1$ ou $x=-2$.
Posons $v_n = \frac{u_n-1}{u_n+2}$.
$v_{n+1} = \frac{u_{n+1}-1}{u_{n+1}+2} = \frac{\frac{3u_n+2}{u_n+4}-1}{\frac{3u_n+2}{u_n+4}+2} = \frac{2u_n-2}{5u_n+10} = \frac{2}{5}v_n$.
$(v_n)$ est géométrique de raison $\frac{2}{5}$, $v_0 = -\frac{1}{2}$.
$v_n = -\frac{1}{2}\left(\frac{2}{5}\right)^n \to 0$, donc $u_n = \frac{1+2v_n}{1-v_n} \to 1$.
\end{exemplebox}

\begin{exercice}[Suite auxiliaire géométrique]
Soit $(u_n)$ la suite définie par $u_0 = 0$ et, pour tout $n \in \mathbb{N}$, $u_{n+1} = \dfrac{3u_n + 2}{u_n + 4}$.
\begin{enumerate}
\item Calcule $u_1$ et $u_2$.
\item Démontre par récurrence que pour tout $n \in \mathbb{N}$, $u_n \geq 0$.
\item On pose, pour tout $n \in \mathbb{N}$, $v_n = \dfrac{u_n - 1}{u_n + 2}$.
\begin{enumerate}
\item Démontre que $(v_n)$ est une suite géométrique dont on précisera la raison et le premier terme.
\item Exprime $v_n$ puis $u_n$ en fonction de $n$.
\item Détermine la limite de $(u_n)$.
\end{enumerate}
\end{enumerate}
\end{exercice}

\courssub{Valeur approchée de la limite par l'inégalité des accroissements finis}

\begin{propriete}[Inégalité des accroissements finis (IAF) -- Exigible au Bac]
Soit $f$ dérivable sur un intervalle $I$. Pour tous $a,b \in I$, il existe $c$ entre $a$ et $b$ tel que :
\[ f(b) - f(a) = f'(c)\,(b-a). \]
Conséquence : si $|f'(x)| \leq k$ sur $I$, alors $|f(b)-f(a)| \leq k\,|b-a|$.
\end{propriete}

\begin{propriete}[Encadrement de la limite]
Si $u_{n+1}=f(u_n)$ avec $f$ dérivable sur $I$ stable, $\ell \in I$ point fixe, et $|f'(x)| \leq k < 1$ sur $I$, alors :
\[ |u_n - \ell| \leq k^n |u_0 - \ell|. \]
Cela donne une vitesse de convergence géométrique et permet de trouver $N$ tel que $|u_N-\ell| \leq \varepsilon$.
\end{propriete}

\begin{exercice}[Encadrement par l'inégalité des accroissements finis]
Soit $f$ la fonction définie sur $[0\,;\,3]$ par $f(x) = \sqrt{2x+3}$.
\begin{enumerate}
\item Étudie les variations de $f$ sur $[0\,;\,3]$ et montre que l'intervalle $[0\,;\,3]$ est stable par $f$.
\item Résous dans $[0\,;\,3]$ l'équation $f(x) = x$.
\item Démontre que pour tout $x \in [0\,;\,3]$, $\dfrac{1}{3} \leq f'(x) \leq \dfrac{1}{\sqrt{3}}$.
\item On définit la suite $(u_n)$ par $u_0 = 0$ et $u_{n+1} = f(u_n)$. À l'aide de l'inégalité des accroissements finis, démontre que pour tout $n \in \mathbb{N}$ :
\[ |u_{n+1} - 3| \leq \frac{1}{\sqrt{3}}\,|u_n - 3|, \qquad \text{puis} \qquad |u_n - 3| \leq 3 \times \left(\frac{1}{\sqrt{3}}\right)^n. \]
\item Détermine un entier $N$ tel que pour tout $n \geq N$, $|u_n - 3| \leq 10^{-2}$.
\end{enumerate}
\end{exercice}

\courssub{Synthèse : méthode complète d'étude d'une suite au Bac}

\begin{methode}[Check-list Bac pour une suite récurrente]
\begin{enumerate}
\item \textbf{Modélisation / Définition :} Écrire $u_{n+1}=f(u_n)$, domaine de $f$, $u_0$.
\item \textbf{Stabilité :} Trouver $I$ tel que $f(I)\subset I$ et $u_0\in I$. Prouver par récurrence $\forall n,\; u_n \in I$.
\item \textbf{Monotonie :} Étudier signe de $u_{n+1}-u_n = f(u_n)-u_n$ sur $I$ (souvent via $f(x)-x$).
\item \textbf{Convergence :} Appliquer le théorème de convergence monotone (croissante+majorée ou décroissante+minorée).
\item \textbf{Limite $\ell$ :} Résoudre $f(\ell)=\ell$ dans $I$. Choisir la valeur compatible avec la monotonie.
\item \textbf{Encadrement (si demandé) :} Majorer $|f'(x)|$ par $k<1$ sur $I$, utiliser IAF : $|u_n-\ell| \leq k^n |u_0-\ell|$. Résoudre $k^N |u_0-\ell| \leq \varepsilon$.
\item \textbf{Suite auxiliaire (si homographique) :} Poser $v_n = \frac{u_n-\alpha}{u_n-\beta}$ (ou $\frac{1}{u_n-\alpha}$), montrer géométrique/arithmétique, expliciter $u_n$.
\end{enumerate}
\end{methode}

\begin{exercice}[Type Bac : pisciculture à Maroua]
Un bassin de pisciculture de la région de l'Extrême-Nord, à Maroua, contient initialement \num{5000} poissons. Chaque mois, la population augmente de $10\,\%$ par reproduction, puis le pisciculteur prélève $400$ poissons pour la vente au marché. On note $p_n$ le nombre de poissons au début du mois $n$ (avec $p_0 = \num{5000}$), après le prélèvement du mois précédent.

\textbf{Partie A : modélisation}
\begin{enumerate}
\item Justifie que pour tout $n \in \mathbb{N}$ : $p_{n+1} = 1{,}1\, p_n - 400$.
\item Calcule $p_1$ et $p_2$.
\end{enumerate}

\textbf{Partie B : étude de la suite}
\begin{enumerate}
\item On pose, pour tout $n \in \mathbb{N}$, $q_n = p_n - \num{4000}$.
\begin{enumerate}
\item Démontre que $(q_n)$ est une suite géométrique de raison $1{,}1$ et précise son premier terme.
\item Exprime $q_n$ puis $p_n$ en fonction de $n$.
\end{enumerate}
\item Étudie le sens de variation de $(p_n)$ et détermine sa limite. Interprète le résultat pour le pisciculteur.
\item À l'aide de la calculatrice, détermine le premier mois où la population dépasse \num{8000} poissons.
\end{enumerate}

\textbf{Partie C : durabilité de l'exploitation}
\begin{enumerate}
\item Le pisciculteur souhaite une population \emph{stable} : il cherche une population $S$ telle que si $p_n = S$, alors $p_{n+1} = S$. Détermine $S$.
\item Que se passe-t-il si la population initiale est de \num{3000} poissons au lieu de \num{5000} ? Justifie en étudiant la nouvelle suite $(q_n)$ associée et commente la viabilité de l'exploitation.
\end{enumerate}
\end{exercice}

\begin{exercice}[Type Bac : étude complète d'une suite récurrente]
On considère la fonction $f$ définie sur $[0\,;\,+\infty[$ par $f(x) = \sqrt{2x+3}$ et la suite $(u_n)$ définie par $u_0 = 0$ et, pour tout $n \in \mathbb{N}$, $u_{n+1} = f(u_n)$.

\textbf{Partie A : étude de la fonction $f$}
\begin{enumerate}
\item Étudie les variations de $f$ sur $[0\,;\,+\infty[$.
\item Démontre que l'intervalle $[0\,;\,3]$ est stable par $f$, c'est-à-dire que pour tout $x \in [0\,;\,3]$, $f(x) \in [0\,;\,3]$.
\item Démontre que l'équation $f(x) = x$ admet une unique solution dans $[0\,;\,3]$, égale à $3$.
\end{enumerate}

\textbf{Partie B : convergence de la suite}
\begin{enumerate}
\item Démontre par récurrence que pour tout $n \in \mathbb{N}$, $0 \leq u_n \leq 3$.
\item Démontre que pour tout $n \in \mathbb{N}$, $u_{n+1} - u_n$ a le même signe que $-u_n^2 + 2u_n + 3$, puis que $(u_n)$ est croissante.
\item Justifie que $(u_n)$ converge et détermine sa limite $\ell$.
\end{enumerate}

\textbf{Partie C : vitesse de convergence}
\begin{enumerate}
\item Démontre que pour tout $x \in [0\,;\,3]$, $\dfrac{1}{3} \leq f'(x) \leq \dfrac{1}{\sqrt{3}}$.
\item En appliquant l'inégalité des accroissements finis entre $u_n$ et $3$, démontre que pour tout $n \in \mathbb{N}$ :
\[ |u_n - 3| \leq 3 \times \left(\frac{1}{\sqrt{3}}\right)^n. \]
\item Détermine le plus petit entier $N$ tel que $u_N$ soit une valeur approchée de $3$ à $10^{-2}$ près.
\end{enumerate}
\end{exercice}

\begin{exercice}[Type Bac : suite auxiliaire et remboursement d'un crédit]
Un commerçant de Douala emprunte une somme à une microfinance. On modélise la situation par la suite $(u_n)$ définie par $u_0 = 0$ et, pour tout $n \in \mathbb{N}$ :
\[ u_{n+1} = \frac{3u_n + 2}{u_n + 4}. \]

\textbf{Partie A : premières propriétés}
\begin{enumerate}
\item Calcule $u_1$ et $u_2$.
\item Démontre par récurrence que pour tout $n \in \mathbb{N}$, $0 \leq u_n \leq 1$.
\item Étudie le signe de $u_{n+1} - u_n$ et déduis-en le sens de variation de $(u_n)$.
\item Justifie que $(u_n)$ converge vers une limite $\ell$ que l'on déterminera.
\end{enumerate}

\textbf{Partie B : suite auxiliaire}
\begin{enumerate}
\item On pose, pour tout $n \in \mathbb{N}$, $v_n = \dfrac{u_n - 1}{u_n + 2}$.
\begin{enumerate}
\item Démontre que pour tout $n \in \mathbb{N}$, $v_{n+1} = \dfrac{2}{5}\, v_n$.
\item En déduire la nature de la suite $(v_n)$, puis exprimer $v_n$ en fonction de $n$.
\end{enumerate}
\item Exprime $u_n$ en fonction de $n$ à l'aide de $v_n$, puis sous forme explicite simplifiée.
\item Retrouve la limite de $(u_n)$ à l'aide de l'expression explicite obtenue.
\end{enumerate}

\textbf{Partie C : synthèse}
\begin{enumerate}
\item Démontre par récurrence que pour tout $n \geq 1$ :
\[ \sum_{k=1}^{n} k^2 = \frac{n(n+1)(2n+1)}{6}. \]
\item Le commerçant affirme : « comme $(u_n)$ est croissante et majorée par $1$, elle atteint forcément la valeur $1$ pour un certain rang ». Cette affirmation est-elle correcte ? Justifie rigoureusement.
\end{enumerate}
\end{exercice}

\newpage

\coursec{Corrigés des exercices}

\begin{corrige}[Exercice 1 — QCM : récurrence et monotonie]
\textbf{1. Réponse (a).} Le principe de récurrence exige exactement deux choses : l'\textbf{initialisation} (vérifier $\mathcal{P}(n_0)$) et l'\textbf{hérédité} (montrer que pour tout $n \geq n_0$, $\mathcal{P}(n) \implies \mathcal{P}(n+1)$). Vérifier deux valeurs (réponse b) ne prouve rien pour les suivantes, et montrer $\mathcal{P}(n+1)$ sans initialisation (réponse c) ne permet pas d'amorcer le raisonnement.

\textbf{2. Réponse (b).} On a $u_{n+1} - u_n = \dfrac{1}{n+2} - \dfrac{1}{n+1} = \dfrac{(n+1)-(n+2)}{(n+1)(n+2)} = \dfrac{-1}{(n+1)(n+2)} < 0$, donc $(u_n)$ est \textbf{décroissante}. Comme $\dfrac{1}{n+1} > 0$, on a $u_n > 3$ pour tout $n$ : la suite est \textbf{minorée par $3$}. (Elle est aussi majorée par $u_0 = 4$, donc bornée.)

\textbf{3. Réponse (b).} Par le théorème de convergence monotone, toute suite croissante et majorée converge ; sa limite $\ell$ vérifie $\ell \leq M$ pour toute majorante $M$, donc ici $\ell \leq 5$. Elle ne converge pas forcément vers $5$ : par exemple $u_n = 3 - \dfrac{1}{n+1}$ est croissante, majorée par $5$, et converge vers $3$.

\textbf{4. Réponse (b).} Si $(u_n)$ converge vers $\ell$, alors $(u_{n+1})$ converge aussi vers $\ell$. Comme $f(x) = \dfrac{x+4}{3}$ est continue, en passant à la limite dans $u_{n+1} = f(u_n)$ on obtient $\ell = f(\ell)$, c'est-à-dire $\ell = \dfrac{\ell + 4}{3}$. (Ce qui donne $3\ell = \ell + 4$, soit $\ell = 2$.)
\end{corrige}

\begin{corrige}[Exercice 2 — Démonstrations par récurrence]
\textbf{1.} Soit $\mathcal{P}(n)$ : « $1^2 + 2^2 + \cdots + n^2 = \dfrac{n(n+1)(2n+1)}{6}$ », pour $n \geq 1$.

\textbf{Initialisation ($n=1$) :} $1^2 = 1$ et $\dfrac{1 \times 2 \times 3}{6} = 1$. Donc $\mathcal{P}(1)$ est vraie.

\textbf{Hérédité :} Soit $n \geq 1$ tel que $\mathcal{P}(n)$ est vraie. Alors :
\begin{align*}
\sum_{k=1}^{n+1} k^2 &= \frac{n(n+1)(2n+1)}{6} + (n+1)^2 = (n+1)\left[\frac{n(2n+1)}{6} + (n+1)\right]\\
&= (n+1)\,\frac{2n^2 + n + 6n + 6}{6} = \frac{(n+1)(2n^2+7n+6)}{6}.
\end{align*}
Or $2n^2 + 7n + 6 = (n+2)(2n+3)$ (vérification : $(n+2)(2n+3) = 2n^2 + 3n + 4n + 6$). Donc :
\[ \sum_{k=1}^{n+1} k^2 = \frac{(n+1)(n+2)(2n+3)}{6}, \]
qui est bien $\mathcal{P}(n+1)$.

\textbf{Conclusion :} Par le principe de récurrence, $\mathcal{P}(n)$ est vraie pour tout $n \geq 1$.

\textbf{2.} Soit $\mathcal{P}(n)$ : « $3^{2n} - 2^n$ est divisible par $7$ », pour $n \geq 1$.

\textbf{Initialisation ($n=1$) :} $3^2 - 2 = 9 - 2 = 7$, divisible par $7$. $\mathcal{P}(1)$ est vraie.

\textbf{Hérédité :} Supposons $3^{2n} - 2^n = 7k$ pour un $k \in \mathbb{Z}$. Alors $3^{2n} = 7k + 2^n$, et :
\begin{align*}
3^{2(n+1)} - 2^{n+1} &= 9 \times 3^{2n} - 2 \times 2^n = 9(7k + 2^n) - 2 \times 2^n\\
&= 63k + 9 \times 2^n - 2 \times 2^n = 63k + 7 \times 2^n = 7(9k + 2^n).
\end{align*}
Comme $9k + 2^n \in \mathbb{Z}$, $3^{2(n+1)} - 2^{n+1}$ est divisible par $7$ : $\mathcal{P}(n+1)$ est vraie.

\textbf{Conclusion :} Par récurrence, $3^{2n} - 2^n$ est divisible par $7$ pour tout $n \geq 1$.

\textbf{3.} Soit $\mathcal{P}(n)$ : « $2^n \geq n^2$ », pour $n \geq 4$.

\textbf{Initialisation ($n=4$) :} $2^4 = 16$ et $4^2 = 16$, donc $2^4 \geq 4^2$. $\mathcal{P}(4)$ est vraie.

\textbf{Hérédité :} Supposons $2^n \geq n^2$ pour un $n \geq 4$. Alors :
\[ 2^{n+1} = 2 \times 2^n \geq 2n^2. \]
Il suffit de montrer que $2n^2 \geq (n+1)^2$, c'est-à-dire $2n^2 \geq n^2 + 2n + 1$, soit $n^2 - 2n - 1 \geq 0$. Or pour $n \geq 4$ : $n^2 - 2n - 1 = n(n-2) - 1 \geq 4 \times 2 - 1 = 7 > 0$. Donc :
\[ 2^{n+1} \geq 2n^2 \geq (n+1)^2, \]
et $\mathcal{P}(n+1)$ est vraie.

\textbf{Conclusion :} Par récurrence, $2^n \geq n^2$ pour tout entier $n \geq 4$.

\textbf{Point de vigilance :} l'initialisation à $n = 4$ est indispensable ici : la propriété est \emph{fausse} pour $n = 3$ ($2^3 = 8 < 9$). L'entier de départ fait partie de l'énoncé de la propriété.
\end{corrige}

\begin{corrige}[Exercice 3 — Vrai ou faux : justifier]
\textbf{1. FAUX.} Contre-exemple : $u_n = n$ est croissante ($u_{n+1} - u_n = 1 > 0$) mais diverge vers $+\infty$. Pour conclure à la convergence, il manque l'hypothèse « majorée ».

\textbf{2. VRAI.} C'est exactement le théorème de convergence monotone : toute suite décroissante et minorée (ici par $0$) converge, et sa limite $\ell$ vérifie $\ell \geq 0$.

\textbf{3. FAUX.} Contre-exemple : $u_n = (-1)^n$. On a $|u_n| = 1$, donc $(u_n)$ est bornée (majorée par $1$, minorée par $-1$). Mais $u_{n+1} - u_n = (-1)^{n+1} - (-1)^n = 2(-1)^{n+1}$ change de signe : la suite n'est ni croissante ni décroissante.

\textbf{4. FAUX.} La condition $f(\ell) = \ell$ est \emph{nécessaire} mais pas \emph{suffisante}. Contre-exemple : $f(x) = 2x$, $u_0 = 1$, $u_{n+1} = 2u_n$. L'équation $f(\ell) = \ell$ donne $\ell = 0$, pourtant $u_n = 2^n \to +\infty$ : la suite diverge. Il faut en plus une hypothèse de stabilité, de monotonie et de bornitude (ou de contraction $|f'| \leq k < 1$).
\end{corrige}

\begin{corrige}[Exercice 4 — Définitions et bornitude]
\textbf{1.} Soit $(u_n)$ une suite réelle.
\begin{itemize}
\item $(u_n)$ est \textbf{majorée} s'il existe un réel $M$ tel que pour tout $n \in \mathbb{N}$, $u_n \leq M$.
\item $(u_n)$ est \textbf{minorée} s'il existe un réel $m$ tel que pour tout $n \in \mathbb{N}$, $u_n \geq m$.
\item $(u_n)$ est \textbf{bornée} si elle est majorée \emph{et} minorée, ce qui équivaut à : il existe $M > 0$ tel que pour tout $n$, $|u_n| \leq M$.
\end{itemize}

\textbf{2. (a)} $u_n = \dfrac{2n+1}{n+1}$. Écrivons $u_n = \dfrac{2(n+1) - 1}{n+1} = 2 - \dfrac{1}{n+1}$.
\begin{itemize}
\item Comme $\dfrac{1}{n+1} > 0$, on a $u_n < 2$ : $(u_n)$ est \textbf{majorée par $2$}.
\item Comme $n + 1 \geq 1$, $\dfrac{1}{n+1} \leq 1$, donc $u_n \geq 2 - 1 = 1$ : $(u_n)$ est \textbf{minorée par $1$}.
\end{itemize}
Conclusion : $(u_n)$ est \textbf{bornée} : pour tout $n$, $1 \leq u_n < 2$.

\textbf{2. (b)} $v_n = (-1)^n n$. Pour les rangs pairs, $v_{2p} = 2p \to +\infty$ : $(v_n)$ n'est \textbf{pas majorée}. Pour les rangs impairs, $v_{2p+1} = -(2p+1) \to -\infty$ : $(v_n)$ n'est \textbf{pas minorée}. Conclusion : $(v_n)$ n'est \textbf{pas bornée}.

\textbf{2. (c)} $w_n = 3 - n^2$. Comme $n^2 \geq 0$, $w_n \leq 3$ : $(w_n)$ est \textbf{majorée par $3$}. Mais $w_n \to -\infty$ : pour tout réel $m$, on peut trouver $n$ tel que $3 - n^2 < m$ (il suffit que $n > \sqrt{3 - m}$) : $(w_n)$ n'est \textbf{pas minorée}. Conclusion : $(w_n)$ n'est \textbf{pas bornée}.
\end{corrige}

\begin{corrige}[Exercice 5 — Calculs directs]
\textbf{1.} Avec $u_0 = 1$ et $u_{n+1} = \dfrac{u_n + 4}{3}$ :
\[ u_1 = \frac{1 + 4}{3} = \frac{5}{3} ; \qquad u_2 = \frac{\frac{5}{3} + 4}{3} = \frac{\frac{17}{3}}{3} = \frac{17}{9} ; \qquad u_3 = \frac{\frac{17}{9} + 4}{3} = \frac{\frac{53}{9}}{3} = \frac{53}{27}. \]

\textbf{2.} On a $u_0 = 1 \approx 1$ ; $u_1 \approx 1{,}67$ ; $u_2 \approx 1{,}89$ ; $u_3 \approx 1{,}96$. Les termes augmentent et semblent se rapprocher de $2$. \textbf{Conjecture :} $(u_n)$ est croissante et converge vers $2$. (Vérification du point fixe : $\ell = \dfrac{\ell + 4}{3} \iff 3\ell = \ell + 4 \iff \ell = 2$.)

\textbf{3.} $a_n = \dfrac{n^2}{n+1}$. Calculons :
\begin{align*}
a_{n+1} - a_n &= \frac{(n+1)^2}{n+2} - \frac{n^2}{n+1} = \frac{(n+1)^3 - n^2(n+2)}{(n+1)(n+2)}\\
&= \frac{n^3 + 3n^2 + 3n + 1 - n^3 - 2n^2}{(n+1)(n+2)} = \frac{n^2 + 3n + 1}{(n+1)(n+2)}.
\end{align*}
Pour tout $n \in \mathbb{N}$, le numérateur $n^2 + 3n + 1 > 0$ et le dénominateur $(n+1)(n+2) > 0$, donc $a_{n+1} - a_n > 0$ : la suite $(a_n)$ est \textbf{strictement croissante}.

\textbf{Point de vigilance :} pour étudier la monotonie d'une suite, on calcule $u_{n+1} - u_n$ (ou le quotient $\frac{u_{n+1}}{u_n}$ si tous les termes sont strictement positifs), et non la dérivée d'une fonction — sauf si la suite est de la forme $u_n = f(n)$ avec $f$ monotone sur $[0\,;\,+\infty[$.
\end{corrige}

\begin{corrige}[Exercice 6 — Monotonie et bornitude]
\textbf{1.} Soit $\mathcal{P}(n)$ : « $1 \leq u_n \leq 2$ ».

\textbf{Initialisation :} $u_0 = 1$, donc $1 \leq u_0 \leq 2$ : $\mathcal{P}(0)$ est vraie.

\textbf{Hérédité :} Supposons $1 \leq u_n \leq 2$ pour un $n \in \mathbb{N}$. Alors :
\[ 1 + 4 \leq u_n + 4 \leq 2 + 4 \quad\Longrightarrow\quad 5 \leq u_n + 4 \leq 6 \quad\Longrightarrow\quad \frac{5}{3} \leq u_{n+1} \leq 2. \]
Comme $\dfrac{5}{3} \geq 1$, on a bien $1 \leq u_{n+1} \leq 2$ : $\mathcal{P}(n+1)$ est vraie.

\textbf{Conclusion :} Par récurrence, pour tout $n \in \mathbb{N}$, $1 \leq u_n \leq 2$.

\textbf{2.} Pour tout $n$ :
\[ u_{n+1} - u_n = \frac{u_n + 4}{3} - u_n = \frac{u_n + 4 - 3u_n}{3} = \frac{4 - 2u_n}{3} = \frac{2(2 - u_n)}{3}. \]
D'après la question 1, $u_n \leq 2$, donc $2 - u_n \geq 0$, donc $u_{n+1} - u_n \geq 0$ : la suite $(u_n)$ est \textbf{croissante}.

\textbf{3.} La suite $(u_n)$ est \textbf{croissante} (question 2) et \textbf{majorée par $2$} (question 1). D'après le \textbf{théorème de convergence monotone}, elle converge vers une limite $\ell \leq 2$.

La fonction $f(x) = \dfrac{x+4}{3}$ est continue sur $\mathbb{R}$. En passant à la limite dans $u_{n+1} = f(u_n)$ :
\[ \ell = \frac{\ell + 4}{3} \iff 3\ell = \ell + 4 \iff 2\ell = 4 \iff \ell = 2. \]
Conclusion : $\lim\limits_{n \to +\infty} u_n = 2$.
\end{corrige}

\begin{corrige}[Exercice 7 — Suite auxiliaire géométrique]
\textbf{1.} $u_1 = \dfrac{3 \times 0 + 2}{0 + 4} = \dfrac{1}{2}$ ; \quad $u_2 = \dfrac{3 \times \frac{1}{2} + 2}{\frac{1}{2} + 4} = \dfrac{\frac{7}{2}}{\frac{9}{2}} = \dfrac{7}{9}$.

\textbf{2.} Soit $\mathcal{P}(n)$ : « $u_n \geq 0$ ».

\textbf{Initialisation :} $u_0 = 0 \geq 0$. \textbf{Hérédité :} si $u_n \geq 0$, alors $3u_n + 2 \geq 2 > 0$ et $u_n + 4 \geq 4 > 0$, donc $u_{n+1} = \dfrac{3u_n+2}{u_n+4} \geq 0$. \textbf{Conclusion :} par récurrence, $u_n \geq 0$ pour tout $n$.

\textbf{3. (a)} Pour tout $n$ :
\begin{align*}
v_{n+1} &= \frac{u_{n+1} - 1}{u_{n+1} + 2} = \frac{\dfrac{3u_n + 2}{u_n + 4} - 1}{\dfrac{3u_n + 2}{u_n + 4} + 2} = \frac{\dfrac{3u_n + 2 - (u_n + 4)}{u_n + 4}}{\dfrac{3u_n + 2 + 2(u_n + 4)}{u_n + 4}} = \frac{2u_n - 2}{5u_n + 10} = \frac{2(u_n - 1)}{5(u_n + 2)} = \frac{2}{5}\, v_n.
\end{align*}
Donc $(v_n)$ est \textbf{géométrique de raison $q = \dfrac{2}{5}$}, de premier terme $v_0 = \dfrac{0 - 1}{0 + 2} = -\dfrac{1}{2}$.

\textbf{3. (b)} $v_n = v_0 \times q^n = -\dfrac{1}{2}\left(\dfrac{2}{5}\right)^n$.

De $v_n = \dfrac{u_n - 1}{u_n + 2}$, on tire $v_n(u_n + 2) = u_n - 1$, soit $u_n(v_n - 1) = -1 - 2v_n$, donc :
\[ u_n = \frac{1 + 2v_n}{1 - v_n} = \frac{1 - \left(\frac{2}{5}\right)^n}{1 + \frac{1}{2}\left(\frac{2}{5}\right)^n}. \]

\textbf{3. (c)} Comme $\left|\dfrac{2}{5}\right| < 1$, $\left(\dfrac{2}{5}\right)^n \to 0$, donc $v_n \to 0$ et :
\[ \lim_{n \to +\infty} u_n = \frac{1 - 0}{1 + 0} = 1. \]
\end{corrige}

\begin{corrige}[Exercice 8 — Encadrement par l'inégalité des accroissements finis]
\textbf{1.} $f(x) = \sqrt{2x+3}$ est dérivable sur $[0\,;\,3]$ et $f'(x) = \dfrac{2}{2\sqrt{2x+3}} = \dfrac{1}{\sqrt{2x+3}} > 0$. Donc $f$ est \textbf{strictement croissante} sur $[0\,;\,3]$. Par conséquent, pour $x \in [0\,;\,3]$ :
\[ f(0) \leq f(x) \leq f(3) \quad\text{avec}\quad f(0) = \sqrt{3} \approx 1{,}73 \geq 0 \quad\text{et}\quad f(3) = \sqrt{9} = 3. \]
Donc $f(x) \in [\sqrt{3}\,;\,3] \subset [0\,;\,3]$ : l'intervalle $[0\,;\,3]$ est \textbf{stable par $f$}.

\textbf{2.} Pour $x \in [0\,;\,3]$ : $f(x) = x \iff \sqrt{2x+3} = x$. Les deux membres étant positifs, cela équivaut à $2x + 3 = x^2$, soit $x^2 - 2x - 3 = 0$. Discriminant : $\Delta = 4 + 12 = 16$ ; racines : $x = \dfrac{2 - 4}{2} = -1$ et $x = \dfrac{2+4}{2} = 3$. Seule $x = 3$ appartient à $[0\,;\,3]$. \textbf{Solution : $x = 3$.}

\textbf{3.} $f'(x) = \dfrac{1}{\sqrt{2x+3}}$ est décroissante sur $[0\,;\,3]$ (le dénominateur est croissant). Donc pour $x \in [0\,;\,3]$ :
\[ f'(3) \leq f'(x) \leq f'(0) \quad\text{soit}\quad \frac{1}{\sqrt{9}} \leq f'(x) \leq \frac{1}{\sqrt{3}} \quad\text{soit}\quad \frac{1}{3} \leq f'(x) \leq \frac{1}{\sqrt{3}}. \]

\textbf{4.} $f$ est dérivable sur $[0\,;\,3]$, et $u_n \in [0\,;\,3]$ pour tout $n$ (stabilité, avec $u_0 = 0$). D'après l'\textbf{inégalité des accroissements finis} appliquée entre $u_n$ et $3$ :
\[ |f(u_n) - f(3)| \leq \frac{1}{\sqrt{3}}\,|u_n - 3|. \]
Or $f(u_n) = u_{n+1}$ et $f(3) = 3$ (point fixe), donc :
\[ |u_{n+1} - 3| \leq \frac{1}{\sqrt{3}}\,|u_n - 3|. \]
Par itération (récurrence immédiate) : $|u_n - 3| \leq \left(\dfrac{1}{\sqrt{3}}\right)^n |u_0 - 3| = 3\left(\dfrac{1}{\sqrt{3}}\right)^n$, car $|u_0 - 3| = 3$.

\textbf{5.} On cherche $N$ tel que $3\left(\dfrac{1}{\sqrt{3}}\right)^N \leq 10^{-2}$, soit $\left(\dfrac{1}{\sqrt{3}}\right)^N \leq \dfrac{1}{300}$. En passant au logarithme (fonction croissante) :
\[ N \ln\left(\frac{1}{\sqrt{3}}\right) \leq -\ln(300) \iff N \times \left(-\frac{\ln 3}{2}\right) \leq -\ln 300 \iff N \geq \frac{2\ln 300}{\ln 3}. \]
Numériquement : $\ln 300 \approx 5{,}704$ et $\ln 3 \approx 1{,}0986$, donc $N \geq \dfrac{2 \times 5{,}704}{1{,}0986} \approx 10{,}38$. Le plus petit entier convenant est $\boxed{N = 11}$.

Vérification : $3 \times \left(\frac{1}{\sqrt{3}}\right)^{11} = \dfrac{3}{3^{5{,}5}} = 3^{-4{,}5} \approx \dfrac{1}{140{,}3} \approx 7{,}1 \times 10^{-3} \leq 10^{-2}$. Et pour $n = 10$ : $3^{-4} = \dfrac{1}{81} \approx 1{,}23 \times 10^{-2} > 10^{-2}$. Donc $N = 11$ est bien minimal.
\end{corrige}

\begin{corrige}[Exercice 9 — Type Bac : pisciculture à Maroua]
\textbf{Partie A}

\textbf{1.} Au début du mois $n$, il y a $p_n$ poissons. La reproduction augmente la population de $10\,\%$ : elle devient $p_n + 0{,}1\,p_n = 1{,}1\,p_n$. Puis le prélèvement retire $400$ poissons. Donc :
\[ p_{n+1} = 1{,}1\,p_n - 400. \]

\textbf{2.} $p_1 = 1{,}1 \times 5000 - 400 = 5500 - 400 = 5100$ ; \quad $p_2 = 1{,}1 \times 5100 - 400 = 5610 - 400 = 5210$.

\textbf{Partie B}

\textbf{1. (a)} $q_{n+1} = p_{n+1} - 4000 = 1{,}1\,p_n - 400 - 4000 = 1{,}1\,p_n - 4400 = 1{,}1(p_n - 4000) = 1{,}1\,q_n$. Donc $(q_n)$ est \textbf{géométrique de raison $1{,}1$}, de premier terme $q_0 = 5000 - 4000 = 1000$.

\textbf{1. (b)} $q_n = 1000 \times (1{,}1)^n$, donc $p_n = 4000 + 1000 \times (1{,}1)^n$.

\textbf{2.} $p_{n+1} - p_n = 1000 \times (1{,}1)^{n+1} - 1000 \times (1{,}1)^n = 1000 \times (1{,}1)^n \times 0{,}1 = 100 \times (1{,}1)^n > 0$ : $(p_n)$ est \textbf{strictement croissante}. Comme $1{,}1 > 1$, $(1{,}1)^n \to +\infty$, donc $\lim p_n = +\infty$.

\textbf{Interprétation :} le prélèvement de $400$ poissons est inférieur au gain net de reproduction ; la population croît sans limite (en pratique, jusqu'aux contraintes de capacité du bassin). L'exploitation est viable et rentable.

\textbf{3.} On résout $p_n > 8000 \iff 1000 \times (1{,}1)^n > 4000 \iff (1{,}1)^n > 4 \iff n \ln 1{,}1 > \ln 4 \iff n > \dfrac{\ln 4}{\ln 1{,}1} \approx \dfrac{1{,}3863}{0{,}0953} \approx 14{,}55$. Le premier mois est donc $n = 15$. Vérification : $(1{,}1)^{14} \approx 3{,}798$ ($p_{14} \approx 7798$) et $(1{,}1)^{15} \approx 4{,}177$ ($p_{15} \approx 8177 > 8000$). \textbf{La population dépasse \num{8000} poissons au mois $15$.}

\textbf{Partie C}

\textbf{1.} $S$ vérifie $S = 1{,}1\,S - 400 \iff 0{,}1\,S = 400 \iff S = 4000$. La population stable est de \textbf{\num{4000} poissons}.

\textbf{2.} Si $p_0 = 3000$, alors $q_0 = 3000 - 4000 = -1000$, et $q_n = -1000 \times (1{,}1)^n$, donc $p_n = 4000 - 1000 \times (1{,}1)^n$. La suite $(p_n)$ est alors \textbf{strictement décroissante} et tend vers $-\infty$ : elle atteint des valeurs négatives, ce qui signifie concrètement que le bassin se vide. En effet, $p_n \leq 0$ dès que $(1{,}1)^n \geq 4$, soit $n \geq 15$ : le stock s'épuise en un peu plus d'un an. \textbf{Conclusion :} avec \num{3000} poissons initiaux, le prélèvement mensuel de $400$ poissons dépasse la capacité de reproduction ; l'exploitation n'est \textbf{pas viable}. Le seuil critique de viabilité est exactement $S = \num{4000}$ poissons.

\textbf{Barème indicatif (sur 10) :} A1 : 1 pt ; A2 : 1 pt ; B1a : 1,5 pt ; B1b : 1 pt ; B2 : 1,5 pt ; B3 : 1 pt ; C1 : 1 pt ; C2 : 2 pts.

\textbf{Points de vigilance :} justifier la raison $1{,}1$ par le calcul explicite de $q_{n+1}$ en fonction de $q_n$ ; pour la question B3, exiger la résolution d'inéquation (ou un tableau de valeurs) et non une simple lecture graphique ; en C2, la conclusion doit être interprétée concrètement (épuisement du stock).
\end{corrige}

\begin{corrige}[Exercice 10 — Type Bac : étude complète d'une suite récurrente]
\textbf{Partie A}

\textbf{1.} $f(x) = \sqrt{2x+3}$ est dérivable sur $[0\,;\,+\infty[$ et $f'(x) = \dfrac{1}{\sqrt{2x+3}} > 0$ pour tout $x \geq 0$. Donc $f$ est \textbf{strictement croissante} sur $[0\,;\,+\infty[$, de $f(0) = \sqrt{3}$ vers $+\infty$.

\textbf{2.} Soit $x \in [0\,;\,3]$. Par croissance de $f$ : $f(0) \leq f(x) \leq f(3)$, c'est-à-dire $\sqrt{3} \leq f(x) \leq \sqrt{9} = 3$. Comme $\sqrt{3} \geq 0$, on a $f(x) \in [0\,;\,3]$ : $[0\,;\,3]$ est \textbf{stable par $f$}.

\textbf{3.} Pour $x \geq 0$, $f(x) = x \iff \sqrt{2x+3} = x \iff 2x+3 = x^2$ (les deux membres sont positifs, l'équivalence est licite) $\iff x^2 - 2x - 3 = 0 \iff (x-3)(x+1) = 0 \iff x = 3$ ou $x = -1$. Seule la solution $x = 3$ appartient à $[0\,;\,3]$. L'équation admet donc une \textbf{unique solution dans $[0\,;\,3]$, égale à $3$}.

\textbf{Partie B}

\textbf{1.} Soit $\mathcal{P}(n)$ : « $0 \leq u_n \leq 3$ ». \textbf{Initialisation :} $u_0 = 0 \in [0\,;\,3]$. \textbf{Hérédité :} si $u_n \in [0\,;\,3]$, alors $u_{n+1} = f(u_n) \in [0\,;\,3]$ par la stabilité démontrée en A2. \textbf{Conclusion :} par récurrence, $0 \leq u_n \leq 3$ pour tout $n$.

\textbf{2.} Pour tout $n$ :
\[ u_{n+1} - u_n = \sqrt{2u_n + 3} - u_n. \]
Comme $u_n \geq 0$, les deux quantités $\sqrt{2u_n+3}$ et $u_n$ sont positives ; multiplier par la quantité conjuguée (positive) :
\[ u_{n+1} - u_n = \frac{(\sqrt{2u_n+3} - u_n)(\sqrt{2u_n+3} + u_n)}{\sqrt{2u_n+3} + u_n} = \frac{2u_n + 3 - u_n^2}{\sqrt{2u_n+3} + u_n}. \]
Le dénominateur est strictement positif (pour $n \geq 1$ ; et pour $n = 0$, $\sqrt{3} > 0$), donc $u_{n+1} - u_n$ a le \textbf{même signe que $-u_n^2 + 2u_n + 3$}.

Or $-x^2 + 2x + 3 = -(x^2 - 2x - 3) = -(x-3)(x+1) = (3-x)(x+1)$. Pour $u_n \in [0\,;\,3]$ : $3 - u_n \geq 0$ et $u_n + 1 > 0$, donc $-u_n^2 + 2u_n + 3 \geq 0$. Par suite $u_{n+1} - u_n \geq 0$ : $(u_n)$ est \textbf{croissante}.

\textbf{3.} $(u_n)$ est \textbf{croissante} (B2) et \textbf{majorée par $3$} (B1) : par le \textbf{théorème de convergence monotone}, elle converge vers une limite $\ell \in [0\,;\,3]$. Par continuité de $f$, $\ell$ vérifie $f(\ell) = \ell$ ; d'après A3, la seule solution dans $[0\,;\,3]$ est $\ell = 3$. Donc $\lim u_n = 3$.

\textbf{Partie C}

\textbf{1.} $f'(x) = \dfrac{1}{\sqrt{2x+3}}$ est décroissante sur $[0\,;\,3]$, donc pour $x \in [0\,;\,3]$ : $f'(3) \leq f'(x) \leq f'(0)$, soit $\dfrac{1}{3} \leq f'(x) \leq \dfrac{1}{\sqrt{3}}$.

\textbf{2.} D'après C1, $|f'(x)| \leq \dfrac{1}{\sqrt{3}}$ sur $[0\,;\,3]$. L'\textbf{inégalité des accroissements finis} appliquée entre $u_n \in [0\,;\,3]$ et $3$ donne :
\[ |f(u_n) - f(3)| \leq \frac{1}{\sqrt{3}}|u_n - 3| \quad\text{soit}\quad |u_{n+1} - 3| \leq \frac{1}{\sqrt{3}}|u_n - 3|, \]
car $f(3) = 3$. Par récurrence immédiate : $|u_n - 3| \leq \left(\dfrac{1}{\sqrt{3}}\right)^n |u_0 - 3| = 3\left(\dfrac{1}{\sqrt{3}}\right)^n$.

\textbf{3.} Il faut $3\left(\dfrac{1}{\sqrt{3}}\right)^N \leq 10^{-2}$, soit $\left(\dfrac{1}{\sqrt{3}}\right)^N \leq \dfrac{1}{300}$, soit $N \geq \dfrac{2\ln 300}{\ln 3} \approx 10{,}38$. Le plus petit entier est $\boxed{N = 11}$ (voir la vérification détaillée dans le corrigé de l'exercice 8).

\textbf{Barème indicatif (sur 10) :} A1 : 1 pt ; A2 : 1 pt ; A3 : 1,5 pt ; B1 : 1 pt ; B2 : 2 pts ; B3 : 1,5 pt ; C1 : 0,5 pt ; C2 : 1 pt ; C3 : 0,5 pt.

\textbf{Points de vigilance :}
\begin{itemize}
\item En A3, l'équivalence $\sqrt{2x+3} = x \iff 2x+3 = x^2$ n'est valable que parce que $x \geq 0$ : le citer.
\item En B2, la justification du signe doit passer par la factorisation $(3 - u_n)(u_n + 1)$ et l'encadrement de B1.
\item En B3, il faut citer explicitement le théorème de convergence monotone \emph{et} la continuité de $f$ pour justifier $\ell = f(\ell)$.
\item En C2, l'IAF s'applique car $f$ est dérivable sur $[0\,;\,3]$ et $|f'|$ y est majorée par $\frac{1}{\sqrt{3}} < 1$ : ces hypothèses doivent être écrites.
\end{itemize}
\end{corrige}

\begin{corrige}[Exercice 11 — Type Bac : suite auxiliaire et remboursement d'un crédit]
\textbf{Partie A}

\textbf{1.} $u_1 = \dfrac{3 \times 0 + 2}{0 + 4} = \dfrac{1}{2}$ ; \quad $u_2 = \dfrac{3 \times \frac{1}{2} + 2}{\frac{1}{2} + 4} = \dfrac{\frac{7}{2}}{\frac{9}{2}} = \dfrac{7}{9}$.

\textbf{2.} Soit $\mathcal{P}(n)$ : « $0 \leq u_n \leq 1$ ».

\textbf{Initialisation :} $u_0 = 0$, donc $\mathcal{P}(0)$ est vraie.

\textbf{Hérédité :} Supposons $0 \leq u_n \leq 1$. Alors $3u_n + 2 \geq 2 > 0$ et $u_n + 4 \geq 4 > 0$, donc $u_{n+1} \geq 0$. De plus :
\[ u_{n+1} - 1 = \frac{3u_n + 2}{u_n + 4} - 1 = \frac{3u_n + 2 - u_n - 4}{u_n + 4} = \frac{2u_n - 2}{u_n + 4} = \frac{2(u_n - 1)}{u_n + 4} \leq 0, \]
car $u_n \leq 1$ et $u_n + 4 > 0$. Donc $u_{n+1} \leq 1$ : $\mathcal{P}(n+1)$ est vraie.

\textbf{Conclusion :} par récurrence, $0 \leq u_n \leq 1$ pour tout $n$.

\textbf{3.} Pour tout $n$ :
\[ u_{n+1} - u_n = \frac{3u_n + 2}{u_n + 4} - u_n = \frac{3u_n + 2 - u_n^2 - 4u_n}{u_n + 4} = \frac{-u_n^2 - u_n + 2}{u_n + 4} = \frac{-(u_n - 1)(u_n + 2)}{u_n + 4} = \frac{(1 - u_n)(u_n + 2)}{u_n + 4}. \]
(Vérification : $(1-u_n)(u_n+2) = u_n + 2 - u_n^2 - 2u_n = -u_n^2 - u_n + 2$.) D'après la question 2, $1 - u_n \geq 0$, $u_n + 2 > 0$ et $u_n + 4 > 0$, donc $u_{n+1} - u_n \geq 0$ : $(u_n)$ est \textbf{croissante}.

\textbf{4.} $(u_n)$ est \textbf{croissante} et \textbf{majorée par $1$} : par le théorème de convergence monotone, elle converge vers $\ell \in [0\,;\,1]$. La fonction $f(x) = \dfrac{3x+2}{x+4}$ est continue sur $[0\,;\,1]$ (le dénominateur ne s'y annule pas), donc $\ell = f(\ell)$ :
\[ \ell = \frac{3\ell + 2}{\ell + 4} \iff \ell^2 + 4\ell = 3\ell + 2 \iff \ell^2 + \ell - 2 = 0 \iff (\ell - 1)(\ell + 2) = 0. \]
Les solutions sont $\ell = 1$ et $\ell = -2$ ; seule $\ell = 1$ appartient à $[0\,;\,1]$. Donc $\lim u_n = 1$.

\textbf{Partie B}

\textbf{1. (a)} Voir le calcul détaillé du corrigé de l'exercice 7, question 3(a) :
\[ v_{n+1} = \frac{u_{n+1} - 1}{u_{n+1} + 2} = \frac{2u_n - 2}{5u_n + 10} = \frac{2(u_n - 1)}{5(u_n + 2)} = \frac{2}{5}\,v_n. \]

\textbf{1. (b)} $(v_n)$ est donc \textbf{géométrique de raison $\dfrac{2}{5}$}, de premier terme $v_0 = \dfrac{0-1}{0+2} = -\dfrac{1}{2}$. D'où :
\[ v_n = -\frac{1}{2}\left(\frac{2}{5}\right)^n. \]

\textbf{2.} De $v_n(u_n + 2) = u_n - 1$, on tire $u_n(v_n - 1) = -2v_n - 1$, soit $u_n = \dfrac{1 + 2v_n}{1 - v_n}$. En remplaçant :
\[ u_n = \frac{1 - \left(\frac{2}{5}\right)^n}{1 + \frac{1}{2}\left(\frac{2}{5}\right)^n} = \frac{2\left[1 - \left(\frac{2}{5}\right)^n\right]}{2 + \left(\frac{2}{5}\right)^n}. \]

\textbf{3.} Comme $0 < \dfrac{2}{5} < 1$, $\left(\dfrac{2}{5}\right)^n \to 0$, donc :
\[ \lim_{n \to +\infty} u_n = \frac{2(1 - 0)}{2 + 0} = 1. \]
On retrouve bien la limite $\ell = 1$ de la partie A.

\textbf{Partie C}

\textbf{1.} Voir la démonstration complète dans le corrigé de l'exercice 2, question 1 : initialisation à $n = 1$ ($1 = \frac{1 \times 2 \times 3}{6}$), hérédité par factorisation de $(n+1)$ et identité $2n^2 + 7n + 6 = (n+2)(2n+3)$, conclusion par le principe de récurrence.

\textbf{2.} \textbf{L'affirmation est fausse.} Le théorème de convergence monotone garantit que $(u_n)$ converge vers sa borne supérieure $\ell = 1$, c'est-à-dire s'en approche aussi près que l'on veut, mais \emph{pas} qu'elle l'atteint. Démonstration : supposons qu'il existe un rang $N$ tel que $u_N = 1$. Si $N \geq 1$, alors $u_N = \dfrac{3u_{N-1} + 2}{u_{N-1} + 4} = 1$ impliquerait $3u_{N-1} + 2 = u_{N-1} + 4$, soit $u_{N-1} = 1$. De proche en proche, on obtiendrait $u_0 = 1$, ce qui contredit $u_0 = 0$. On peut aussi le voir sur l'expression explicite : $u_n = 1$ exigerait $\left(\frac{2}{5}\right)^n = 0$, impossible pour un entier $n$. La suite tend vers $1$ \textbf{sans jamais l'atteindre} : $1$ est une limite, pas une valeur de la suite.

\textbf{Barème indicatif (sur 10) :} A1 : 0,5 pt ; A2 : 1,5 pt ; A3 : 1,5 pt ; A4 : 1,5 pt ; B1 : 1,5 pt ; B2 : 1 pt ; B3 : 0,5 pt ; C1 : 1 pt ; C2 : 1 pt.

\textbf{Points de vigilance :}
\begin{itemize}
\item En A2, la majoration par $1$ doit être démontrée (par exemple via le signe de $u_{n+1} - 1$), pas seulement constatée sur les premiers termes.
\item En A4, citer le théorème de convergence monotone \emph{avec ses deux hypothèses}, puis la continuité de $f$ pour justifier $\ell = f(\ell)$, et \textbf{écarter la solution $-2$} car $\ell \in [0\,;\,1]$.
\item En B1, le calcul de $v_{n+1}$ doit être mené jusqu'au bout en factorisant pour faire apparaître $v_n$.
\item En C2, la réponse « vrai car la limite est $1$ » est un piège classique : une limite n'est pas nécessairement atteinte.
\end{itemize}
\end{corrige}

\newpage

\coursec{Fiche bilan}

\begin{popfigure}
\begin{center}\tcbox[colback=popdark,colframe=popdark,coltext=white,arc=9pt,boxsep=4pt,left=6pt,right=6pt]{\bfseries\small Suites numériques}\end{center}
\vspace{-2pt}
\begin{tcbraster}[raster columns=3,raster equal height=rows,raster column skip=6pt,raster row skip=6pt,enhanced,blankest]
\begin{tcolorbox}[enhanced,colback=white,colframe=popblue,boxrule=0.9pt,arc=3pt,title={Raisonnement par récurrence},fonttitle=\bfseries\scriptsize,coltitle=white,colbacktitle=popblue,left=4pt,right=4pt,top=2pt,bottom=2pt,boxsep=1pt,toptitle=1.5pt,bottomtitle=1.5pt,halign=left]
\begin{itemize}[nosep,leftmargin=*,itemsep=1pt,font=\scriptsize]
\item Initialisation : rang $n_0$
\item Hérédité : $P(n)\Rightarrow P(n+1)$
\item Conclusion : $\forall n\ge n_0$
\end{itemize}
\end{tcolorbox}
\begin{tcolorbox}[enhanced,colback=white,colframe=poporange,boxrule=0.9pt,arc=3pt,title={Bornitude},fonttitle=\bfseries\scriptsize,coltitle=white,colbacktitle=poporange,left=4pt,right=4pt,top=2pt,bottom=2pt,boxsep=1pt,toptitle=1.5pt,bottomtitle=1.5pt,halign=left]
\begin{itemize}[nosep,leftmargin=*,itemsep=1pt,font=\scriptsize]
\item Majorée : $\exists M,\ u_n\le M$
\item Minorée : $\exists m,\ u_n\ge m$
\item Bornée : majorée et minorée
\end{itemize}
\end{tcolorbox}
\begin{tcolorbox}[enhanced,colback=white,colframe=popgreen,boxrule=0.9pt,arc=3pt,title={Monotonie},fonttitle=\bfseries\scriptsize,coltitle=white,colbacktitle=popgreen,left=4pt,right=4pt,top=2pt,bottom=2pt,boxsep=1pt,toptitle=1.5pt,bottomtitle=1.5pt,halign=left]
\begin{itemize}[nosep,leftmargin=*,itemsep=1pt,font=\scriptsize]
\item Différence $u_{n+1}-u_n$
\item Rapport $\frac{u_{n+1}}{u_n}$ si $u_n>0$
\item Fonction $f$ pour $u_{n+1}=f(u_n)$
\end{itemize}
\end{tcolorbox}
\begin{tcolorbox}[enhanced,colback=white,colframe=poppurple,boxrule=0.9pt,arc=3pt,title={Théorème de convergence monotone},fonttitle=\bfseries\scriptsize,coltitle=white,colbacktitle=poppurple,left=4pt,right=4pt,top=2pt,bottom=2pt,boxsep=1pt,toptitle=1.5pt,bottomtitle=1.5pt,halign=left]
\begin{itemize}[nosep,leftmargin=*,itemsep=1pt,font=\scriptsize]
\item Croissante + majorée $\Rightarrow$ converge
\item Décroissante + minorée $\Rightarrow$ converge
\item Limite = majorant/minorant optimal
\end{itemize}
\end{tcolorbox}
\begin{tcolorbox}[enhanced,colback=white,colframe=poppink,boxrule=0.9pt,arc=3pt,title={Suite récurrente $u_{n+1}=f(u_n)$},fonttitle=\bfseries\scriptsize,coltitle=white,colbacktitle=poppink,left=4pt,right=4pt,top=2pt,bottom=2pt,boxsep=1pt,toptitle=1.5pt,bottomtitle=1.5pt,halign=left]
\begin{itemize}[nosep,leftmargin=*,itemsep=1pt,font=\scriptsize]
\item Intervalle stable $I$ : $f(I)\subset I$
\item Monotonie via croissance de $f$
\item Équation point fixe $\ell=f(\ell)$
\end{itemize}
\end{tcolorbox}
\begin{tcolorbox}[enhanced,colback=white,colframe=popgold,boxrule=0.9pt,arc=3pt,title={IAF et approximation},fonttitle=\bfseries\scriptsize,coltitle=white,colbacktitle=popgold,left=4pt,right=4pt,top=2pt,bottom=2pt,boxsep=1pt,toptitle=1.5pt,bottomtitle=1.5pt,halign=left]
\begin{itemize}[nosep,leftmargin=*,itemsep=1pt,font=\scriptsize]
\item $|f'(x)|\le k<1$ sur $I$
\item $|u_{n+1}-\ell|\le k|u_n-\ell|$
\item Rang $n$ pour précision $10^{-p}$
\end{itemize}
\end{tcolorbox}
\end{tcbraster}

\legende{Carte mentale de la leçon M04 : les six compétences clés du programme officiel.}
\end{popfigure}

\coursec{Raisonnement par récurrence}

\begin{definition}[Principe de récurrence]
Soit $P(n)$ une propriété dépendant d'un entier naturel $n$. Si :
\begin{enumerate}
  \item \textbf{Initialisation :} $P(n_0)$ est vraie pour un certain entier $n_0$ (souvent $0$ ou $1$),
  \item \textbf{Hérédité :} pour tout entier $n\ge n_0$, $P(n)$ vraie $\implies$ $P(n+1)$ vraie,
\end{enumerate}
alors $P(n)$ est vraie pour \textbf{tout} entier $n\ge n_0$.
\end{definition}

\begin{attention}[Rédaction rigoureuse au Bac]
\begin{itemize}
  \item Ne jamais écrire « $P(n)$ est vraie » sans quantifier $n$. Écrire : « Soit $n\ge n_0$ tel que $P(n)$ soit vraie (hypothèse de récurrence) ».
  \item L'hérédité doit être démontrée \emph{pour tout $n\ge n_0$}, pas seulement pour un exemple.
  \item La conclusion doit être explicite : « Par le principe de récurrence, $P(n)$ est vraie pour tout entier $n\ge n_0$. »
  \item Attention aux indices : si la suite commence à $u_1$, l'initialisation se fait à $n=1$.
\end{itemize}
\end{attention}

\begin{methode}[Démontrer une inégalité ou propriété sur une suite par récurrence]
\begin{enumerate}
  \item \textbf{Initialisation :} Vérifier la propriété pour le premier indice ($n=0$ ou $n=1$). Calculer explicitement.
  \item \textbf{Hypothèse de récurrence :} Supposer la propriété vraie pour un $n\ge n_0$ arbitraire.
  \item \textbf{Hérédité :} En déduire la propriété pour $n+1$ en utilisant la relation de récurrence (ex: $u_{n+1}=f(u_n)$) et l'hypothèse.
  \item \textbf{Conclusion :} Appliquer le principe de récurrence.
\end{enumerate}
\end{methode}

\begin{exemplebox}[Récurrence sur une suite définie par récurrence]
Soit $(u_n)$ définie par $u_0=2$ et $u_{n+1}=\dfrac{u_n+6}{2}$. Démontrer que $\forall n\in\mathbb{N},\ 2\le u_n\le 4$.
\begin{itemize}
  \item \textbf{Initialisation ($n=0$) :} $u_0=2$. On a bien $2\le 2\le 4$.
  \item \textbf{Hypothèse :} Soit $n\in\mathbb{N}$ tel que $2\le u_n\le 4$.
  \item \textbf{Hérédité :} $u_{n+1}=\dfrac{u_n+6}{2}$. Comme $2\le u_n\le 4$, on a $8\le u_n+6\le 10$, donc $4\le \dfrac{u_n+6}{2}\le 5$.
  \item \textbf{Correction :} L'encadrement obtenu est $4\le u_{n+1}\le 5$, ce qui ne donne pas $u_{n+1}\le 4$. La propriété $u_n\le 4$ n'est \textbf{pas} héréditaire.
  \item \textbf{Recherche d'un majorant héréditaire :} On cherche $M$ tel que $u_n\le M \implies \frac{u_n+6}{2}\le M \iff u_n\le 2M-6$. Il faut $M\ge 2M-6 \iff M\le 6$. Prenons $M=6$.
  \item \textbf{Nouvelle hérédité :} $2\le u_n\le 6 \implies 8\le u_n+6\le 12 \implies 4\le u_{n+1}\le 6$. Donc $2\le u_{n+1}\le 6$.
  \item \textbf{Conclusion :} $\forall n\in\mathbb{N},\ 2\le u_n\le 6$.
\end{itemize}
\end{exemplebox}

\coursec{Suites majorées, minorées, bornées}

\begin{definition}[Majorée, minorée, bornée]
Soit $(u_n)$ une suite à valeurs dans $\mathbb{R}$.
\begin{itemize}
  \item $(u_n)$ est \cle{majorée} s'il existe un réel $M$ tel que $\forall n\in\mathbb{N},\ u_n\le M$. $M$ est un \cle{majorant}.
  \item $(u_n)$ est \cle{minorée} s'il existe un réel $m$ tel que $\forall n\in\mathbb{N},\ u_n\ge m$. $m$ est un \cle{minorant}.
  \item $(u_n)$ est \cle{bornée} si elle est à la fois majorée et minorée, c'est-à-dire $\exists m,M\in\mathbb{R},\ \forall n,\ m\le u_n\le M$.
\end{itemize}
\end{definition}

\begin{propriete}[Majorant/Minorant et récurrence]
Pour prouver qu'une suite récurrente $u_{n+1}=f(u_n)$ est majorée (resp. minorée) par $M$ (resp. $m$), on utilise la récurrence en vérifiant que $f([m,M])\subset [m,M]$ (intervalle stable).
\end{propriete}

\begin{exemplebox}[Bornitude de $u_{n+1}=\sqrt{u_n+2}$]
$u_0=0$, $u_{n+1}=\sqrt{u_n+2}$. Montrer que $0\le u_n\le 2$.
\begin{itemize}
  \item $f(x)=\sqrt{x+2}$ est croissante sur $[-2,+\infty[$.
  \item $f([0,2])=[\sqrt{2},\sqrt{4}]=[\sqrt{2},2]\subset [0,2]$. L'intervalle $[0,2]$ est stable.
  \item Récurrence : $u_0=0\in[0,2]$. Si $u_n\in[0,2]$, $u_{n+1}=f(u_n)\in f([0,2])\subset[0,2]$.
\end{itemize}
\end{exemplebox}

\coursec{Suites monotones et théorème de convergence monotone}

\begin{definition}[Monotonie]
\begin{itemize}
  \item $(u_n)$ est \cle{croissante} si $\forall n,\ u_{n+1}\ge u_n$.
  \item $(u_n)$ est \cle{décroissante} si $\forall n,\ u_{n+1}\le u_n$.
  \item $(u_n)$ est \cle{monotone} si elle est croissante ou décroissante.
  \item \cle{Strictement} croissante/décroissante si les inégalités sont strictes.
\end{itemize}
\end{definition}

\begin{methode}[Étudier la monotonie]
\begin{enumerate}
  \item \textbf{Méthode de la différence :} Étudier le signe de $u_{n+1}-u_n$. Factoriser si possible.
  \item \textbf{Méthode du rapport :} Si $\forall n,\ u_n>0$, étudier $\dfrac{u_{n+1}}{u_n}$ par rapport à $1$.
  \item \textbf{Méthode par récurrence :} Si la suite est définie par $u_{n+1}=f(u_n)$ avec $f$ croissante, comparer $u_1$ et $u_0$ puis propager par récurrence.
\end{enumerate}
\end{methode}

\begin{propriete}[Théorème de convergence monotone -- Exigible au Bac]
\begin{itemize}
  \item Toute suite \textbf{croissante} et \textbf{majorée} converge vers sa borne supérieure (son plus petit majorant).
  \item Toute suite \textbf{décroissante} et \textbf{minorée} converge vers sa borne inférieure (son plus grand minorant).
\end{itemize}
\textbf{Attention :} Les deux hypothèses (monotonie \textbf{ET} bornitude) sont indispensables. Une suite croissante non majorée diverge vers $+\infty$.
\end{propriete}

\begin{experience}[Démonstration guidée du théorème (croissante + majorée)]
Soit $(u_n)$ croissante et majorée. L'ensemble $A=\{u_n,\ n\in\mathbb{N}\}$ est non vide et majoré dans $\mathbb{R}$.
\begin{enumerate}
  \item D'après l'axiome de la borne supérieure (propriété fondamentale de $\mathbb{R}$), $A$ admet un supremum $\ell=\sup A$.
  \item Montrer que $u_n\to\ell$ : Soit $\varepsilon>0$. $\ell-\varepsilon$ n'est pas un majorant de $A$, donc $\exists n_0,\ u_{n_0}>\ell-\varepsilon$.
  \item Comme $(u_n)$ est croissante, $\forall n\ge n_0,\ u_n\ge u_{n_0}>\ell-\varepsilon$.
  \item De plus $\forall n,\ u_n\le \ell < \ell+\varepsilon$.
  \item Donc $\forall n\ge n_0,\ |u_n-\ell|<\varepsilon$. CQFD.
\end{enumerate}
\end{experience}

\begin{exemplebox}[Application du théorème]
$u_0=1$, $u_{n+1}=\sqrt{u_n+1}$.
\begin{itemize}
  \item $f(x)=\sqrt{x+1}$ croissante sur $[0,+\infty[$. $u_1=\sqrt{2}>1=u_0 \implies$ suite croissante (par récurrence).
  \item Majorée ? $f(x)\le x \iff \sqrt{x+1}\le x \iff x+1\le x^2 \iff x^2-x-1\ge 0$. Racine positive $\varphi=\frac{1+\sqrt{5}}{2}\approx 1,618$. Pour $x\in[1,\varphi]$, $f(x)\le \varphi$. Par récurrence $u_n\le \varphi$.
  \item Suite croissante majorée $\implies$ convergente. Limite $\ell=\varphi$.
\end{itemize}
\end{exemplebox}

\coursec{Suites récurrentes $u_{n+1}=f(u_n)$ : étude de la convergence}

\begin{propriete}[Monotonie via la fonction $f$]
Soit $u_{n+1}=f(u_n)$ avec $f$ continue sur un intervalle $I$.
Si $f$ est \textbf{croissante} sur $I$ et $I$ est \textbf{stable} par $f$ ($f(I)\subset I$) et $u_0\in I$, alors :
\begin{itemize}
  \item si $u_1\ge u_0$, la suite est croissante ;
  \item si $u_1\le u_0$, la suite est décroissante.
\end{itemize}
\emph{Preuve par récurrence immédiate.}
\end{propriete}

\begin{methode}[Étude complète d'une suite $u_{n+1}=f(u_n)$ au Bac]
\begin{enumerate}
  \item \textbf{Domaine et continuité :} Préciser l'ensemble de définition de $f$ et sa continuité.
  \item \textbf{Intervalle stable $I$ :} Trouver un intervalle $I=[m,M]$ (ou $[m,+\infty[$, etc.) tel que $u_0\in I$ et $f(I)\subset I$. \textbf{Prouver par récurrence} que $\forall n,\ u_n\in I$.
  \item \textbf{Monotonie :} Étudier le signe de $u_{n+1}-u_n = f(u_n)-u_n$ sur $I$ (ou utiliser la croissance de $f$ et comparer $u_1$ à $u_0$). Conclure : croissante ou décroissante.
  \item \textbf{Convergence :} Appliquer le théorème de convergence monotone (monotonie + bornitude fournie par $I$). La suite converge vers $\ell\in I$.
  \item \textbf{Valeur de la limite :} Passer à la limite dans $u_{n+1}=f(u_n)$ (justifier par continuité de $f$) : $\ell=f(\ell)$. Résoudre l'équation. Choisir la solution compatible avec l'encadrement $m\le \ell\le M$ (et le sens de variation).
\end{enumerate}
\end{methode}

\begin{attention}[Pièges classiques]
\begin{itemize}
  \item Oublier de prouver la stabilité de l'intervalle $I$ par récurrence.
  \item Oublier de justifier la continuité de $f$ avant de passer à la limite.
  \item Résoudre $\ell=f(\ell)$ mais garder une racine extérieure à $I$ ou incompatible avec la monotonie (ex: limite négative pour une suite positive).
  \item Confondre « suite convergente » et « suite ayant une limite finie » (c'est la même chose en Terminale, mais « diverger vers $+\infty$ » n'est pas converger).
\end{itemize}
\end{attention}

\begin{exoresolu}[Étude classique : $u_0=3$, $u_{n+1}=\frac{1}{2}\left(u_n+\frac{5}{u_n}\right)$]
\textbf{Énoncé :} Étudier la convergence de $(u_n)$ et déterminer sa limite. Donner une valeur approchée à $10^{-4}$ près.
\tcblower
\textbf{Solution :}
\begin{enumerate}
  \item $f(x)=\frac{1}{2}(x+\frac{5}{x})$ définie et continue sur $\mathbb{R}^*_+$. $u_0=3>0 \implies \forall n,\ u_n>0$.
  \item \textbf{Intervalle stable :} $f'(x)=\frac{1}{2}(1-\frac{5}{x^2})$. $f$ décroissante sur $]0,\sqrt{5}]$, croissante sur $[\sqrt{5},+\infty[$. Minimum $f(\sqrt{5})=\sqrt{5}$.
      $u_0=3>\sqrt{5}$. Pour $x\ge\sqrt{5}$, $f(x)\ge\sqrt{5}$. De plus $f(x)\le x \iff \frac{1}{2}(x+\frac{5}{x})\le x \iff x+\frac{5}{x}\le 2x \iff \frac{5}{x}\le x \iff x^2\ge 5$, vrai pour $x\ge\sqrt{5}$.
      Donc $f([\sqrt{5},3])\subset [\sqrt{5},3]$. Par récurrence, $\forall n,\ \sqrt{5}\le u_n\le 3$.
  \item \textbf{Monotonie :} Sur $[\sqrt{5},3]$, $f(x)\le x \implies u_{n+1}\le u_n$. Suite décroissante.
  \item \textbf{Convergence :} Décroissante et minorée par $\sqrt{5}$ $\implies$ converge vers $\ell\ge\sqrt{5}$.
  \item \textbf{Limite :} $\ell=\frac{1}{2}(\ell+\frac{5}{\ell}) \iff 2\ell^2=\ell^2+5 \iff \ell^2=5 \iff \ell=\sqrt{5}$ (car $\ell>0$).
  \item \textbf{Approximation (IAF) :} Voir section suivante.
\end{enumerate}
\end{exoresolu}

\coursec{Valeur approchée de la limite par l'inégalité des accroissements finis (IAF)}

\begin{propriete}[Inégalité des accroissements finis (IAF)]
Soit $f$ dérivable sur un intervalle $I$. Pour tout $a,b\in I$, il existe $c$ entre $a$ et $b$ tel que $f(b)-f(a)=f'(c)(b-a)$.
\emph{Conséquence :} Si $|f'(x)|\le k$ sur $I$, alors $|f(b)-f(a)|\le k|b-a|$ ($f$ est $k$-lipschitzienne).
\end{propriete}

\begin{propriete}[Encadrement de l'erreur pour $u_{n+1}=f(u_n)$]
Soit $(u_n)$ définie par $u_{n+1}=f(u_n)$ convergente vers $\ell$, avec $f$ dérivable sur un intervalle $I$ stable contenant $\ell$ et tous les $u_n$.
Si $\exists k\in[0,1[$ tel que $\forall x\in I,\ |f'(x)|\le k$, alors :
\[
|u_{n+1}-\ell| = |f(u_n)-f(\ell)| \le k\,|u_n-\ell|
\]
Par récurrence immédiate : $\forall n\in\mathbb{N},\ |u_n-\ell|\le k^n\,|u_0-\ell|$.
\end{propriete}

\begin{methode}[Déterminer un rang $n$ pour une précision $10^{-p}$]
\begin{enumerate}
  \item Trouver un intervalle $I$ stable contenant $\ell$ et tous les $u_n$.
  \item Majorer $|f'(x)|$ sur $I$ par une constante $k<1$ (souvent $k=\max_{x\in I}|f'(x)|$).
  \item Écrire l'encadrement : $|u_n-\ell|\le k^n\,|u_0-\ell|$.
  \item Résoudre $k^n\,|u_0-\ell| \le 10^{-p}$ : $n \ge \dfrac{\ln(10^{-p}) - \ln|u_0-\ell|}{\ln k} = \dfrac{-p\ln 10 - \ln|u_0-\ell|}{\ln k}$.
  \item Prendre le plus petit entier $n$ satisfaisant cette inégalité. Calculer $u_n$ (à la calculatrice) qui donne la valeur approchée.
\end{enumerate}
\end{methode}

\begin{exemplebox}[Application IAF sur l'exemple précédent]
$f(x)=\frac{1}{2}(x+\frac{5}{x})$, $f'(x)=\frac{1}{2}(1-\frac{5}{x^2})$. Sur $I=[\sqrt{5},3]$, $f'$ est croissante (car $f''>0$).
$|f'(x)|\le \max(|f'(\sqrt{5})|, |f'(3)|) = \max(0, \frac{1}{2}(1-\frac{5}{9})) = \frac{2}{9} \approx 0,222$.
On peut prendre $k=\frac{2}{9}$. $|u_0-\ell|=|3-\sqrt{5}|\approx 0,764$.
Erreur $\le (\frac{2}{9})^n \times 0,764$.
Pour $10^{-4}$ : $(\frac{2}{9})^n \le \frac{10^{-4}}{0,764} \approx 1,31\times 10^{-4}$.
$n\ln(2/9) \le \ln(1,31\times 10^{-4}) \implies n \ge \frac{-8,94}{-1,504} \approx 5,94$.
Donc $n=6$ suffit. $u_6 \approx 2,23607$ (valeur exacte $\sqrt{5}\approx 2,2360679$).
\end{exemplebox}

\coursec{Synthèse : Méthode complète d'étude d'une suite au Bac}

\begin{aretenir}[Fiche de révision : Les 5 étapes types]
\begin{center}
\begin{tabularx}{\linewidth}{|c|X|}
\hline
\rowcolor{popblueL} \textbf{Étape} & \textbf{Action et rédaction attendue} \\
\hline
1. \textbf{Domaine / $f$} & Définir $f$, ensemble de définition, continuité, dérivée si besoin IAF. \\
\hline
2. \textbf{Intervalle stable $I$} & Trouver $I=[m,M]$ tel que $u_0\in I$ et $f(I)\subset I$. \textbf{Prouver par récurrence} $\forall n,\ u_n\in I$. \\
\hline
3. \textbf{Monotonie} & Étudier signe de $u_{n+1}-u_n = f(u_n)-u_n$ sur $I$ (tableau de signes). Conclure : croissante/décroissante. \\
\hline
4. \textbf{Convergence} & Invoquer \textbf{Théorème de convergence monotone} (monotonie + bornitude par $I$). La suite converge vers $\ell\in I$. \\
\hline
5. \textbf{Limite $\ell$} & $\ell=f(\ell)$ (justifier passage à la limite par continuité de $f$). Résoudre. Garder racine dans $I$ compatible avec monotonie. \\
\hline
6. \textbf{Approximation (si demandé)} & Majorer $|f'|$ par $k<1$ sur $I$. Écrire $|u_n-\ell|\le k^n|u_0-\ell|$. Trouver $n$ pour précision $10^{-p}$. Calculer $u_n$. \\
\hline
\end{tabularx}
\end{center}
\end{aretenir}

\begin{attention}[Checklist finale avant de rendre la copie -- Points de vigilance MINESEC]
\begin{itemize}
  \item[$\square$] \textbf{Récurrence} : Initialisation explicite, Hérédité « Soit $n\ge n_0$, supposons $P(n)$... montrons $P(n+1)$ », Conclusion formelle.
  \item[$\square$] \textbf{Bornitude} : Majorant/Minorant \textbf{explicites} (valeurs numériques ou expressions). Pas de « la suite est bornée car elle converge » (raisonnement circulaire !).
  \item[$\square$] \textbf{Monotonie} : Signe de $u_{n+1}-u_n$ factorisé correctement. Si rapport, vérifier $u_n>0$ strictement.
  \item[$\square$] \textbf{Stabilité de $I$} : Vérifier $f(m)\ge m$ et $f(M)\le M$ (si $f$ croissante) ou utiliser les extrema de $f$ sur $I$.
  \item[$\square$] \textbf{Continuité de $f$} : Mentionner explicitement « $f$ est continue sur $I$ donc on peut passer à la limite ».
  \item[$\square$] \textbf{Équation $\ell=f(\ell)$} : Résolution algébrique complète. Écartement des racines parasites justifié par l'encadrement $m\le \ell\le M$.
  \item[$\square$] \textbf{IAF} : $k$ doit être \textbf{strictement inférieur à 1}. Vérifier que le max de $|f'|$ sur $I$ est bien $<1$. Utiliser $\ln$ pour trouver $n$ (ou essais successifs si calculatrice interdite).
  \item[$\square$] \textbf{Conclusion finale} : Phrase complète : « La suite $(u_n)$ converge vers $\ell = \dots$ » ou « La suite diverge vers $+\infty$ ».
\end{itemize}
\end{attention}

\begin{savaistu}[Le saviez-vous ? -- Contexte camerounais et histoire]
\begin{itemize}
  \item \textbf{Algorithme d'Héron (Babylone, ~1700 av. J.-C.) :} La suite $u_{n+1}=\frac{1}{2}(u_n+\frac{a}{u_n})$ calcule $\sqrt{a}$. C'est la méthode utilisée dans les processeurs des téléphones (MTN/Orange) et ordinateurs pour les racines carrées !
  \item \textbf{Intérêt composé :} Un placement de $C_0$ FCFA à $t\%$ par an : $C_{n+1}=C_n(1+\frac{t}{100})$. Suite géométrique, modélisation de l'épargne ou des crédits (CFC, banques).
  \item \textbf{Modèle logistique (Population) :} $u_{n+1}=r u_n(1-u_n)$ modélise la croissance d'une population (ex: poissons dans un barrage, cultures) avec ressources limitées. Chaos pour $r>3,57$ !
  \item \textbf{Mathématiques au Cameroun :} Le théorème de convergence monotone repose sur l'\textbf{axiome de la borne supérieure} (propriété d'archimède + complétude de $\mathbb{R}$), distinction fondamentale avec $\mathbb{Q}$ (une suite croissante majorée de rationnels peut ne pas converger dans $\mathbb{Q}$, ex: approximations de $\sqrt{2}$).
\end{itemize}
\end{savaistu}

\begin{exercice}[Entraînement Bac -- Suite récurrente classique]
Soit la suite $(u_n)$ définie par $u_0=1$ et $u_{n+1}=\frac{2u_n+3}{u_n+2}$.
\begin{enumerate}
  \item Montrer que $\forall n\in\mathbb{N},\ 1\le u_n\le \sqrt{3}$.
  \item Étudier la monotonie de $(u_n)$.
  \item En déduire la convergence et calculer la limite $\ell$.
  \item À l'aide de l'IAF, déterminer un entier $n$ tel que $|u_n-\ell|<10^{-3}$.
\end{enumerate}
\end{exercice}

\begin{corrige}[Exercice 1]
\begin{enumerate}
  \item $f(x)=\frac{2x+3}{x+2}$ sur $[1,+\infty[$. $f'(x)=\frac{1}{(x+2)^2}>0$ (croissante). $f(1)=\frac{5}{3}\approx 1,66$, $f(\sqrt{3})=\frac{2\sqrt{3}+3}{\sqrt{3}+2}=\sqrt{3}$ (vérification : $(2\sqrt{3}+3)= \sqrt{3}(\sqrt{3}+2)=3+2\sqrt{3}$). Donc $f([1,\sqrt{3}])=[\frac{5}{3},\sqrt{3}]\subset [1,\sqrt{3}]$. Récurrence : $u_0=1\in I$. Si $u_n\in I$, $u_{n+1}=f(u_n)\in f(I)\subset I$.
  \item $u_{n+1}-u_n = \frac{2u_n+3}{u_n+2}-u_n = \frac{2u_n+3-u_n^2-2u_n}{u_n+2} = \frac{3-u_n^2}{u_n+2}$. Comme $u_n\le \sqrt{3}$, $3-u_n^2\ge 0$. Dénominateur $>0$. Donc $u_{n+1}-u_n\ge 0$. Suite croissante.
  \item Croissante et majorée par $\sqrt{3}$ $\implies$ converge vers $\ell\in[1,\sqrt{3}]$. $\ell=f(\ell)=\frac{2\ell+3}{\ell+2} \iff \ell^2+2\ell=2\ell+3 \iff \ell^2=3 \iff \ell=\sqrt{3}$ (car $\ell>0$).
  \item $f'(x)=\frac{1}{(x+2)^2}$. Sur $I=[1,\sqrt{3}]$, $f'$ décroissante, max en $1$ : $k=f'(1)=\frac{1}{9}$. $|u_0-\ell|=|\sqrt{3}-1|\approx 0,732$. Erreur $\le (\frac{1}{9})^n \times 0,732 < 10^{-3} \iff (\frac{1}{9})^n < 1,366\times 10^{-3}$. $n=2 : 1/81\approx 0,012$. $n=3 : 1/729\approx 0,00137$. $n=4 : 1/6561\approx 1,5\times 10^{-4}$. Donc $n=4$ suffit.
\end{enumerate}
\end{corrige}

\findecours

\end{document}
